% concatenated from pebble-motion-problem-on-trees-20260401.tex
\documentclass[letter,11pt]{article}
%\documentclass[a4paper,11pt]{article}
% \documentclass[a4paper]{article}
%\usepackage[latin1]{inputenc}
\usepackage{graphicx}
\usepackage{xcolor}
\usepackage{amsmath}
\usepackage{amsfonts}
\usepackage{cite}
\usepackage{ascmac}
\usepackage{comment}
\usepackage{amsthm}
% just comment this out if you don't have natbib
%\usepackage[round]{natbib}
%\author{Tomoki Nakamigawa\addressmark{1}}
%\title
%{}
%\address{\addressmark{1}

\usepackage{hyperref}
\usepackage{appendix}

%\usepackage{hyperref}

\usepackage{ulem}

\oddsidemargin=0mm
\evensidemargin=0mm
%%%\oddsidemargin=9.6mm
%%%\evensidemargin=9.6mm
%%%\def\sloppy{\tolerance=9999 \hfuzz=.5pt \vfuzz=.5pt}
%%%\renewcommand{\baselinestretch}{1.2}
\textwidth=164mm
\textheight=221mm
\topmargin=-15mm
%\marginparsep=18pt

%%\headheight=0truemm 
%%\headsep=0truemm
%%\doublespacing





\newtheorem{thm}{Theorem}
\newtheorem*{thminformal}{Theorem [informal]}
%\renewcommand{\thethm}{\arabic{thm}.}
%\renewcommand{\thethm}{\arabic{section}.\arabic{thm}.}
\newtheorem{ex}{Example}
\newtheorem{la}[thm]{Lemma}
\newtheorem{subla}[thm]{Sublemma}
\newtheorem{cor}[thm]{Corollary}
\newtheorem{prop}[thm]{Proposition}
\newtheorem{obs}[thm]{Observation}
\newtheorem{conj}{Conjecture}
\renewcommand{\theconj}{\arabic{conj}}
%\renewcommand{\theconj}{}
%\newcommand{\owari}{\hfill\hbox{\rule{6pt}{10pt}}}
\newcommand{\white}{\qquad{\framebox{\rule{0pt}{4pt}}}}
\newcommand{\owari}{\hfill\white}
%\newcommand{\qed}{\hfill$\BBox$}
\def\GG{\tilde{G}}
\def\cs{\choose}
\def\ov{\overline}





\theoremstyle{definition}
\newtheorem{rem}{Remark}







%%%%%%%%%%%%%%%%%%%%%%%%%%%%%%%%%%%%%%%%%%%%%%%%%%%%%%%%%%%%%%%
%%%Torikeshi-sen
%%%%%%%%%%%%%%%%%%%%%%%%%%%%%%%%%%%%%%%%%%%%%%%%%%%%%%%%%%%%%%%
%\usepackage{ulem}
%%\renewcommand{\reviewstrike}[1]{\sout{#1}}
%%%%%%%%%%%%%%%%%%%%%%%%%%%%%%%%%%%%%%%%%%%%%%%%%%%%%%%%%%%%%%%
%%%%%%%%%%%%%%%%%%%%%%%%%%%%%%%%%%%%%%%%%%%%%%%%%%%%%%%%%%%%%%%
%%%%%%%%%%%%%%%%%%%%%%%%%%%%%%%%%%%%%%%%%%%%%%%%%%%%%%%%%%%%%%%


%\usepackage{tikz}
\def\P{{$\mathcal{P}$}}
\def\NP{{$\mathcal{NP}$}}
\def\diam{\mathrm{diam}}
\def\dist{\mathrm{dist}}
\def\blue#1{\textcolor{blue}{#1}}
\def\red#1{\textcolor{red}{#1}}
\def\green#1{\textcolor{green}{#1}}
\def\purple#1{\textcolor{purple}{#1}}

%	\theoremstyle{definition}
%	\newtheorem{problem}[theorem]{Problem}
%	\theoremstyle{claimstyle}
%\newcounter{claim}
%\newtheorem{subclm}[claim]{Subclaim}


%%%%%%%%%%%%%%%%%%%%%%%%%%%%%%%%%%%%%%%%%%%%%%%%%%%%%%%%%%%%%%%%%%%%%%%%%%%%%%%%%%%%
%%%%%%%%%%%%%%%%%%%%%%%%%%%%%%%%%%%%%%%%%%%%%%%%%%%%%%%%%%%%%%%%%%%%%%%%%%%%%%%%%%%%
%%%%%
%%%%% Tadashi SAKUMA's macro START
%%%%%
%%%%%%%%%%%%%%%%%%%%%%%%%%%%%%%%%%%%%%%%%%%%%%%%%%%%%%%%%%%%%%%%%%%%%%%%%%%%%%%%%%%%
%%%%%%%%%%%%%%%%%%%%%%%%%%%%%%%%%%%%%%%%%%%%%%%%%%%%%%%%%%%%%%%%%%%%%%%%%%%%%%%%%%%%

%\newtheorem{lem}[thm]{Lemma}
%\newtheorem{prob}{Problem}
\theoremstyle{definition}
\newtheorem{alg}{Algorithm}
\newtheorem{prob}{Problem}
  \renewcommand{\theprob}{}


\DeclareMathOperator{\MaximumDegree}{\mathit{\Delta}}
\DeclareMathOperator{\function}{\varphi}
\DeclareMathOperator{\flip}{flip}
\DeclareMathOperator{\diff}{Diff}
\DeclareMathOperator{\Nb}{\it{N}}
%\DeclareMathOperator{\Nb}{Neighbor}
\DeclareMathOperator{\modu}{modulo}
\DeclareMathOperator{\Od}{O}
\DeclareMathOperator{\od}{o}
\DeclareMathOperator{\Eod}{\Theta}
\DeclareMathOperator{\IA}{IA}
\DeclareMathOperator{\BT}{BT}
\DeclareMathOperator{\Peb}{Peb}
\DeclareMathOperator{\col}{col}
\DeclareMathOperator{\IC}{IC}
\DeclareMathOperator{\Irr}{Irr}
\DeclareMathOperator{\id}{id}
\DeclareMathOperator{\Alg}{Alg}
\DeclareMathOperator{\Cen}{Centroid}
\DeclareMathOperator{\CenE}{CentroidEdge}
\DeclareMathOperator{\CenV}{CentroidVertex}
\DeclareMathOperator{\Joi}{Joint}
\DeclareMathOperator{\JoiV}{JointVertex}
\DeclareMathOperator{\JoiE}{JointEdge}
\DeclareMathOperator{\NS}{NodeSet}
\DeclareMathOperator{\PacCon}{PackedConfig}
\DeclareMathOperator{\Pac}{Packing}
\DeclareMathOperator{\CutOff}{CutOff}
\DeclareMathOperator{\Extraction}{Extraction}
\DeclareMathOperator{\Main}{Main}
\DeclareMathOperator{\PendantCycle}{PendantCycle}
\DeclareMathOperator{\SP}{SP}
\DeclareMathOperator{\D}{D}
\DeclareMathOperator{\CL}{CL}
%\DeclareMathOperator{\OPT}{OPT}


\DeclareMathOperator{\aapeq}{\underset{\textrm{\tiny{AAP}}}{\cong}}
\DeclareMathOperator{\sgaeq}{\underset{\textrm{\tiny{SGA}}}{\cong}}

\DeclareMathOperator{\C}{C}
%\DeclareMathOperator{\s}{s}
%\DeclareMathOperator{\cs}{cs}

\newcommand{\1}{\boldmath$~1~$\unboldmath}
\newcommand{\A}{\boldmath$A$\unboldmath}
\newcommand{\Poly}{\boldmath$~P~$\unboldmath}
\newcommand{\x}{\boldmath$~x~$\unboldmath}
\newcommand{\p}{\boldmath$~p~$\unboldmath}
\newcommand{\pp}{\boldmath$p$\unboldmath}
\newcommand{\vS}{\boldmath$S$\unboldmath}
\newcommand{\Rea}{\mathbb{R}}
\newcommand{\Z}{\mathbb{Z}}
\newcommand{\N}{\mathbb{N}}
\newcommand{\Q}{\mathbb{Q}}
\newcommand{\norm}[1]{\parallel\hspace*{-0.3em}{#1}\hspace*{-0.3em}\parallel}

\newcommand{\G}{\mathcal{G}^{\mathrm{cs}}}
\newcommand{\Ord}{\mathrm{O}}
\newcommand{\ord}{\mathrm{o}}

\newcommand{\OPT}{\mathrm{OPT}}

%\newcommand{\owari}{\hfill\hbox{\rule{6pt}{10pt}}}

\newcommand{\qedgo}{\vspace{0.35cm}}
\newcommand{\kuttuke}{\hspace{-0.1em}}
\newcommand{\tume}{\hspace*{0.2em}}
\newcommand{\piriod}{\hspace{-1.8mm}{\bf .} \hspace{5pt}}

\def\setdel{\!\setminus\!\!} \def\mbzero {\mbox{{\rm {\bf 0}}}}
\def\mbone {\mbox{{\rm {\bf 1}}}}
\def\P{{$\mathcal{P}$}}
\def\NP{{$\mathcal{NP}$}}
\def\APX{{$\mathcal{APX}$}}
\def\p{\mathcal{P}}
\def\np{\mathcal{NP}}
\def\EXT{\rm{EXPTIME}}
\def\calA{\mathcal{A}}
\def\calB{\mathcal{B}}
\def\calC{\mathcal{C}}
\def\co{\textrm{co-}}
\def\YES{\textrm{YES}}
\def\NO{\textrm{NO}}
\def\deg{\textrm{deg}}
%\def\BT{\textrm{BT}}


%\def\Cen{\textrm{Centroid}}
%\def\CenV{\textrm{CentroidVertex}}
%\def\CenE{\textrm{CentroidEdge}}
%\def\Join{\textrm{Joint}}
%\def\JoinV{\textrm{JointVertex}}
%\def\JoinE{\textrm{JointEdge}}
%\def\Alg{\textrm{Alg}}



\newcommand{\ol}[1]{\overline{#1}}
%%\newcommand{\ul}[1]{\underline{#1}}
\newcommand{\ff}[1]{\function({#1})}
\newcommand{\Deg}[1]{\MaximumDegree({#1})}
\newcommand{\olra}[1]{\overleftrightarrow{#1}}
\newcommand{\ora}[1]{\overrightarrow{#1}}
\newcommand{\ola}[1]{\overleftarrow{#1}}
\newcommand{\rev}[1]{\overleftarrow{#1}}
\newcommand{\either}[2]{$~{\rm either}~$ {{#1}} $~{\rm or}~$ {#2}}

\newcommand{\NOT}[1]{~\ooalign{\hfil/\hfil\crcr{#1}}~}

\def\BBox{\rule{2mm}{3mm}} \def\QED{\hfill$\BBox$}

%\newcommand{\maru}[1]{{\ooalign{\hfill$\scriptsize#1$\hfill\crcr\scalebox{1.3}{$\bigcirc$}}}}
\newcommand{\maru}[1]{{\ooalign{\hfill$\small#1$\hfill\crcr\scalebox{1.3}{$\bigcirc$}}}}

%%%%%%%%%%%%%%%%%%%%%%%%%%%%%%%%%%%%%%%%%%%%%%%%
%%%%%%%%%%%%%%%%%%%%%%%%%%%%%%%%%%%%%%%%%%%%%%%%
%%%
%%% For algorithm enumeration environment
%%%
%%%%%%%%%%%%%%%%%%%%%%%%%%%%%%%%%%%%%%%%%%%%%%%%
%%%%%%%%%%%%%%%%%%%%%%%%%%%%%%%%%%%%%%%%%%%%%%%%

\newenvironment{alg-enumerate}{%
  \begin{enumerate}%
   \renewcommand{\theenumi}{\arabic{enumi}}%
   \renewcommand{\labelenumi}{\rm{{\maru{\theenumi}}\tume\/}}%
   \renewcommand{\labelenumii}{\rm{$(${\maru{\theenumi}}-\theenumii$)$\tume\/}}%
 }{\end{enumerate}%
}
%%%%%%%%%%%%%%%%%%%%%%%%%%%%%%%%%%%%%%%%%%%%%%%%
%%%%%%%%%%%%%%%%%%%%%%%%%%%%%%%%%%%%%%%%%%%%%%%%
%%%
%%% For algorithm enumeration environment
%%%
%%%%%%%%%%%%%%%%%%%%%%%%%%%%%%%%%%%%%%%%%%%%%%%%
%%%%%%%%%%%%%%%%%%%%%%%%%%%%%%%%%%%%%%%%%%%%%%%%



\newenvironment{(R)-enumerate}{%
  \begin{enumerate}%
   \renewcommand{\theenumi}{\Roman{enumi}}%
   \renewcommand{\theenumii}{--\roman{enumii}}%
   \renewcommand{\labelenumi}{$($\theenumi$)$}%
   \renewcommand{\labelenumii}{$($\theenumi\theenumii$)$}%
 }{\end{enumerate}%
}

\newenvironment{(A)-enumerate}{%
  \begin{enumerate}%
   \renewcommand{\theenumi}{\arabic{enumi}}%
   \renewcommand{\theenumii}{--\roman{enumii}}%
   \renewcommand{\labelenumi}{$($\theenumi$)$}%
   \renewcommand{\labelenumii}{$($\theenumi\theenumii$)$}%
 }{\end{enumerate}%
}

\newenvironment{claimenumerate}{%
 \begin{quote}
  \begin{enumerate}%
   \renewcommand{\theenumi}{\arabic{enumi}}%
   \renewcommand{\theenumii}{\roman{enumii}}%
   \renewcommand{\labelenumi}{{\bf Claim\tume\theenumi\tume \/}}%
   \renewcommand{\labelenumii}{{\bf Claim\tume\theenumii\tume \/}}%
   }{\end{enumerate}\end{quote}}

\newenvironment{Proof}{\expandafter
\paragraph{{\bf Proof\/.}}}{\hspace{5mm}\owari\qedgo}

\newenvironment{Proofof}[1]{\expandafter
\paragraph{{\it Proof of #1.}}}{\hspace{5mm}\owari\qedgo}


%namelist environment (Nelson H.~F.~Beebe~and~Tadashi~Sakuma)
%form: \begin{namelist}{width}
  \newcommand{\namelistlabel}[1]{\hspace*{10pt}\mbox{#1}:\hspace{\fill}}
  \newenvironment{namelist}[1]{%
    \begin{list}{}
      {\let\makelabel\namelistlabel
      \settowidth{\labelwidth}{#1}
      \setlength{\leftmargin}{1.1\labelwidth}
      \setlength{\labelsep}{5pt}
      \setlength{\parsep}{0pt}}
  }{%
  \end{list}}

\def\NZQ{\mathbb}               % the font for N,Z,Q,R,C
\def\NN{{\NZQ N}} \def\QQ{{\NZQ Q}} \def\ZZ{{\NZQ Z}} \def\RR{{\NZQ R}} 
\def\CC{{\NZQ C}} \def\AA{{\NZQ A}} \def\BB{{\NZQ B}} \def\PP{{\NZQ P}} 
\def\FF{{\NZQ F}} \def\GG{{\NZQ G}} \def\HH{{\NZQ H}} 
\def\CAL{\mathcal}
\def\cA{{\CAL A}} \def\cB{{\CAL B}} \def\cC{{\CAL C}} \def\cD{{\CAL D}} 
\def\cE{{\CAL E}} \def\cF{{\CAL F}} \def\cG{{\CAL G}} \def\cH{{\CAL H}} 
\def\cI{{\CAL I}} \def\cJ{{\CAL J}} \def\cK{{\CAL K}} \def\cL{{\CAL L}} 
\def\cM{{\CAL M}} \def\cN{{\CAL N}} \def\cO{{\CAL O}} \def\cP{{\CAL P}} 
\def\cQ{{\CAL Q}} \def\cR{{\CAL R}} \def\cS{{\CAL S}} \def\cT{{\CAL T}} 
\def\cU {{\CAL U}} \def\cV{{\CAL V}} \def\cW{{\CAL W}} \def\cX{{\CAL X}} 
\def\cY{{\CAL Y}} \def\cZ{{\CAL Z}} 
\def\FRAC{\mathfrak}
\def\fA{{\FRAC A}} \def\fB{{\FRAC B}} \def\fC{{\FRAC C}} \def\fD{{\FRAC D}} 
\def\fE{{\FRAC E}} \def\fF{{\FRAC F}} \def\fG{{\FRAC G}} \def\fH{{\FRAC H}} 
\def\fI{{\FRAC I}} \def\fJ{{\FRAC J}} \def\fK{{\FRAC K}} \def\fL{{\FRAC L}} 
\def\fM{{\FRAC M}} \def\fN{{\FRAC N}} \def\fO{{\FRAC O}} \def\fP{{\FRAC P}} 
\def\fQ{{\FRAC Q}} \def\fR{{\FRAC R}} \def\fS{{\FRAC S}} \def\fT{{\FRAC T}} 
\def\fU{{\FRAC U}} \def\fV{{\FRAC V}} \def\fW{{\FRAC W}} \def\fX{{\FRAC X}} 
\def\fY{{\FRAC Y}} \def\fZ{{\FRAC Z}} 
\def\SCR{\mathscr}
\def\sA{{\SCR A}} \def\sB{{\SCR B}} \def\sC{{\SCR C}} \def\sD{{\SCR D}} 
\def\sE{{\SCR E}} \def\sF{{\SCR F}} \def\sG{{\SCR G}} \def\sH{{\SCR H}}
\def\sI{{\SCR I}} \def\sJ{{\SCR J}} \def\sK{{\SCR K}} \def\sL{{\SCR L}}
\def\sM{{\SCR M}} \def\sN{{\SCR N}} \def\sO{{\SCR O}} \def\sP{{\SCR P}}
\def\sQ{{\SCR Q}} \def\sR{{\SCR R}} \def\sS{{\SCR S}} \def\sT{{\SCR T}}
\def\sU{{\SCR U}} \def\sV{{\SCR V}} \def\sW{{\SCR W}} \def\sX{{\SCR X}}
\def\sY{{\SCR Y}} \def\sZ{{\SCR Z}}
\def\red{\textcolor{red}}
\def\purple{\textcolor{purple}}
\def\green{\textcolor{green}}


%%%%%%%%%%%%%%%%%%%%%%%%%%%%%%%%%%%%%%%%%%%%%%%%%%%%%%%%%%%%%%%%%%%%%%%%%%%%%%%%%%%%
%%%%%%%%%%%%%%%%%%%%%%%%%%%%%%%%%%%%%%%%%%%%%%%%%%%%%%%%%%%%%%%%%%%%%%%%%%%%%%%%%%%%
%%%%%
%%%%% Tadashi SAKUMA's macro END
%%%%%
%%%%%%%%%%%%%%%%%%%%%%%%%%%%%%%%%%%%%%%%%%%%%%%%%%%%%%%%%%%%%%%%%%%%%%%%%%%%%%%%%%%%
%%%%%%%%%%%%%%%%%%%%%%%%%%%%%%%%%%%%%%%%%%%%%%%%%%%%%%%%%%%%%%%%%%%%%%%%%%%%%%%%%%%%













\def\blue#1{\textcolor{blue}{#1}}
\def\red#1{\textcolor{red}{#1}}
\def\green#1{\textcolor{green}{#1}}
\def\purple#1{\textcolor{purple}{#1}}

\newcounter{aaa}
%\renewcommand{\thefootnote}{\fnsymbol{footnote}}


%\parskip=10pt

\begin{document}

\thispagestyle{empty}

\begin{center}
%{\textcolor{blue}{\Large The Structure of Cubic Complexity in Pebble Motion}}
 {\Large Structural Origins of Cubic Complexity in Pebble Motion}
% {\Large Characterizing Cubic Complexity in Pebble Motion}
%    {\Large On the order of the shortest solution sequences for the pebble motion problems}
\end{center}

%%%%%%%%%%%%%%%%%%%%%%%%%%%%%%%%%%%%%%%%%%%%%%%%%%%%%%%%%%%%%%%%%%%%%%%%%%%%%%%%%
%%%%%%%%%%%%%%%%%%%%%%%%%%%%%%%%%%%%%%%%%%%%%%%%%%%%%%%%%%%%%%%%%%%%%%%%%%%%%%%%%
%%%%%%%%%%%%%%%%%%%%%%%%%%%%%%%%%%%%%%%%%%%%%%%%%%%%%%%%%%%%%%%%%%%%%%%%%%%%%%%%%
%%%%%%%%%%%%%%%%%%%%%%%%%%%%%%%%%%%%%%%%%%%%%%%%%%%%%%%%%%%%%%%%%%%%%%%%%%%%%%%%%
%%%%%%%%%%%%%%%%%%%%%%%%%%%%%%%%%%%%%%%%%%%%%%%%%%%%%%%%%%%%%%%%%%%%%%%%%%%%%%%%%
%%%%%%%%%%%%%%%%%%%%%%%%%%%%%%%%%%%%%%%%%%%%%%%%%%%%%%%%%%%%%%%%%%%%%%%%%%%%%%%%%
%%%%%%%%%%%%%%%%%%%%%%%%%%%%%%%%%%%%%%%%%%%%%%%%%%%%%%%%%%%%%%%%%%%%%%%%%%%%%%%%%
%%%%%%%%%%%%%%%%%%%%%%%%%%%%%%%%%%%%%%%%%%%%%%%%%%%%%%%%%%%%%%%%%%%%%%%%%%%%%%%%%
%%%%%%%%%%%%%%%%%%%%%%%%%%%%%%%%%%%%%%%%%%%%%%%%%%%%%%%%%%%%%%%%%%%%%%%%%%%%%%%%%
%%%%%%%%%%%%%%%%%%%%%%%%%%%%%%%%%%%%%%%%%%%%%%%%%%%%%%%%%%%%%%%%%%%%%%%%%%%%%%%%%
%\maketitle 
\begin{center}
{\large Tomoki Nakamigawa
\footnote{This work was supported by JSPS KAKENHI Grant Number JP16K05260.\\
Email: {\texttt{nakamigwt@gmail.com}}}
%Email: {\texttt{nakami@info.shonan-it.ac.jp}}}
}
\end{center}
%\maketitle 
\begin{center}
{
Department of Information Science \\%
Shonan Institute of Technology \\%
1-1-25 Tsujido-Nishikaigan, Fujisawa 251-8511, Japan
}
\end{center}
\begin{center}
{\large Tadashi Sakuma
\footnote{This work was supported by JSPS KAKENHI Grant Number JP16K05260, JP26400185, JP18K03388.\\
Email: {\texttt{sakuma@sci.kj.yamagata-u.ac.jp}}}
}
\end{center}
\begin{center}
{
Faculty of Science \\%
Yamagata University \\%
1-4-12 Kojirakawa, Yamagata 990-8560, Japan
}
\end{center}
%%%%%%%%%%%%%%%%%%%%%%%%%%%%%%%%%%%%%%%%%%%%%%%%%%%%%%%%%%%%%%%%%%%%%%%%%%%%%%%%%
%%%%%%%%%%%%%%%%%%%%%%%%%%%%%%%%%%%%%%%%%%%%%%%%%%%%%%%%%%%%%%%%%%%%%%%%%%%%%%%%%
%%%%%%%%%%%%%%%%%%%%%%%%%%%%%%%%%%%%%%%%%%%%%%%%%%%%%%%%%%%%%%%%%%%%%%%%%%%%%%%%%
%%%%%%%%%%%%%%%%%%%%%%%%%%%%%%%%%%%%%%%%%%%%%%%%%%%%%%%%%%%%%%%%%%%%%%%%%%%%%%%%%
%%%%%%%%%%%%%%%%%%%%%%%%%%%%%%%%%%%%%%%%%%%%%%%%%%%%%%%%%%%%%%%%%%%%%%%%%%%%%%%%%
%%%%%%%%%%%%%%%%%%%%%%%%%%%%%%%%%%%%%%%%%%%%%%%%%%%%%%%%%%%%%%%%%%%%%%%%%%%%%%%%%
%%%%%%%%%%%%%%%%%%%%%%%%%%%%%%%%%%%%%%%%%%%%%%%%%%%%%%%%%%%%%%%%%%%%%%%%%%%%%%%%%
%%%%%%%%%%%%%%%%%%%%%%%%%%%%%%%%%%%%%%%%%%%%%%%%%%%%%%%%%%%%%%%%%%%%%%%%%%%%%%%%%
%%%%%%%%%%%%%%%%%%%%%%%%%%%%%%%%%%%%%%%%%%%%%%%%%%%%%%%%%%%%%%%%%%%%%%%%%%%%%%%%%
%%%%%%%%%%%%%%%%%%%%%%%%%%%%%%%%%%%%%%%%%%%%%%%%%%%%%%%%%%%%%%%%%%%%%%%%%%%%%%%%%


\vspace{10pt}



\begin{abstract}
The pebble motion problem (PMP) asks whether one configuration of labeled pebbles on a graph 
can be transformed into another by moving pebbles to adjacent unoccupied vertices. 
It is a fundamental model of graph reconfiguration and is closely related to multi-agent path finding (MAPF).

A central open problem since Kornhauser, Miller, and Spirakis (FOCS 1984) is to understand the origin 
of the classical $\Theta(N^3)$ worst-case behavior. 
While it is known that every feasible instance on an $N$-vertex graph admits a solution sequence of 
length $\Ord(N^3)$, it has remained unclear which instances actually require cubic complexity.

In this paper, we resolve the long-standing complexity of the pebble motion problem on trees. 
We show that every feasible instance on an $N$-vertex tree admits a solution sequence of length 
$\Ord(N^2 \log N)$, computable by an output-sensitive algorithm. 
Since a lower bound of $\Omega(N^2)$ is known, this establishes that the $\Theta(N^3)$ phenomenon 
does not occur on trees and nearly closes the gap $\Omega(N^2)\le \OPT(N)\le \Ord(N^3)$ up to 
a logarithmic factor.

Building on this result, we extend our approach to general graphs by applying the tree algorithm 
to breadth-first spanning trees. 
This yields an efficient framework that produces $o(N^3)$-length solution sequences for a broad class 
of instances, including the classical square-grid example, where we recover the $\Ord(N^{3/2})$ bound 
observed by Kornhauser, Miller, and Spirakis.

Finally, by analyzing the behavior of this algorithm, we obtain strong structural restrictions governing 
when $\Theta(N^3)$ complexity can arise.
We show that such behavior is possible only under highly constrained conditions, specifically when 
$\Theta(N)$ degree-two vertices lie on cycles of length $\Theta(N)$, with each cycle being 
the shortest containing the corresponding vertex.
\end{abstract}
Keywords: Pebble motion problem, Multi-agent path finding, Token swapping, Reconfiguration, $15$-puzzle, Trees










%\begin{abstract}
%The pebble motion problem (PMP) involves transforming a configuration of 
%labeled pebbles on a graph by moving them to adjacent unoccupied vertices. 
%It serves as a foundational model for reconfiguration processes and is closely 
%related to multi-agent path finding (MAPF).
%
%Let $G$ be a graph with $N$ vertices and $n$ pebbles placed on distinct vertices. 
%Kornhauser, Miller, and Spirakis (FOCS 1984) showed that for general graphs, every 
%feasible instance admits a solution sequence of length $\Ord(N^3)$ and presented 
%instances requiring $\Omega(N^3)$ moves. 
%However, they noted that the number of instances requiring $\Theta(N^3)$ moves 
%is likely small, leaving the design of an efficient algorithm as an open problem.
%
%Despite decades of study, the question of which instances require $\Theta(N^3)$ 
%moves, particularly for trees, remains unresolved.
%
%In this paper, we address these long-standing issues directly. 
%We show that every feasible instance of the pebble motion problem on 
%a tree with $N$ vertices admits a solution sequence of length 
%$\Ord(N^2 \log N)$, which can be computed using an output-sensitive 
%algorithm. 
%Since a lower bound of $\Omega(N^2)$ for trees is known, our result nearly 
%closes the long-standing gap $\Omega(N^2) \le \OPT(N) \le \Ord(N^3)$, 
%with only a logarithmic factor separating them.
%
%Next, we extend this result to general graphs by applying our method to breadth-first 
%spanning trees. 
%Our algorithm computes the $\Ord(N^{3/2})$ move bound for the square grid example 
%discussed in Kornhauser, Miller, and Spirakis (FOCS 1984), addressing one of the open 
%problems they posed, and generalizes this to a broader class of graphs.
%
%Finally, we prove a structural restriction that explains the $\Theta(N^3)$ phenomenon 
%for general graphs. 
%Our analysis reveals that such solution sequences arise only under highly constrained 
%conditions, specifically when $\Theta(N)$ degree-two vertices lie on cycles of length 
%$\Theta(N)$, with each cycle being the shortest containing the corresponding vertex.
%\end{abstract}
%Keywords: Pebble motion problem, Multi-agent path finding, Token swapping,  
%Reconfiguration, $15$-puzzle, Trees


\newpage



\section{Introduction}\label{intro}

\subsection{Definition of the Pebble Motion Problem}

Let $G$ be a finite undirected graph with no multiple edges or loops.
We denote its vertex and edge sets by $V(G)$ and $E(G)$, and write
$N := |V(G)|$ for the number of vertices.
Let $P=\{1,\ldots,n\}$ be a set of pebbles with $n \leq N-1$.
Let $\D(G)$ denote the diameter of $G$. 
For each vertex $v \in V(G)$, let $\CL(G,v)$ denote the length of a
shortest cycle containing $v$ if such a cycle exists, and set
$\CL(G,v) := 1$ otherwise. Define $\CL(G) := \max_{v \in V(G)} \CL(G,v)$.


In the context of the Pebble Motion Problem, the graph $G$ serves as the
\textit{board graph} on which the pebbles move.
When the board graph $G$ is a tree, we refer to it as the
\textit{board tree}.
A \textit{configuration} of $P$ on the board graph $G$ is a function 
$f : V(G) \to \{0,1,\ldots,n\}$ such that $|f^{-1}(i)| = 1$ for each
$1 \le i \le n$.
Thus $f^{-1}(i)$ denotes the vertex occupied by pebble $i$, while
$f^{-1}(0)$ denotes the set of unoccupied vertices.
A \textit{move} shifts a pebble from a vertex to an adjacent
unoccupied vertex.
The \textit{Pebble Motion Problem} (PMP) asks whether one configuration
can be transformed into another by a sequence of such moves.
The classical ``15-puzzle'' introduced by Loyd~\cite{loyd1959mathematical}
is a standard example in which the board graph is a $4\times4$ grid 
(cf. \cite{aa887662-c3e2-30b0-a475-bcd8c655206a, 
WOS:000083999100001}).

The pebble motion problem is a fundamental model of
reconfiguration on graphs.
It is closely related to the \textit{Multi-Agent Path Finding} (MAPF)
problem studied extensively in artificial intelligence and robotics.





\subsection{Background}

The algorithmic theory of the pebble motion problem was initiated by
Kornhauser, Miller, and Spirakis in their seminal
FOCS~1984 paper~\cite{715921}.
They showed that for general graphs, every feasible instance admits
a solution sequence of length $\Ord(N^3)$ and also exhibited families
of instances requiring $\Omega(N^3)$ moves.
Thus, the worst-case length of a shortest solution sequence is $\Theta(N^3)$.

The same paper also considered the special case in which the board
graph is a tree, referred to as the \textit{Pebble Motion Problem on Trees} (PMT).
They informally stated that the length of the shortest solution
sequence for PMT should be $\Ord(N^2)$.
However, the proof sketch relied on an assumption that does not hold
in general.
Indeed, Auletta, Monti, Parente, and Persiano~\cite{WOS:000077614700003} later remarked
that they were unable to verify this claim.
Consequently, subsequent work relied only on the $\Ord(N^3)$ bound
implied by the general algorithm of~\cite{715921}.
Even the most recent algorithm of Ardizzoni, Saccani, Consolini, Locatelli, 
and Nebel~\cite{WOS:001177663400002} still has worst-case complexity 
$\Ord(N^3)$ on trees.
As a consequence, the true complexity of the pebble motion problem on
trees has remained unresolved for more than four decades.

At the same time, the tree case plays a central role in the theory of
PMP and MAPF.
Several practical algorithms for general MAPF reduce the problem,
via appropriate techniques, to instances of the pebble motion problem
on trees and solve them on those trees.
Understanding the true complexity of the pebble motion problem on
trees is therefore a fundamental question.

The best known lower bound is $\Omega(N^2)$, obtained from instances
where $\Theta(N)$ pebbles must traverse a path of length $\Theta(N)$.
Let $\OPT(N)$ denote the maximum length of a shortest solution
sequence over all trees with $N$ vertices and all feasible instances.
Thus the worst-case complexity has remained within the gap
\[ 
\Omega(N^2) \le \OPT(N) \le \Ord(N^3).
\]
It is important to note that edges in the graph serve as essential tools 
for enabling the exchange of pebbles with empty vertices. 
As the number of edges decreases, the difficulty of the problem increases 
due to the reduced opportunities for such exchanges. 
Trees provide only limited opportunities for such exchanges, 
which makes the pebble motion problem inherently more difficult
on trees than on other graph structures.
This difficulty manifests itself in extreme cases in the form of {\it infeasibility}. 
Such situations arise when the number of empty vertices is insufficient 
relative to the length of a corridor (i.e., a maximal path consisting of 
consecutive degree-two vertices), so that the arrangement of pebbles in 
a tree makes it impossible to reach the target configuration. 
This phenomenon does not arise in $2$-connected graphs other than cycles,  
provided that the number of empty vertices is at least two.

Beyond the existence of worst-case instances requiring
$\Theta(N^3)$ moves on general graphs, no sharper structural understanding 
of the complexity of the pebble motion problem had been obtained.
Although the algorithms for $2$-connected graphs and for trees
are structurally quite different, reflecting the fundamental
differences between these graph classes,
no asymptotic separation between the two settings had been established. 

Already in the paper~\cite{715921},
it was suggested in the concluding remarks that
``... it seems that only a small fraction of the graph puzzles actually
require $\Ord(N^3)$ moves.''
In light of this observation, and given the growing importance of
multi-agent systems, including MAPF and other distributed particle systems,
as well as the role of the pebble motion problem as a
fundamental underlying model for such systems,
this state of affairs has been unsatisfactory.

The present paper is motivated by the goal of identifying the structural
origin of cubic behavior in general graphs, while resolving the
long-standing complexity gap for trees.


\subsection{Our Contribution}


In this paper we provide a structural explanation for 
the classical cubic phenomenon in pebble motion.
In particular, we show that $\Theta(N^3)$-length solution sequences 
cannot arise on trees, and hence must arise from additional
structures present in general graphs.

As a consequence of this insight, we obtain the following result.
\begin{quote}
\textbf{Theorem\tume\ref{main} (informal).}
Every feasible instance of the pebble motion problem
on a tree with $N$ vertices admits a solution sequence
of length $\Ord(N^2 \log N)$, constructible in time linear in its length.
\end{quote}
Together with the $\Omega(N^2)$ lower bound,
this establishes a near-quadratic worst-case complexity for trees.

We next sketch the main idea of the tree algorithm. 
Our algorithm follows a divide-and-conquer strategy based on a centroid
decomposition of the input tree.
The decomposition recursively partitions the tree into balanced subtrees.
A key ingredient is Theorem\tume\ref{sbs3}, proved in Section\tume\ref{base-cases},
which provides a near-quadratic procedure for rearranging pebbles inside a tree of size 
$\Ord(k)$, where $k:=N-n$ is the number of empty vertices.
This result serves as a fundamental building block
in the recursive construction.
Using the decomposition, we partition the tree into subtrees of size
$\Ord(k)$, obtaining $\Ord(N/k)$ such subtrees.
We then repeatedly choose a peripheral subtree $X(t)$, called the
target subtree, and place on it the pebbles prescribed by the target
configuration.
To do so, we push out of the surrounding subtrees the pebbles that do
not belong to $X(t)$.
During this process, at least one pebble destined for $X(t)$
moves along a shortest path without backtracking,
thereby making $\Theta(k)$ progress toward its destination
in each iteration (a formal proof is given in
Section\tume\ref{refinement}). 
Applying Theorem\tume\ref{sbs3} to each such subtree costs
$\Ord(k^2\log k)$ moves.
Since there are $\Ord(N/k)$ such subtrees associated with one target 
subtree, the work spent per target subtree is $\Ord(Nk\log k)$.
Repeating this over all $\Ord(N/k)$ target subtrees yields a total of
$\Ord(N^2\log k)$ moves.
Furthermore, the sequence can be produced in time $\Ord(N+L)$, where
$L$ denotes the length of the sequence.




This tree-based algorithm also forms the basis of our approach to the
\textsc{Pebble Motion Problem} on general graphs.
Given a board graph $G$, we first construct a breadth-first search
spanning tree $T_G$ of $G$ and solve the corresponding PMT instance
on $T_G$.
However, if $T_G$ contains a corridor (i.e., a maximal path of degree-$2$ 
vertices in $G$) whose length exceeds the number of empty vertices, 
the feasibility conditions for PMT may be violated.
To maintain feasibility in such cases, we slide pebbles along a
shortest cycle of $G$ that contains the corridor, thereby restoring
the conditions required for the PMT procedure to proceed.
This approach yields a unified algorithmic framework for the pebble
motion problem on arbitrary graphs.
In particular, when the number of empty vertices is constant,
every feasible instance admits a solution sequence of length
\[
\Ord(n\,\CL(G)\,\D(G)).
\]
This bound recovers the classical square-grid example discussed in
the concluding remarks of
Kornhauser, Miller, and Spirakis~\cite{715921}.
Indeed, for a $\sqrt N \times \sqrt N$ grid we have
$\CL(G)=4$ and $\D(G)=\Theta(\sqrt N)$, which yields the bound
$\Ord(N^{3/2})$.
Thus our framework provides a general algorithmic extension
of that example to broader graph classes.

Beyond these algorithmic results, our analysis also sheds light on
the classical $\Theta(N^3)$ phenomenon for general graphs.
Kornhauser, Miller, and Spirakis~\cite{715921} observed that cubic behavior arises
from highly constrained configurations in which empty vertices must
repeatedly traverse long cycles.
We prove the following structural restriction.
\begin{quote}
\textbf{Theorem\tume\ref{Restriction} (informal).}
Let $k:=N-n$ be the number of empty vertices.
If there exists a feasible instance whose shortest solution
sequence has length $\Theta(N^3)$, then the following conditions must hold.
\begin{enumerate}
\item $k=\Ord(1)$.
\item The graph $G$ contains a set $S$ of $\Theta(N)$ degree-two vertices,
each of which lies in a sequence of at least $N-n-2$ consecutive
degree-two vertices on a path of $G$.\label{CLL3}
\item For every vertex $v \in S$, the length of the shortest cycle in $G$
that passes through $v$ is $\Theta(N)$.
\end{enumerate}
\end{quote}
Thus cubic behavior can arise only under extremely restricted
structural conditions.
These conditions are satisfied by the classical constructions
of~\cite{715921}, such as a graph consisting of a long cycle
with a single pendant edge where pebbles occupy all but two
vertices.


\subsection{Summary of Contributions}

The main contributions of this paper are as follows.
\begin{enumerate}
\item We show that the classical $\Ord(N^2)$ claim for trees in
	Kornhauser–Miller–Spirakis (FOCS 1984)~\cite{715921} relies 
	on an assumption that does not hold in general, and we provide 
	a counterexample. 

\item We present the first algorithm for the pebble motion problem on
	trees that computes a solution sequence of near-quadratic length,
	namely $\Ord(N^2\log N)$, thereby nearly closing the long-standing
	quadratic--cubic gap.

\item We develop an efficient algorithm for the pebble motion problem on general 
      graphs which outputs a sequence of length $L$ in $\Ord(|E(G)| + L)$ time.
      This yields an efficient framework that generalizes the
      square-grid example mentioned in the concluding remarks of
      Kornhauser, Miller, and Spirakis~\cite{715921}, where an
      $\Ord(N^{3/2})$ bound was observed.

\item We show that $\Theta(N^3)$ solution sequences arise only under
      highly constrained structural conditions, providing a refined
      explanation of when cubic behavior occurs.
\end{enumerate}


\subsection{Organization}

The remainder of this paper is organized as follows:

Section\tume\ref{pre} presents preliminary results, including background and related work relevant to this study.  

Section\tume\ref{results} describes the main results established in this paper and explains their significance.

Section\tume\ref{base-cases} proves a key proposition corresponding to the base case of the algorithm, namely the case where the input instance has a sufficiently large number of unoccupied spaces.

Section\tume\ref{splitting-tree} presents a method for performing a centroid decomposition of the input board tree.
This decomposition reduces a general instance into smaller subinstances that satisfy the conditions established in Section\tume\ref{base-cases}.

Section\tume\ref{MainAlg} presents the Main Algorithm proposed in this paper.

Section\tume\ref{Main1Proof} provides a proof of our first main theorem, Theorem\tume\ref{main}, 
while Section\tume\ref{Main2Proof} establishes the second main theorem, Theorem\tume\ref{main2}.

Section\tume\ref{refinement} presents refinements of 
Theorem\tume\ref{main} and Theorem\tume\ref{main2}.
Specifically, Subsection\tume\ref{refinement-tree}
proves Theorem\tume\ref{refinement1}, which refines
Theorem\tume\ref{main}. 
Subsection\tume\ref{refinement-general}
establishes Theorem\tume\ref{refinement2},
refining Theorem\tume\ref{main2}, together with its
corollary Theorem\tume\ref{refinement2-general-constant-blanc}.
It also proves Theorem\tume\ref{refinement3}, from which
the structural restriction stated in
Theorem\tume\ref{Restriction} follows.

%\textcolor{red}{Section\tume\ref{ThetaNcube} gives a complete structural
%characterization of when the pebble motion problem on a
%general graph admits shortest solution sequences of length
%$\Theta(N^3)$.}

Section\tume\ref{final} concludes the paper with some final remarks.

All propositions not reproduced from cited literature are accompanied by their proofs, 
either immediately after the statement in the main text, unless stated otherwise, or in the Appendix.

%%%%%%%%%%%%%%%%%%%%%%%%%%%%%%%%%%%%%%%%%%%%%%%%%%%%%%%%%%%%%%%%%%%%%%%%%%%%%%%%
%%%%%%%%%%%%%%%%%%%%%%%%%%%%%%%%%%%%%%%%%%%%%%%%%%%%%%%%%%%%%%%%%%%%%%%%%%%%%%%%
%%%%%%%%%%%%%%%%%%%%%%%%%%%%%%%%%%%%%%%%%%%%%%%%%%%%%%%%%%%%%%%%%%%%%%%%%%%%%%%%

\section{Preliminaries}\label{pre}

The Pebble Motion Problem (PMP) and closely related models,
most notably the Multi-Agent Path Finding (MAPF) problem,
have been extensively studied in several areas of computer science,
including
discrete mathematics \cite{
loyd1959mathematical,
zbMATH03446929,
aa887662-c3e2-30b0-a475-bcd8c655206a, 
WOS:000083999100001,
WOS:000254460900008, 
zbMATH05966071,
WOS:000301306200015, 
zbMATH06431549, 
WOS:000378029000042, 
zbMATH07360934, 
WOS:001155044500001}, 
computational geometry \cite{9981807, a16110498}, 
algorithmic theory \cite{715921, 
365740, 
zbMATH04155895, 
WOS:000077614700003, 
WOS:000168095200003, 
WOS:000254460900008, 
WOS:000280552200005, 
WOS:000378029000042, 
WOS:001177663400002,  
WOS:001155044500001, 
ARDIZZONI2025104372}, 
and artificial intelligence \cite{10.5555/3463952.3464166,
10.5555/3463952.3464075, 
Huang_Li_Koenig_Dilkina_2022, 
10.1609/aaai.v39i22.34501, 
yang2025mapfworldactionworldmodel}. 
These models capture the fundamental task of reconfiguring
a collection of agents (pebbles) moving on the vertices of a graph
while avoiding collisions.

Such models arise in a variety of applications, 
such as memory allocation \cite{715921}, 
warehouse robotics \cite{WOS:000521238104083, 
WOS:000559287300006, 
LiAAAI21lifelong, 
WOS:000767667600009, 
WOS:000838038300001, 
ijcai2023p611, 
DBLP:conf/ecai/Lehoux-Lebacque24}, 
traffic control \cite{WOS:000455723500030, 
WOS:000860325800004, 
WOS:001136907801043, 
WOS:000958827700001, 
WOS:001199810300001}, 
game AI \cite{WOS:000345499600001, 
WOS:000356269500008, 
WOS:000441027500010}, 
and even in certain biomedical diagnostic systems \cite{WOS:001303734200001}.


The problem has a long history in the study of graph reconfiguration
and token swapping processes.
One of the earliest systematic studies is due to Wilson~\cite{zbMATH03446929},
who completely characterized the feasibility of instances
with a single empty vertex ($N-n=1$).
This was later extended by Kornhauser, Miller, and Spirakis~\cite{715921},
who established a polynomial-time feasibility test
for the general case $N-n \ge 2$ on undirected graphs.
In both results, feasibility can be determined in time linear in the input size.

When the board graph is directed, the problem becomes substantially more difficult.
In this setting, deciding reachability between configurations is
\NP-complete in general~\cite{WOS:001155044500001}.
Nevertheless, for feasible instances on undirected graphs
or strongly connected digraphs,
polynomial-time algorithms are known that produce solution sequences
of polynomial length~\cite{715921, WOS:000077614700003,
WOS:001177663400002, WOS:001155044500001, ARDIZZONI2025104372}.


A closely related research direction is the Multi-Agent Path Finding (MAPF) problem,
which studies coordinated motion planning for multiple agents on graphs
under various movement and optimization models.
MAPF has attracted considerable attention in recent years,
particularly in the artificial intelligence community,
due to its practical importance in large-scale robotic coordination
and automated warehouse systems.
A wide range of algorithmic techniques has been developed,
including search-based methods, conflict-based decomposition,
and compilation-based approaches
(see, e.g.,~\cite{10.5555/3463952.3464166,
Ma_2021, 
Huang_Li_Koenig_Dilkina_2022, 
DBLP:journals/corr/abs-2305-16303,
10.1609/aaai.v39i22.34501}).

Despite these algorithmic advances,
computing shortest solution sequences remains computationally difficult.
Ratner and Warmuth~\cite{zbMATH04155895} proved that,
even when $N-n=1$ and the board graph is restricted to square grid graphs,
finding an optimal solution is \NP-hard.
Polynomial-time algorithms for computing shortest sequences
are known only under strong relaxations,
such as the unlabeled setting~\cite{WOS:000254460900008}
or substantially restricted problem formulations,
such as the single-robot motion planning problem on trees with 
unlabeled obstacles~\cite{365740}.

Among the graph classes studied for PMP and MAPF,
trees remain particularly important in both theory and applications.

From an algorithmic perspective,
many MAPF algorithms employ procedures for PMT
as a core subroutine
(see, e.g.,~\cite{Khorshid2011APA, 
Krontiris2013FromFT, 
WOS:000948128106003,
WOS:001177663400002}).
Consequently, improvements in the efficiency of PMT algorithms
can directly translate into performance gains
for a wide range of MAPF solvers.

For this reason, understanding the computational complexity of PMT
remains a fundamental problem.


\section{Results}\label{results}

\subsection{Asymptotic Bounds on the Shortest Solution Length for PMT}


To understand the $\Theta(N^3)$ worst-case behavior of PMP,
we first address a long-standing open problem, 
dating back to the FOCS 1984 paper of Kornhauser, Miller, and Spirakis~\cite{715921},
namely determining the asymptotically optimal bound on the
shortest solution length for PMT and the complexity of
algorithms achieving such bounds.

\noindent
{\bf{Open Problem.}} {\it Determine the asymptotically optimal upper bound 
on the length of the shortest solution sequence for the pebble motion problem
on $N$-vertex trees.}
%%%%%%%%%%%%%%%%%%%%%%%%%%%%%%%%%%%%%%%%%%%%%%%%%%%%%%%%%%%%%%%%%%%%



We note that, in the worst case, moving $n$ pebbles from one end of an $N$-vertex path to the other may require $\Omega(Nn)$ moves. 
Moreover, the parameter $n$ itself can be as large as $\Theta(N)$ in worst-case instances. 
The order $\Omega(N^2)$ therefore provides an immediate lower bound for the above problem.

Kornhauser, Miller, and Spirakis~\cite{715921} claimed that the upper bound
for this problem is $\Ord(N^2)$.
Their argument relies on the statement that on a tree
a pebble can be moved to any vertex using at most $\Ord(N)$ moves.
However, this statement does not hold in general
(see Remark~\ref{counterexample} in the Appendix for a counterexample),
and therefore the argument does not establish the claimed bound.
This problem has therefore remained open for over four decades.
In fact, attempts to lower the classical $\Ord(N^3)$ upper bound of this length 
(c.f.~\cite{WOS:000077614700003}) have been made up to the present day
\cite{WOS:000077614700003, 
WOS:001177663400002}. 
However, to the best of our knowledge, the current best result, as shown in the 2024 
paper\cite{WOS:001177663400002}, still does not fall below the $\Ord(N^3)$ threshold.

We prove that the pebble motion problem on trees
admits a solution sequence of length $\tilde{\Ord}(N^2)$,
thereby nearly closing the long-standing quadratic--cubic gap.
We now describe our results more precisely.

To formalize the structural constraints underlying feasibility,
we introduce the notion of an isthmus. 
A path $I=v_1 \cdots v_d$ of a connected graph $G$ is called an \textit{isthmus} if 
(1) every edge of $I$ is a bridge of $G$, 
(2) every vertex of $I$ is a cutvertex of $G$, and
(3) ${\rm deg}_G (v_i)=2$ for $1 < i < d$.
An isthmus with $d$ vertices is called a $d$-isthmus.



We restrict attention to {\it feasible} instances,
namely those in which any pair of configurations can be transformed into each other.
Given an instance $(G, P)$, let $n := |P|$ be the number of pebbles placed on $G$, 
and let $N := |V(G)|$ denote the number of vertices in $G$.
We say that the instance $(G, P)$ is \textit{transitive} if, for any configuration $f$ and 
any vertex $u \in V(G)$, the pebble located at $u$ can be moved to any other vertex via a finite 
sequence of moves.
The following classical result of \cite{715921} shows that,
for any instance $(G, P)$ such that the number of unoccupied vertices, $N - n$, is at least $2$,
the following are equivalent:




\medskip
\noindent
\textbf{Theorem A} (Kornhauser, Miller, Spirakis~\cite{715921}).  {\it{
Let $2 \leq q \leq N - 1$ and assume $n + q = N$.  
Let $G$ be a connected graph with $N$ vertices that is not a cycle, 
and let $P$ be a set of $n$ pebbles placed on the vertices of $G$.  
Then the following are equivalent:
\begin{enumerate}
\item $G$ has no $q$-isthmus;
\item The instance $(G, P)$ is transitive;
\item The instance $(G, P)$ is feasible.
\end{enumerate}}}
\medskip


%%%%%%%%%%%%%%%%%%%%%%%%%%%%%%%%%%%%%%%%%%%%%%%%%%%%%%%%%%%%%%%%%%%%%%%%%%%%%%%%%%%%%%%%%%
From this characterization, it follows that no pebble can move across an isthmus
whose size is greater than or equal to the number of empty vertices $N-n$.
Consequently, in the presence of such isthmuses,
it suffices to handle the reconfiguration within the individual
feasible connected components induced by these isthmuses.
Thus, we restrict the instances of the following combinatorial problem,
the \textsc{Pebble Motion Problem on Trees}, to those satisfying these conditions.

\begin{prob}[\textsc{Pebble Motion Problem on Trees}]\label{PMPonTree} \hfill
\begin{namelist}{INSTANCE}
\item[{\it Instance}] \rm
A tree $G$ with $N$ vertices serving as the board graph (i.e., the board tree);  
an integer $k$, defined as the number of vertices contained in a longest isthmus of $G$;  
a set of pebbles $P := \{1, \ldots, n\}$ with $n \le N-1$;  
an initial configuration $f$ and a goal configuration $g$ of $P$ on $G$,  
where the condition $n + k < N$ holds.

\item[{\it Task}] \rm
Compute a sequence of moves that transforms $f$ into $g$, 
where each move shifts a pebble from its current vertex
to an adjacent unoccupied vertex of $G$.
\end{namelist}
\end{prob}

\subsection{Results on Trees}
\setcounter{thm}{0}
The following is one of our main theorems. 
\begin{thm}\label{main}
A solution sequence to the \textsc{Pebble Motion Problem on Trees} can be computed 
in time linear in the length of the sequence, which is $\Ord(nN + n^2 \log(1+\min\{n, k\})).$
\end{thm}
As a consequence of Theorem~\ref{main}, the worst-case complexity of 
the problem is $\tilde{\Ord}(N^2)$. 

In Section\tume\ref{refinement}, this theorem is further refined into the following result.
While this does not improve the worst-case complexity, it provides a finer characterization of the problem.
\begin{thm}\label{refinement1}
A solution sequence of length 
$\Ord\left(n\D(G) + \min\left\{k n \D(G),\ n^{2} \log\big(1+\min\{n, k\}\big)\right\}\right)$
to the \textsc{Pebble Motion Problem on Trees} can be computed in time 
$$\Ord\left(N+n\D(G) + \min\left\{k n \D(G),\ n^{2} \log\big(1+\min\{n, k\}\big)\right\}\right).$$
\end{thm}
The main result of \cite{WOS:001177663400002} -- which, as of 2024, remains the most recent known 
result on this problem -- is that the length of the shortest solution sequences for 
the \textsc{Pebble Motion Problem on Trees} is in $\Ord(N^2 + nkN)$.
In contrast, Theorem\tume\ref{refinement1} above constitutes a genuine improvement over their result.

By applying the new results of this paper, we obtain that the mere fact that the size $k$ of 
a longest isthmus is in $\Ord(a^{N/n})$ for some positive constant $a$ already implies that 
the length of the shortest solution sequences belongs to $\Ord(nN)$.

We conjecture the following.
\begin{conj}\label{conjecture}
There exist instances of the \textsc{Pebble Motion Problem on Trees} for 
which the length of the shortest solution sequences is $\Theta(N^2\log N)$.
\end{conj}
Indeed, considering the combinatorial structure of this problem (see Remark\tume\ref{counterexample} 
for the simplest yet most representative example), any further improvement of the complexity bound 
presented above appears to be highly challenging.


\subsection{Results on General Graphs}

We now turn to the \textsc{Pebble Motion Problem on General Graphs}.
Beyond improved upper bounds, our main contribution in this setting is
to identify structural conditions under which the classical $\Theta(N^3)$ 
behavior can arise.

Our analysis isolates the graph-theoretic features that force long
reconfiguration sequences.
In particular, we show that cubic-length solution sequences can occur
only under highly constrained configurations, and require the presence of 
long cycles interacting with a small number of empty vertices.
In the absence of such structures, significantly shorter solution
sequences necessarily exist.

Throughout this section, we restrict attention to feasible instances.
By convention, for an instance of the \textsc{Pebble Motion Problem on General Graphs},
we let the parameter $k$ denote the size of a maximum isthmus of $G$,
and set $k=1$ in the special case that $G$ is $2$-edge-connected.
We formalize this perspective through a general algorithmic framework
that extends the tree-based approach.

As shown in Theorems~\ref{main2} and \ref{refinement2-general-constant-blanc},
this yields refined upper bounds on the length of solution sequences,
expressed in terms of structural parameters such as
$\CL(G)$ and $\D(G)$.
In particular, when the number of empty vertices is constant,
every feasible instance admits a solution sequence of length
$\Ord(n\,\CL(G)\,\D(G))$,
which recovers, as a special case, the square-grid example
mentioned in the concluding remarks of
Kornhauser, Miller, and Spirakis~\cite{715921},
where it was observed that only a small fraction of instances
appear to require cubic-length solutions.

Finally, we prove that the existence of $\Theta(N^3)$-length optimal
solutions imposes stringent structural constraints on the board graph.
Specifically, such behavior can occur only if the graph contains a large
set of degree-two vertices arranged along long cycles.
This sharply restricts the structural regime in which cubic behavior can occur.
\begin{thm}\label{main2}
Let $(G, P)$ be a given instance of the \textsc{Pebble Motion Problem on General Graphs}.  
Suppose that, in addition to the instance, %a spanning breadth-first search tree $T_G$ of $G$, and 
a shortest cycle through each vertex $v \in V(G)$ (or $\emptyset$ if none exists), 
is provided. 
Let $\theta$ be an arbitrary integer such that $k \leq \theta \leq N - n - 1$.
Then a solution sequence of length 
$\Ord(n\D(G)+\frac{n^2\min\{n,\CL(G)\}}{\theta}+n^2\log(1+\min\{n, \theta\}))$ 
can be computed in time %linear in its length, which is  
$\Ord(|E(G)|+n\D(G)+\frac{n^2\min\{n,\CL(G)\}}{\theta}+n^2\log(1+\min\{n, \theta\})).$ 
\end{thm}

In a manner similar to Theorem\tume\ref{refinement1}, Theorem\tume\ref{main2} is also refined in this paper 
(for further details, see Theorem\tume\ref{refinement2} in Section\tume\ref{refinement}).
However, since the general form of the result is considerably complex, we present here only a representative case:
specifically, under the assumptions that the number of unoccupied spaces, 
$N - n$ is bounded by some constant.
Under this assumption, we establish the asymptotic upper bound on the length of the shortest solution sequences as follows:
\begin{thm}\label{refinement2-general-constant-blanc}
Assume that (1) the number of unoccupied spaces $N - n$ is constant with respect to $N$;  
and  (2) a shortest cycle passing through each vertex $v$ of $G$ (or $\emptyset$ if no such cycle exists) is 
provided in advance.
Then, a solution sequence of length 
$\Ord(n \CL(G) \D(G))$
to the \textsc{Pebble Motion Problem on General Graphs} can be 
computed in time 
$\Ord(|E(G)| + n \CL(G) \D(G)).$
\end{thm}
To find, for every vertex $v$ of a given $N$-vertex graph $G$, 
a ``shortest'' cycle (one of them) that passes through $v$, 
it would require a total computational effort of $\Ord(|E(G)||V(G)|)$. 
However, if it suffices to find, for each vertex $v$, some (not necessarily shortest) 
cycle that passes through $v$, then the required computational cost reduces to 
$\Ord(|E(G)|)$, which matches the asymptotic cost of computing
a single breadth-first search spanning tree of $G$. 
This observation leads to the corollary below.
\begin{cor}\label{main2-cor}
Let $\theta$ be an arbitrary integer such that $k \leq \theta \leq N - n - 1$.
Then a solution to the \textsc{Pebble Motion Problem on General Graphs} can be computed in 
$\Ord(|E(G)|+n\D(G)+\frac{n^3}{\theta}+n^2\log(1+\min\{n, \theta\}))$ time, and the length of the resulting 
sequence is bounded by $\Ord(n\D(G)+\frac{n^3}{\theta}+n^2\log(1+\min\{n, \theta\})).$
\end{cor}

Theorem\tume\ref{main2} and its refinement, Theorem\tume\ref{refinement2}
indicate that $\Ord(N^3)$-length optimal solutions in the
\textsc{Pebble Motion Problem on General Graphs}
can occur only under highly restricted structural conditions
on the board graph $G$.
The following theorem makes this restriction precise;
its proof will be given in Section\tume\ref{refinement}.

Let $V(G,q)$ denote the set of degree-two vertices in $G$ that lie on paths 
in $G$ and are part of a contiguous sequence of at least $q$ degree-two vertices.
\begin{thm}\label{Restriction}
Let $(G,P)$ be a feasible instance of the \textsc{Pebble Motion Problem on General Graphs}.
If the length of an optimal solution sequence is $\Theta(N^3)$, then the board graph $G$ must satisfy 
the following conditions:
\begin{enumerate}
\item $N-n=\Ord(1)$;
\item $\displaystyle \sum_{v\in V(G,N-n-2)} \CL(G,v)=\Theta(N^2)$;\label{CLL}
\item The graph $G$ contains a set $S (\subseteq V(G,N-n-2))$ of $\Theta(N)$ vertices;\label{CLL3}
\item For every vertex $v \in S$, the length of the shortest cycle in $G$
      that passes through $v$ is $\Theta(N)$.\label{CLL4}
\end{enumerate}
\end{thm}

\noindent
Note that, in Theorem\tume\ref{Restriction}, Condition\tume\ref{CLL} immediately implies 
Conditions\tume\ref{CLL3} and \ref{CLL4}.

The algorithmic bounds stated above are obtained
by constructing new algorithms for the
\textsc{Pebble Motion Problem on Trees} (PMT)
and for the \textsc{Pebble Motion Problem on General Graphs} (PMP).
The overall structure of these algorithms was outlined
in the introduction.
Their complete description will be given in
Section~\ref{MainAlg}, together with several refinements
presented in the proofs in Section~\ref{refinement}.

Finally, the techniques developed in this paper
suggest potential improvements for MAPF algorithms,
offering a more efficient alternative to existing approaches.

%%%%%%%%%%%%%%%%%%%%%%%%%%%%%%%%%%%%%%%%%%%%%%%%%%%%%%%%%%%%%%%%%%%%%%%%%%%%%%%%
%%%%%%%%%%%%%%%%%%%%%%%%%%%%%%%%%%%%%%%%%%%%%%%%%%%%%%%%%%%%%%%%%%%%%%%%%%%%%%%%
%%%%%%%%%%%%%%%%%%%%%%%%%%%%%%%%%%%%%%%%%%%%%%%%%%%%%%%%%%%%%%%%%%%%%%%%%%%%%%%%

\section{Solving the Problem on Trees with Ample Blank Space}\label{base-cases}

%%%nakami

This section describes the base cases, where the size of the board $N$ is significantly larger than the number of pebbles $n$, which are the cornerstone of the proof of Theorem\tume\ref{main}.

%In the base case where $n \ll N$, we will demonstrate that for given configurations $f, g \in (G, P)$, $f$ can be transformed into $g$ with $\Ord (nN + n^2 \log n)$ moves, and that such a sequence of moves can also be computed with the same time complexity $\Ord (nN + n^2 \log n)$. 

The main result of this section is the following theorem.

\begin{thm}\label{sbs3}
If the puzzle $(G, P)$ satisfies $N \ge 3n$, then a solution sequence of length $\Ord (n\D(G) + n^2 \log n)$ for the \textsc{Pebble Motion Problem on Trees} can be computed in time $\Ord (N + n\D(G) + n^2 \log n)$.
\end{thm}

We will show the outline of the proof of Theorem\tume\ref{sbs3}.
Let $\cF(G, P)$ be the set of all configurations on $(G, P)$.
\begin{itemize}
\item Under the condition $n \ll N$, there exists a subtree $H$ of $G$ such that the puzzle $(H, P)$ is feasible and $(H, P)$ can be divided into a family of sub-puzzles on $H$.  
\item For given configurations $f, g \in \cF(G, P)$, choose a pair of configurations $f_H, g_H \in \cF(H, P)$ such that $f$ and $g$ can be transformed to $f_H$ and $g_H$ respectively with $\Ord (n\D(G))$ moves.
\item In order to transfer $f_H$ to $g_H$ with $\Ord (n\D(G) + n^2 \log n)$ moves, we apply binary comparisons, where the puzzle $(H, P)$ is divided into sub-puzzles, in which each sub-puzzle has a half number of original pebbles, based on a divide-conquer strategy. 
\item The algorithm works similarly to the merge sorting algorithm:
 it splits the puzzle into two halves until each sub-puzzle has one pebble, and merges the pebbles to build a target configuration.  
\end{itemize}

%We propose an algorithm to move unlabeled pebbles from a set of vertices to another. 
%For a vertex $v \in V(G)$, let us denote the set of neighbors of $v$ by $N(v)$.
Firstly, we state a basic lemma concerning unlabeled pebble motion problem on trees.

\begin{la}\label{sbs0}
Let $S_1, S_2 \subset V(G)$ with $|S_1| = |S_2| = n$.
Then there exists a sequence of moves of length $\Ord (n\D(G))$ which transfers all the pebbles from $S_1$ to $S_2$, and it can be computed in time $\Ord (N + n\D(G))$.
\end{la}

Although Lemma\tume\ref{sbs0} can be derived from the proof of Theorem 2.1 of  \cite{WOS:000254460900008}, we will show a proof in the Appendix for the convenience of the reader.
Lemma\tume\ref{sbs0} is used in subsequent lemmas concerning labeled pebbles.


%Let $T$ be a tree with $N$ vertices.
%Let $P$ be a set of $n$ pebbles.
Let us denote $N + n\D(G)$ by $\alpha(G, P)$.
For a configuration $\pi \in \cF(G,P)$, the {\it support} of $\pi$ is defined as $\pi^{-1}(P)$, which is the subset of vertices of $G$ occupied by $P$.
We denote the support of $\pi$ by ${\rm sup}(\pi)$.   
For two configurations $\pi, \pi^{\ast} \in \cF(G,P)$, we say that $\pi$ and $\pi^{\ast}$ are {\it strongly related}, if there exists an algorithm which transfers $\pi$ to $\pi^{\ast}$ with a time complexity of $\Ord (\alpha(G,P))$.
For any configuration $\pi$ and for any $S \subset V(G)$ with $|S| = n$, by applying Lemma\tume\ref{sbs0}, we have a configuration $\pi^{\ast}$ such that ${\rm sup}(\pi^{\ast})$ is $S$ and $\pi^{\ast}$ is strongly related to $\pi$. 

To prove Theorem\tume\ref{sbs3}, the next two cases are particularly important.
The first key case is where $G$ has a rivet, a vertex of degree at least $3$, dividing $G$ into two parts with a suitable size.

\begin{la}\label{sbs1}
Let $G$ be a tree having a vertex $o$ with $\deg_G o \ge 3$ and let $G$ have subtrees $G_1$ and $G_2$ such that $V(G_1) \cap V(G_2) = \{ o \}$, and $|V(G_i)| \ge n + 1$ for $i=1, 2$. 
%Then the pebble motion problem on $G$ is computted in time $\Ord (Nn + n^2 \log n)$.
Then a solution sequence of length $\Ord (n\D(G) + n^2 \log n)$ for the \textsc{Pebble Motion Problem on Trees} can be computed in time $\Ord (N + n\D(G) + n^2 \log n)$.
\end{la}

The second key case for the proof of Theorem\tume\ref{sbs3} is where $G$ has a long isthmus with respect to the number of pebbles.

\begin{la}\label{sbs2}
Let $G$ be a tree with an isthmus $I$ of size at least $n$.
Let $G_1$ and $G_2$ be the two subtrees induced by $V(G) \setminus V_0$, where $V_0$ is a set of inner vertices of $I$.
%If $|V(G_i)| \ge n/2$ for $i = 1,2$, then the pebble motion problem on $G$ can be computed in time $\Ord (Nn + n^2 \log n)$.
If $|V(G_i)| \ge n/2 + 1$ for $i = 1,2$, then a solution sequence of length $\Ord (n\D(G) + n^2 \log n)$ for the \textsc{Pebble Motion Problem on Trees} can be computed in time $\Ord (N + n\D(G) + n^2 \log n)$.
\end{la}

Combining Lemma\tume\ref{sbs0}, Lemma\tume\ref{sbs1} and Lemma\tume\ref{sbs2}, we can prove Theorem\tume\ref{sbs3}. 

Before closing the section, we note the following lemma which states that starting from a given configuration, we can exchange the positions of two specified pebbles in time $\Ord (N + n\D(G))$.

Let $G$ be a tree.
We assume that the puzzle $(G, P)$ is feasible.
For a configuration $\pi \in \cF(G, P)$ and two pebbles $p, q \in P$, a configuration $\pi^\ast = \pi(p\leftrightarrow q)$ is defined as $\pi^\ast \circ \pi^{-1} (p) = q$, $\pi^\ast \circ \pi^{-1} (q) = p$ and $\pi^\ast \circ \pi^{-1} (r) = r$ for all $r \in P \setminus \{ p, q \}$.

\begin{la}\label{sbs4}
If the puzzle $(G, P)$ satisfies $N \ge 3n$, then for any $\pi \in \cF(G, P)$ and $p, q \in P$, there exists a sequence of moves of length $\Ord (n\D(G))$ which transforms $\pi$ to $\pi(p\leftrightarrow q)$, and it can be calculated in time $\Ord (N + n\D(G))$.
\end{la}

Lemma\tume\ref{sbs4} will be used in Section\tume\ref{refinement-tree}. 
 
%%%

\section{Decomposing a Tree into Smaller Subtrees}\label{splitting-tree}

To prove the main theorem, we need to decompose an instance puzzle 
into smaller sub-puzzles that satisfy the conditions of the base cases 
in the previous section. In this section we develop several structural
tools for this purpose.

The following three lemmas play complementary roles in our algorithm.
Lemma\tume\ref{centroid-splitting} provides a balanced way to split a tree
into two subtrees while sharing at most one vertex.
Lemma\tume\ref{feasible-subtree} allows us to construct subtrees of
prescribed size while preserving the bound on the maximum isthmus.
Finally, Lemma\tume\ref{bunkatsu} repeatedly applies 
Lemma\tume\ref{centroid-splitting} to obtain a global decomposition of a
tree into many smaller subtrees of bounded size, which will serve as the
domains of the subproblems solved by Theorem\tume\ref{sbs3}.

We first recall a standard property of tree centroids.
Every tree $T$ has either a vertex $v$ (called a {\it centroid vertex} of $T$)
such that the order of each connected component of $T-v$ does not exceed $(|V(T)|-1)/2$,
or an edge $e$ (called a {\it centroid edge} of $T$) such that each of the two connected
components of $T-e$ has order exactly $|V(T)|/2$.
\begin{la}[\hspace{-0.3pt}\cite{MR396308}]\label{centroid}
A centroid vertex or a centroid edge of a tree can be found in linear time.
\end{la}
Using this fact, we derive the following lemma.
\begin{la}\label{centroid-splitting}
Let $T$ be an arbitrary tree with $N$ vertices. 
Then one can find, in $\Ord(N)$ time, a pair of subtrees $(T_1, T_2)$ of $T$ such that all of 
the following three conditions hold:
\begin{enumerate}
    \item $|V(T_1)\cap V(T_2)| \leq 1$;
    \item $V(T_1) \cup V(T_2) = V(T)$;
    \item $(N+2)/3 \leq |V(T_1)| \leq |V(T_2)| \leq (N+1)/2$.
\end{enumerate}
\end{la}
This Lemma\tume\ref{centroid-splitting} is used when recursively splitting a given tree 
into two subtrees as balanced as possible, while sharing at most one vertex.

%%%%%%%%%%%%%%%%%%%%%%%%%%%%%%%%%%%%%%%%%%%%%%%%
%%%%%%%%%%%%%%%%%%%%%%%%%%%%%%%%%%%%%%%%%%%%%%%%

Consider a tree $X$ (the board tree) serving as the input board graph of a PMT. 
If the order $N$ of a tree $X$ (the board tree) is considerably larger than the number of pebbles $n = |P|$ 
to be placed on it, solving the PMT on an appropriate subtree $Y$ of $X$ can be more efficient---both 
in computation time and in the number of moves---than solving it directly on $X$. Here, $Y$ is chosen 
so that its order is at least $n + k + 1$, where $k$ is the size of the maximum isthmus of $X$, and 
the size of the maximum isthmus of $Y$ is at most $k$. The solution of the PMT on $(Y, P)$ can 
then be transferred back to the original instance $(X, P)$. 
The following Lemma\tume\ref{feasible-subtree} is used to find such a suitable $Y$ within $X$. 
In our proposed algorithm, situations repeatedly arise in which one must find a subtree $Z$ of $X$ 
that contains a specified subtree $W$, has a prescribed order, and does not contain any isthmus 
larger than the maximum isthmus of $X$, and then solve the PMT on $Z$. In these cases as well, 
this lemma is used extensively. 
\begin{la}\label{feasible-subtree}
Let $T$ be a tree such that the size of a maximum isthmus of $T$ is at most $k$. 
For any subtree $T^{\prime}$ of $T$ such that the size of a maximum isthmus of $T^{\prime}$ is also at most $k$, the following statements hold:
\begin{itemize}
\item For any natural number $d$ satisfying $|V(T^{\prime})| \leq d \leq |V(T)|$, 
we can find, in $\Ord(d)$, a $d$-vertex subtree $T^{\prime\prime}$ of $T$ such that the size of a maximum isthmus 
of $T^{\prime\prime}$ is at most $k$ and that $T^{\prime\prime}$ contains $T^{\prime}$ 
as a subgraph.
\end{itemize}
\end{la}

%%PMTの入力board graphである木（board tree）$X$の位数$N$が、$X$の上に配置するpebble達$P$の個数$n=|P|$に比べて小さいときには、
%%$X$の上のPMTを直接扱うよりも、$X$の最大isthmusのサイズ$k$と$n$の和でバウンドされた$n+k+1$の位数を持つ$X$の適切なsub tree$Y$で、
%%再び$Y$の最大isthmusのサイズも$k$を超えないものを見つけて、$(Y,P)$というインスタンスに対するPMTを解き、それを$(X,P)$のインスタンスに
%%戻してやる方が計算効率も上がるし、moveの手数も抑えることが出来る。以下の補題\ref{feasible-subtree}は、$X$からそのようなうまい$Y$を
%%見つける際に用いられる。また、我々が提案するアルゴリズムにおいては、指定された$X$の部分木$Z$を含み、かつ指定された位数を持ち、
%%更に、$X$の最大isthmusのサイズを越える最大isthmusを持たないような部分木$W$を見つけて、$W$の上のPMTを解かなければならない
%%シチュエーションが繰り返される。そうした際にも、この補題は頻繁に利用される。
%\begin{la}\label{feasible-subtree}
%Let $T$ be a tree such that the size of a maximum isthmus of $T$ is at most $k$. 
%For any subtree $T^{\prime}$ of $T$ such that the size of a maximum isthmus of $T^{\prime}$ is also at most $k$, the following statements hold.
%%最大isthmus長が$k$以下である任意の木$T$の任意の部分木$T^{\prime}$について、以下の言明が成り立つ。
%\begin{itemize}
%\item{For any natural number $d$ satisfying $|V(T^{\prime})|\leq d \leq |V(T)|$, 
%we can find, in $\Ord(d)$, a $d$-vertex subtree $T^{\prime\prime}$ of $T$ such that the size of a maximum isthmus 
%of $T^{\prime\prime}$ is at most $k$ and that $T^{\prime\prime}$ contains $T^{\prime}$ 
%as a subgraph.}
%%$T^{\prime}$の最大isthmus長も$k$以下であれば、$|V(T^{\prime})|\leq d \leq |V(T)|$を満たす任意の自然数$d$に対して、最大isthmus長が$k$以下であり$T^{\prime}$を部分グラフとして含む$T$の位数$d$の部分木が線形時間の手間で見つかる。}
%\end{itemize}
%\end{la}

%%%%%%%%%%%%%%%%%%%%%%%%%%%%%%%%%%%%%%%%%%%%%%%%
%%%%%%%%%%%%%%%%%%%%%%%%%%%%%%%%%%%%%%%%%%%%%%%%
%%%%%%%%%%%%%%%%%%%%%%%%%%%%%%%%%%%%%%%%%%%%%%%%
%%%%%%%%%%%%%%%%%%%%%%%%%%%%%%%%%%%%%%%%%%%%%%%%


By repeatedly applying Lemma\tume\ref{centroid-splitting}, we obtain the
following structural decomposition.
\begin{la}\label{bunkatsu}
For any tree $T$ with $n$ vertices and any positive integer $k \le n$,
there exists an algorithm that constructs a family $\mathcal{F}$ of at most
$4n/k$ subtrees of order $k$ such that
\[
\bigcup_{H\in\mathcal{F}} V(H)=V(T)
\]
and any two subtrees in $\mathcal{F}$ share at most one vertex.
\end{la}

\noindent
Lemma\tume\ref{bunkatsu} provides a decomposition of the board tree
into subtrees of sizes to which Theorem\tume\ref{sbs3} applies.
Our algorithm repeatedly solves PMT subproblems associated with these 
subtrees.
In the algorithm of Lemma\tume\ref{bunkatsu}, the subtrees covering $T$ are
stored and maintained in a binary tree data structure $\BT$.
Each node $x$ of $\BT$ corresponds to a subtree $X(x)$ of $T$ that
arises during the recursive decomposition.
When a subtree is split during the recursive decomposition,
the resulting pieces may share either a vertex or an edge.
We refer to this shared object as the \textit{joint} of the subtree.
Accordingly, the joint is represented by either a vertex
(a \textit{joint vertex}) or an edge (a \textit{joint edge}).
For each node $x$ of $\BT$, let $p(x)$ denote its parent node and
$c_i(x)$ $(i=1,2)$ denote its child nodes.
Each node $x$ stores the following three objects as its attributes:
\begin{itemize}
\item $X(x)$,
\item $\Cen(X(x)) := (\CenV(X(x)), \CenE(X(x)))$,
\item $\Joi(X(x)) := (\JoiV(X(x)), \JoiE(X(x)))$.
\end{itemize}
First, $X(x)$ is a subtree of $T$ determined by the inductive
procedure described below.
For a subtree $X(x)$, the pair
$\Cen(X(x)) = (\CenV(X(x)),\CenE(X(x)))$ stores its centroid
information.
If $X(x)$ has a centroid vertex, it is stored in $\CenV(X(x))$;
otherwise $\CenV(X(x)):=\emptyset$.
Similarly, if $X(x)$ has a centroid edge, it is stored in
$\CenE(X(x))$; otherwise $\CenE(X(x)):=\emptyset$.
If $x$ is a leaf node of $\BT$, we define
$\Cen(X(x)):=(\emptyset,\emptyset)$ to avoid unnecessary computation.
We now describe how the objects associated with the nodes of $\BT$
are defined.
\begin{alg-enumerate}
\item
For the root node $r$ of $\BT$, let
$X(r):=T$ and $\Joi(X(r)):=\Cen(X(r))$.
\item
Suppose that $X(x)$ is stored at a non-leaf node $x$ of $\BT$.
The objects $X(c_i(x))$ and $\Joi(X(c_i(x)))$ $(i=1,2)$ to be stored
in the child nodes $c_i(x)$ are defined inductively by the following
procedure:
\begin{alg-enumerate}
\item
Using Lemma\tume\ref{centroid-splitting}, find two subtrees
$X_1$ and $X_2$ of $X(x)$ satisfying
$V(X_1) \cap V(X_2)=\{\CenV(X(x))\}$ and
$V(X_1) \cup V(X_2)=V(X(x))$.
If $X(x)$ has a centroid edge, $X_1$ and $X_2$ are the two subtrees
obtained by removing that edge from $X(x)$.
\item
Of the two subtrees $X_1$ and $X_2$ of $X(x)$,
the one that contains $\JoiV(X(x))$ or one of the endpoints
of $\JoiE(X(x))$ is stored at the second child node $c_2(x)$ of $x$.
The other subtree is stored at the first child node $c_1(x)$.
\label{No2Child}
\item
$\Joi(X(c_2(x))):=\Joi(X(x))$.
\item
$\Joi(X(c_1(x))):=\Cen(X(x))$.
\end{alg-enumerate}
\end{alg-enumerate}
Strictly speaking, symbols such as $p^{\BT}(x)$, $c^{\BT}_i(x)$,
and $X^{\BT}(x)$ should be used to indicate that these functions are
defined on the nodes of the binary tree $\BT$.
However, to avoid notational clutter, we omit such superscripts
throughout the paper whenever no confusion arises, even when
multiple binary trees (such as $\BT_1$, $\BT_2$, etc.) are considered
simultaneously.

Then, the following lemmas hold. 
%このとき、以下の事実が成り立ちます。
\begin{la}\label{children}
For any child node $c_i(r) (i=1,2)$ of the root node $r$ of this binary tree $\BT$, the following statement holds:
%二分木$\BT$のrootノード$r$のどの子ノード$c_i(r) (i=1,2)$についても、以下の言明が成り立つ：
\begin{itemize}
\item{Regarding the leaf node $\ell(i)$ that can be reached by successively selecting only the first child of each descendant of the node $c_i(r)$, even if we delete the part $X(\ell(i))-\JoiV(X(\ell(i)))$ from the graph $X(r)$, the resulting graph $X(r)-(X(\ell(i))-\JoiV(X(\ell(i))))$ is connected. 
%$X(\ell(i))-\JoiV(X(\ell(i)))$を$X(r)$から削除して得られるグラフ$X(r)-(X(\ell(i))-\JoiV(X(\ell(i))))$は連結である。
}
\end{itemize}
\end{la}


\begin{la}\label{BT}
The effort required to create the binary tree $\BT$ from a given tree $T$ with $n$-vertices using the algorithm from Lemma\tume\ref{bunkatsu} is $\Od(n\log n)$. 
%与えられた$n$頂点の木$T$からLemma~\ref{bunkatsu}のアルゴリズムを用いて二分木リスト$\BT$を作るのに必要な手間は$\Od(n\log n)$である。
\end{la}

It is worth noting that Lemma\tume\ref{children} is essential for ensuring the correctness of the Main Algorithm, 
which will be presented in the next section. 
Lemma\tume\ref{BT}, on the other hand, provides a complexity bound for constructing the binary tree database 
used in the algorithm's complexity analysis.


\section{Main Algorithm}\label{MainAlg}


In this section, we present the main algorithm for solving the \textsc{Pebble Motion Problem on Trees} in $\Ord(nN+n^2\log(\min\{n,k\}))$ 
without any additional constraints on the instances.

%%%%%%%%%%%%%%%%%%%%%%%%%%%%%%%%%%%%%%%%%%%%%%%%
%%%%%%%%%%%%%%%%%%%%%%%%%%%%%%%%%%%%%%%%%%%%%%%%
First, we state the following lemma, which handles the case where the
maximum isthmus size is bounded by a constant.


\begin{la}\label{constant}
If the size $k$ of a longest isthmus in $G$ is bounded by a positive constant $M$, i.e., $k \leq M \in \Ord(1)$, 
then the solution sequence of length $\Ord(n\D(G)) (\subset \Ord(nN))$ to 
the \textsc{Pebble Motion Problem on Trees} can be computed in time $\Ord(N + n\D(G)) (\subset \Ord(nN))$.
\end{la}


%\begin{Proofof}{Lemma\tume\ref{constant}}
%In this case, by applying Lemma~\ref{feasible-subtree} and consecutively solving at most 
%$\Ord(N)$ feasible sub-puzzles --- each on a board tree with at most $2M+2$ vertices --- 
%along the path connecting any two pebbles, it is possible to swap their positions 
%without altering the positions of the other pebbles. Since the size of each board tree is 
%bounded by the constant $2M+2$ in these sub-puzzles, the effort to rearrange any 
%configuration on such a board tree is also in constant order $\Ord(1)$. 
%In general, after performing a suitable preprocessing of $\Ord(n\D(G))$ steps, 
%the entire configuration can be completed by repeating the operation of 
%'swapping the positions of two pebbles' at most $n-1$ times. 
%%%%%%%%%%%%%%%%%%%%%%%%%%%%%%%%%%%%%%%%%%%%%%%%%%%%%%%%%%%%%%%%%%%%%%%%%
%%%%%%%%%%%%%%%%%%%%%%%%%%%%%%%%%%%%%%%%%%%%%%%%%%%%%%%%%%%%%%%%%%%%%%%%%
%%%%%%%%%%%%%%%%%%%%%%%%%%%%%%%%%%%%%%%%%%%%%%%%%%%%%%%%%%%%%%%%%%%%%%%%%
%%%%%%%%%%%%%%%%%%%%%%%%%%%%%%%%%%%%%%%%%%%%%%%%%%%%%%%%%%%%%%%%%%%%%%%%%
%%%%%%%%%%%%%%%%%%%%%%%%%%%%%%%%%%%%%%%%%%%%%%%%%%%%%%%%%%%%%%%%%%%%%%%%%
%%%%%%%%%%%%%%%%%%%%%%%%%%%%%%%%%%%%%%%%%%%%%%%%%%%%%%%%%%%%%%%%%%%%%%%%%
%%%%%%%%%%%%%%%%%%%%%%%%%%%%%%%%%%%%%%%%%%%%%%%%%%%%%%%%%%%%%%%%%%%%%%%%%
%Moreover, in general, for any vertex $u$ of $G$ occupied by a pebble $A$, 
%and any other vertex $v$ occupied either by a different pebble $B$ or 
%by an unoccupied blank space, the positions of pebble $A$ and pebble $B$ 
%(or the blank space) can be exchanged—without altering the positions of 
%the other pebbles—by performing at most $\Ord\left(\frac{\D(G)}{k}\right)$
%sub-puzzles of size $\Ord(k)$ along a shortest path in $G$ connecting $u$ 
%and $v$, whose length is at most $\D(G)$.
%%%%%%%%%%%%%%%%%%%%%%%%%%%%%%%%%%%%%%%%%%%%%%%%%%%%%%%%%%%%%%%%%%%%%%%%%
%%%%%%%%%%%%%%%%%%%%%%%%%%%%%%%%%%%%%%%%%%%%%%%%%%%%%%%%%%%%%%%%%%%%%%%%%
%%%%%%%%%%%%%%%%%%%%%%%%%%%%%%%%%%%%%%%%%%%%%%%%%%%%%%%%%%%%%%%%%%%%%%%%%
%%%%%%%%%%%%%%%%%%%%%%%%%%%%%%%%%%%%%%%%%%%%%%%%%%%%%%%%%%%%%%%%%%%%%%%%%
%%%%%%%%%%%%%%%%%%%%%%%%%%%%%%%%%%%%%%%%%%%%%%%%%%%%%%%%%%%%%%%%%%%%%%%%%
%%%%%%%%%%%%%%%%%%%%%%%%%%%%%%%%%%%%%%%%%%%%%%%%%%%%%%%%%%%%%%%%%%%%%%%%%
%%%%%%%%%%%%%%%%%%%%%%%%%%%%%%%%%%%%%%%%%%%%%%%%%%%%%%%%%%%%%%%%%%%%%%%%%
%From these observations, it immediately follows that, in this case, even if scanning 
%the entire structure of the input tree $G$ requires up to $\Ord(N)$ time, 
%the \textsc{Pebble Motion Problem on Trees} can still be computd in time  
%$\Ord(N+n\D(G))$, and the length of the resulting 
%solution sequence is bounded by $\Ord(n\D(G)).$
%\end{Proofof}
%%%%%%%%%%%%%%%%%%%%%%%%%%%%%%%%%%%%%%%%%%%%%%%%%
%%%%%%%%%%%%%%%%%%%%%%%%%%%%%%%%%%%%%%%%%%%%%%%%%

Therefore, in the remainder of this paper, we assume $k \gg 1$ for all subsequent discussions.


Then, the outline of the remaining part of the main algorithm is as follows:
%メインアルゴリズムの概要は次の通りである。
\begin{enumerate}
\item{For the input $(G,P)$, use Lemma\tume\ref{feasible-subtree} to find a subtree $T$ of $G$ such that it has order $n+k+1$ and the size of its maximum isthmus is $k$ or less. Then, move all pebbles in the initial configuration $f$ on $G$ onto $T$ with $\Ord(N+n\D(G))$ effort and $\Ord(n\D(G))$ moves, and call this configuration, which is strongly related to $f$, $f^{\prime}$.}
%, while sequentially recording the sequence of pebble movements that occur during this process.}
%\item{入力$(G,P)$のboard tree $G$に対して、位数が$n+k+1$であり、かつ、その最大isthmus長が$k$以下であるような$G$の部分木$T$をLemma\tume\ref{feasible-subtree}を用いて見つけ、$G$上の初期配置$f$におけるすべての碁石を$T$の上に$\Ord(Nn)$の手間で移動させる（そしてこの作業で生じたPPPの移動の系列を逐次記録しておく）。}
\item{Similarly, for the goal configuration $g$ on $G$, convert it in advance to a configuration $g^{\prime}$ on $T$ that is strongly related to $g$ and matches the operations from Step\tume\ref{kiriotoshi}.
%, while also sequentially recording the sequence of pebble movements that occur during this process. 
%同様にして、$G$上のゴール配置$g$についても、$g$とstrongly relatedであり、かつステップ\ref{kiriotoshi}の操作にマッチするような$T$上の配置$g^{\prime}$に予め変換しておく。
}
\item{By using the algorithm from Lemma\tume\ref{bunkatsu}, decompose the tree $T$ into subtrees of sufficiently small order with respect to $k$.
%Lemma\tume\ref{bunkatsu}の手法を用いて、$T$を、その最大isthmus長$k$に対して十分小さな部分木達に分解する。
}\label{bunkai}
\item{
Among the subtrees obtained from the procedure in Step\tume\ref{bunkai}, there is a tree $T^{\prime}$ such that removing all of $T^{\prime}$ except at most one vertex from $T$ keeps the remaining graph connected. Therefore, starting from the configuration $f^{\prime}$, repeat the following sequence of operations: find one such tree $T^{\prime}$, place pebbles on it without empty space according to the configuration $g^{\prime}$, and then cut off that portion from the main body. 
%上記\ref{bunkai}の手続きで得られた部分木達の中には、その高々$1$点を残して他のすべてを$T$から削除しても、残されたグラフが連結に保たれるような木が常に存在するので、そのような木$T^{\prime}$上に、配置$g^{\prime}$の通りに碁石を隙間なく配置してその部分をグラフから切り落とす操作を繰り返す。
}\label{kiriotoshi}
\end{enumerate}
In Step\tume\ref{kiriotoshi} of the outline above, in order to place pebbles according to the configuration $g^{\prime}$ on $T^{\prime}$, it is first necessary to gather the set of pebbles $g^{\prime}(V(T^{\prime}))\setminus\{0\}$ that are scattered throughout the entire graph $T$ near the vicinity of $T^{\prime}$. To achieve this, the following method will be used.
%なお、上記のステップ\tume\ref{kiriotoshi}において、「$T^{\prime}$上に配置$g^{\prime}$の通りに碁石を配置する」ためには、まず$T$全体に散在している$(g^{\prime})^{-1}(V(T^{\prime}))$という碁石の集合を、$T^{\prime}$付近に集めてくる必要がある。そのためには、下記の手法を用いる。
\begin{itemize}
\item{Using the method of Lemma\tume\ref{bunkatsu} to decompose $T$, there will always be a subtree $T^{\prime\prime}$, aside from $T^{\prime}$ such that removing all but at most one vertex from $T$ keeps the remaining graph connected. On such a subtree $T^{\prime\prime}$, tightly pack any pebbles that do not belong to $g^{\prime}(V(T^{\prime}))$ and then "temporarily" cut them off from the whole. Once the pebbles are arranged on $T^{\prime}$ according to the configuration $g^{\prime}$, all parts that were "temporarily" removed in this operation will be restored.
%Lemma\tume\ref{bunkatsu}の手法を用いて分解すると、$T$には、自身の高々$1$点以外のすべてを$T$から削除しても残されたグラフが連結に保たれるような部分木が、$T^{\prime}$以外にも、常に存在する。そうした部分木$T^{\prime\prime}$の上に$(g^{\prime})^{-1}(V(T^{\prime}))$以外の碁石を隙間なく詰め込んで、$T$から一旦切り落とす（その頂点集合を灰色に着色する）という操作を繰り返す。配置$g^{\prime}$の通りに$T^{\prime}$上に碁石が配置されたら、この操作でテンポラリに削除した部分はすべて復元する。
}
\end{itemize}

%%%%%%%%%%%%%%%%%%%%%%%%%%%%%%%%%%%%%%%%%%%%%%%%
%%%%%%%%%%%%%%%%%%%%%%%%%%%%%%%%%%%%%%%%%%%%%%%%
%%%%%%%%%%%%%%%%%%%%%%%%%%%%%%%%%%%%%%%%%%%%%%%%

In the following sections of this paper, we will assume that each vertex $v$ of the input tree $G$ (and consequently all of its subgraphs) is associated with the following two attributes.
%以降の議論においては、木$G,T$の各頂点$v$に次の2つの属性を持たせる。
\begin{namelist}{********}
\item[$\Peb(v)$]{A function that returns the pebble $x (\in \{1,\ldots,n\})$ or a blank $0$ placed on vertex $v$.}
%%頂点$v$に乗っている碁石$x (\in \{1,\ldots,n\})$もしくは空白$0$を返す関数
\item[$\col(v)$]{A function that returns the color of vertex $v$. The possible colors are white, gray, or black.
%頂点$v$の色を返す関数．頂点の色は、白色、灰色、黒色のいずれか．
Regarding the meaning of the color, white indicates a vertex that has been completely removed from the graph, gray signifies a vertex that is temporarily removed during the execution of the algorithm but will be restored later, and black represents a vertex that remains in the graph.
%なお、treeの各頂点$v$の色の意味についてであるが、白色はグラフから完全に削除された頂点であることを、灰色はアルゴリズムの遂行上一旦はグラフから削除されるがのちに復元される頂点であることを、黒色はグラフに残された頂点であることを示している。
}
\end{namelist}
In other words, the graphs are treated not just as simple graphs, but also as objects that also contain information about the configurations of pebbles on them.
%すなわち、$G$や$T$は単なるグラフであるのみならず、その上の碁石の配置の情報も併せ持つオブジェクトとして扱います。
In this algorithm, the shapes such as the graphs $G$ and $T$, as well as the graphs $X(x)$ stored in each node $x$ of $\BT$ created by the operations in Step\tume\ref{bunkai} of the outline above, do not change; however, the values of the functions $\Peb(v)$ and $\col(v)$ at each vertex $v$ of the graphs continuously change.
%このアルゴリズムにおいては、$G$や$T$あるいは$\BT$の各ノード$x$に格納された$X(x)$などの形状は変化しないが、treeの各頂点$v$の持つ関数$\Peb(v), \col(v)$の値は刻々と変化していく。


%%%%%%%%%%%%%%%%%%%%%%%%%%%%%%%%%%%%%%%%%%%%%%%%
%%%%%%%%%%%%%%%%%%%%%%%%%%%%%%%%%%%%%%%%%%%%%%%%
%%%%%%%%%%%%%%%%%%%%%%%%%%%%%%%%%%%%%%%%%%%%%%%%


According to Theorem\tume{A}, if the size of the pebble set $P$ is $n$ and the
size of the maximum isthmus of a tree $T$ is at most $k$, then the
pebble motion problem on $(T,P)$ is feasible whenever $|V(T)| \geq n+k+1$.
Therefore, given an instance $(G,P)$, our algorithm first uses
Lemma\tume\ref{feasible-subtree} to find a subtree $T$ of $G$ whose order
is $n+k+1$ and whose maximum isthmus size is at most $k$.
Next, we transform the goal configuration $g$ on $G$ into a
``suitable'' configuration $g'$ on $T$ that is strongly related to $g$.
This transformation requires at most $\Ord(Nn)$ moves.
The procedure for constructing $g'$ is given by the algorithm
$\Alg.\PacCon(G,T,\BT,g)$ below.
%According to Theorem A, if the size of the pebble set $P$ is $n$ and the size of the maximum isthmus of the graph $T$ is at most $k$, then if the order of $T$ is at most $n + k + 1$, it is guaranteed that the pebble motion problem on $(T, P)$ is feasible.
%%the size of the set of pebbles $P$が$n$である場合、Theorem Aによれば、
%%最大isthmus長が$k$以下でありさえすれば、board tree $T$の位数は、
%%高々$n+k+1$あれば、$(T,P)$の上のpebble motion problemはfeasibleとなる。
%Therefore, the algorithm presented in this paper first uses Lemma\tume\ref{feasible-subtree} to find a subtree $T$ of the tree $G$ from the input $(G, P)$ such that its order is $n + k + 1$ and its size of a maximum isthmus is at most $k$.
%%よって、本論文に示すアルゴリズムでは、まず入力$(G,P)$のboard tree $G$に対して、
%%位数が$n+k+1$であり、かつ、その最大isthmus長が$k$以下であるような$G$の
%%部分木$T$をLemma\tume\ref{feasible-subtree}を用いて見つける。
%Furthermore, we rearrange the goal configuration $g$ on $G$ into a "suitable" configuration $g^{\prime}$ on $T$ that is strongly related to $g$ with a time complexity and number of moves of at most $\Ord(Nn)$. The algorithm used for this purpose is $\Alg.\PacCon(\BT, g)$ shown below.
%%さらに、$(G,P)$のgoal configuration $g$を$\Ord(Nn)$の手数で（$g$とstrongly relatedな）$T$上の「適切な」配置$g^{\prime}$に変換する。
%%そのために利用するアルゴリズムが、以下に示す$\Alg.$詰込配置$(\BT, g)$である。
\begin{alg}[$\Alg.\PacCon(G, T, \BT, g)$]\hfill
\begin{namelist}{OUTPUT}
\item[\sl Input]
A tree $G$ and its subtree $T$, the binary tree $\BT$ storing the
subtrees covering $T$ obtained by Lemma\tume\ref{bunkatsu},
and a configuration $g$ on $(G,P)$.

\item[\sl Output]
A configuration $g'$ on $(T,P)$ and a sequence of moves
$\phi[r]\cdot\psi$ transforming $g$ into $g'$.
\end{namelist}

\begin{alg-enumerate}

\item
Set $\Peb := g$.
For each $u\in V(T)$ set $\col(u):=\mathrm{black}$,
and for each $v\in V(G)\setminus V(T)$ set $\col(v):=\mathrm{white}$.

\item
Rearrange the configuration $g$ on $(G,P)$ into a configuration on
$(T,P)$ that is strongly related to $g$, and record the performed
sequence of moves as $\psi$.

\item
Call $\Alg.\Pac(T,\BT,r)$ for the root node $r$ of $\BT$.

\item
Return the move sequence $\phi[r]\cdot\psi$ and the configuration
$g' := \phi[r]\cdot\psi(g)$ on $(T,P)$.

\end{alg-enumerate}
\end{alg}
%\begin{alg}[$\Alg.\PacCon(G, T, \BT, g)$]\hfill
%\begin{namelist}{OUTPUT}
%\item[\sl Input]{A tree $G$ and its subtree $T$, the binary tree $\BT$ that stores the set of subtrees covering the tree $T$ obtained from the algorithm stated in Lemma\tume\ref{bunkatsu}, and a configuration $g$ on the pair $(G,P)$.}
%\item[\sl Output]{A configuration $g^{\prime} $ on $(T,P)$ and a sequence of moves $\phi[r]\cdot\psi$ to rearrange $g$ into $g^{\prime}$.}
%\end{namelist}
%\begin{alg-enumerate}
%\item{Set $\Peb:=g$. Set $\col(u):={\textrm{black}}$ for each $u\in V(T)$. Set $\col(v):={\textrm{white}}$ for each $v\in V(G) \setminus V(T)$.}
%\item{Rearrange the configuration $g$ on $(G,P)$ to some configuration on $(T,P)$ which is strongly related to $g$, and remember the sequence of moves performed here as $\psi$.}
%%配置$g$のすべての碁石を$T$上に移動させると共に、ここで行った一連のmoveの系列を$\psi$として記憶せよ}
%\item{Call $\Alg.\Pac(T, \BT, r)$ for the root node $r$ of $\BT$. }
%%$\BT$のrootノード$r$に対して、サブルーチン$\Alg.$詰込$(\BT,r)$を呼ぶ}
%\item{Return the sequence of moves $\phi[r]\cdot\psi$ and the configuration $g^{\prime} := \phi[r] \cdot \psi(g)$ on $(T,P)$.}
%%moveの系列$\phi[r]\cdot\psi$を出力して停止する}
%\end{alg-enumerate}
%\end{alg}
%%%
The procedure $\Alg.\Pac$ recursively packs pebbles toward the first
child subtrees of $\BT$ so that the configuration becomes compact
with respect to the decomposition structure.
\begin{alg}[$\Alg.\Pac(T,\BT,x)$]\hfill
\begin{namelist}{OUTPUT}

\item[\sl Input]
A tree $T$, the binary tree $\BT$ storing the set of subtrees covering
$T$ obtained from Lemma\tume\ref{bunkatsu},
and a node $x$ of $\BT$.

\item[\sl Output]
A sequence of moves $\phi[x]$.

\end{namelist}

\begin{alg-enumerate}

\item
If $x$ is a leaf node of $\BT$, return the empty sequence.

If there is no pebble on $X(x)$, or
if every vertex of $X(x)-\JoiV(X(x))$ is occupied,
return the empty sequence.

\item
Move pebbles into vertices of 
$X(c_1(x))-\JoiV(X(c_1(x)))$ 
until either
\begin{enumerate}
\item all vertices of $X(c_1(x))-\JoiV(X(c_1(x)))$ are occupied, or  
\item no pebble remains in $X(x)$ that can be moved there.
\end{enumerate}
Record the performed moves as $\phi[x]$.

\item
Compute $\phi[c_1(x)] := \Alg.\Pac(T,\BT,c_1(x))$
and $\phi[c_2(x)] := \Alg.\Pac(T,\BT,c_2(x))$.

\item
Update
\[
\phi[x] := \phi[c_2(x)]\cdot\phi[c_1(x)]\cdot\phi[x]
\]
and return $\phi[x]$.

\end{alg-enumerate}
\end{alg}

%\begin{alg}[$\Alg.\Pac(T, \BT, x)$]\hfill
%\begin{namelist}{OUTPUT}
%\item[\sl Input]{A tree $T$, the binary tree $\BT$ that stores the set of subtrees covering the tree $T$ obtained from the algorithm stated in Lemma\tume\ref{bunkatsu},
%and a node $x$ of $\BT$.}
%\item[\sl Output]{A sequence of moves $\phi[x]$.}
%\end{namelist}
%\begin{alg-enumerate}
%\item{If the node $x$ is a leaf node of $\BT$, else if there is no pebble on $X(x)$, else if there is no unoccupied vertex in $X(x)-\JoiV(X(x))$, then return a zero-length sequence of moves (i.e., an identity transformation) as $\phi[x]$.}
%%ノード$x$が$\BT$のleafノードであるか、$X(x)-\JoiV(X(x))$上に空白がないか、あるいは$X(x)$上に碁石が一つも乗っていない場合には、$\phi[x]$をmoveの長さ$0$の系列（何もしない恒等写像）として返して停止する
%\item{Fill the pebbles in the current configuration into the portion of the tree $X (c_1(x))$ except for $\JoiV(X(c_1(x)))$ % (i.e. $X(c_1(x))-\JoiV(X(c_1(x)))$) 
%as much as possible, and record the sequence of moves performed here as $\phi[x]$.}
%%現在の碁石の配置に対して、碁石を$\BT$のノード$x$の第1子$c_1(x)$に格納されている木$X(c_1(x))$の$\JoiV(X(c_1(x)))$以外の部分($X(c_1(x))-\JoiV(X(c_1(x)))$)に詰められるだけ詰め込むと共に、ここで行った一連のmoveの系列を$\phi[x]$として記憶する
%\item{Call $\Alg.\Pac(T, \BT, c_1(x))$ and $\Alg.\Pac(T, \BT, c_2(x))$.}
%%\item{$\Alg.$詰込$(\BT, c_1(x))$と$\Alg.$詰込$(\BT, c_2(x))$を呼ぶ}
%\item{Update $\phi[x]:=\phi[c_2(x)]\cdot\phi[c_1(x)]\cdot\phi[x]$ and return $\phi[x]$.}
%%\item{$\phi[x]:=\phi[c_2(x)]\cdot\phi[c_1(x)]\cdot\phi[x]$に更新して、$\phi[x]$を返して停止する}
%\end{alg-enumerate}
%\end{alg}
\noindent
In the main algorithm, the initial configuration $f$ is first rearranged to the configuration 
$g^{\prime} = \phi[r] \cdot \psi(g)$, which is the output of $\Alg.\PacCon(\BT, g)$.  
Then, by applying the inverse sequence of moves $\psi^{-1}\cdot\phi[r]^{-1}$ from $g^{\prime}$ to $g$, 
we obtain $g$ as $\psi^{-1} \cdot \phi[r]^{-1}(g^{\prime})$. 
%メインアルゴリズムでは、まず初期配置$f$を$g^{\prime}:=\phi[r]\cdot\psi(g)$に配置換えしてから、
%それを$\psi^{-1}\cdot\phi[r]^{-1}$というmoveの逆系列に掛けることにより、$g=\psi^{-1}\cdot\phi[r]^{-1}(g^{\prime})$を実現する。

%%%%%%%%%%%%%%%%%%%%%%%%%%%%%%%%%%%%%%%%%%%%%%%%
%%%%%%%%%%%%%%%%%%%%%%%%%%%%%%%%%%%%%%%%%%%%%%%%

After reducing the instance on $(G,P)$ to the one on $(T,P)$, we will use the following algorithm $\Alg.\CutOff(G,T,\BT,f,g^{\prime})$ to perform the following steps repeatedly: 1) find a subtree at the tip of tree $T$ that, when removed, does not cause the remaining parts to become disconnected (we will refer to this subtree as the {\it target} subtree), 2) fill the target subtree with pebbles according to the configuration $g^{\prime}$, and 3) cut that portion off from the whole. 
%$(T,P)$上のpebble motion problemに落とし込んだ後は、下記の$\Alg.$切り落とし$(\BT,f,h)$というアルゴリズムを用いて、$g^{\prime}$の配置通りに、$T$の端にあるsubtree（taget subtreeと呼ぶ）にpebbleを詰め込んで切り落とすという作業を繰り返す。
\begin{alg}[$\Alg.\CutOff(G,T,\BT,f,g^{\prime})$]\hfill
\begin{namelist}{OUTPUT}
\item[\sl Input]{An instance tree $G$ of the \textsc{Pebble Motion Problem on Trees}, 
its subtree $T$ of order $n+k+1$ whose maximum isthmus size is at most $k$, 
the binary tree $\BT$ storing the family of subtrees covering $T$ obtained in Lemma~\ref{bunkatsu}, 
a configuration $f$ on $(G,P)$, and an output configuration $g'$ on $(T,P)$ of $\Alg.\PacCon(G,T,\BT,g)$.}

\item[\sl Output]{A sequence of moves $\sigma$ that rearranges $f$ into $g'$.}
\end{namelist}

\begin{alg-enumerate}

\item
Set $\Peb := f$.  
Set $\col(u) := \mathrm{black}$ for each $u \in V(T)$, and  
$\col(v) := \mathrm{white}$ for each $v \in V(G)\setminus V(T)$.

\item
Rearrange the configuration $f$ on $G$ to a configuration on $T$ 
that is strongly related to $f$, and store the sequence of moves as $\sigma$.

\item
If $|V(X(r))|\le 3(k+1)/2$ holds for the root node $r$ of $\BT$, then  
compute the sequence of moves $\sigma_0$ using Theorem\tume\ref{sbs3} so that the remaining configuration becomes $g'$.  
Update $\sigma := \sigma_0\cdot\sigma$ and return $\sigma$.  
(\#\# Note that $\sigma(f)=g'$.)  
\label{initial-check}

\item
Set $t:=c_1(r)$.  
While $t$ is not a leaf node of $\BT$, update $t:=c_1(t)$.

\item
Call $\Alg.\Extraction(\BT,g',t)$.

\item
Update $\sigma := \sigma_t\cdot\sigma$.

\item
Set $\col(v):=\mathrm{white}$ for each 
$v\in V(X(t)-\JoiV(X(t)))$.  
\label{t-eliminate}

\item
Set $x:=p(t)$ and $y:=c_2(x)$.  
(\#\# Note that $t=c_1(x)$.)

\item
If $X(r)-(X(t)-\JoiV(X(t)))$ would contain an isthmus of size at least $k+1$,  
choose an edge $uv$ from the set
\[
\{uv\mid u\in V(X(t)-\JoiV(X(t))),\; v\in V(X(y))\},
\]
change $\col(v)$ (which was set to white in Step~\ref{t-eliminate}) back to black,  
find a leaf node $w$ of $\BT$ such that $v\in V(X(w))$, and update
\[
X(w):=X(w)\cup uv .
\]
\label{hige}

\item
Delete the node $t$ from $\BT$ and update $x:=y$.  
\label{del-t}

\item
If $x$ is the root node of $\BT$, then set
\[
\Joi(X(x)) := \Cen(X(x)).
\]

\item
Go to Step~\ref{initial-check}.

\end{alg-enumerate}
\end{alg}
%\begin{alg}[$\Alg.\CutOff(G,T,\BT,f,g^{\prime})$]\hfill
%\begin{namelist}{OUTPUT}
%\item[\sl Input]{An instance tree $G$ of the \textsc{Pebble Motion Problem on Trees}, and its subtree $T$ of order $n + k + 1$ with the maximum isthmus size less than or equal to $k$, the binary tree $\BT$ that stores the set of subtrees covering the tree $T$ obtained from the algorithm stated in Lemma\tume\ref{bunkatsu}, a configuration $f$ on $(G,P)$, and an output configuration $g^{\prime}$ on $(T,P)$ of $\Alg.\PacCon(G, T, \BT, g)$.}
%\item[\sl Output]{A sequence of moves $\sigma$ to rearrange $f$ into $g^{\prime}$.}
%\end{namelist}
%\begin{alg-enumerate}
%\item{Set $\Peb:=f$. Set $\col(u):={\textrm{black}}$ for each $u\in V(T)$. Set $\col(v):={\textrm{white}}$ for each $v\in V(G) \setminus V(T)$.}
%%グラフの頂点のPeb関数をfにセットし、G-Tの各頂点を白く、Tの各頂点を黒く彩色せよ。
%\item{Rearrange the configuration $f$ on $G$ to some configuration on $T$ which is strongly related to $f$, and remember the sequence of moves performed here as $\sigma$.}
%%\item{$N$頂点からなるinitial board tree $G$上のパズルの初期配置$f$における碁石のすべてを位数$n+k+1$のボードグラフ$T$の上に集めると共に、当該処理に要したmoveの系列を$\sigma$として記憶する}
%\item{If $|V(X(r))| \leq 3(k+1)/2$ holds for the root node $r$ of $\BT$, then by using Theorem\tume\ref{sbs3}, 
%calculate the sequence of moves $\sigma_0$ to rearrange the last remaining parts according to the configuration $g^{\prime}$, 
%update $\sigma:=\sigma_0\cdot\sigma$, and return the sequence of moves $\sigma$ that have been performed so far. (\#\# Note that $\sigma(f)=g^{\prime}$.) }\label{initial-check}
%%\label{initial-check}\item{$\BT$のrootノード$r$に格納されている木$X(r)$の位数が$3(k+1)/2$以下ならば、
%%中上川先生のアルゴリズムを使って、残された$T$の部分パズルを配置$h$に直すmoveの系列$\sigma_0$を計算し、
%%これまでに施した一連のmoveの系列$\sigma:=\sigma_0\cdot\sigma$を返して停止する($\sigma(f)=g^{\prime}$であることに注意)}\label{initial-check}
%\item{Find the leaf node $t$ that is reached by traversing the descendants of the first child $c_1(r)$ of the root node $r$ of $\BT$ by selecting only its first child node from generation to generation.}
%%二分木$\BT$のrootノード$r$の第1子$c_1(r)$の子孫を代々その第1子ノードのみを選択して辿ることにより行き着くleafノード$t$を見つける}
%\item{Call $\Alg.\Extraction(\BT,g^{\prime},t)$.}
%%\item{サブルーチン$\Alg.\Extraction(\BT,h,t)$を呼ぶ}
%\item{Update $\sigma:=\sigma_t\cdot\sigma$.}
%%\item{$\sigma:=\sigma_t\cdot\sigma$に更新する}
%\item{Set $col(v)={\textrm{white}}$ for each $v \in V(X(t)-\JoiV(X(t)))$.}\label{t-eliminate}
%%\item{$X(t)-\JoiV(X(t))$の頂点を白く着色する}\label{t-eliminate}
%\item{Set $x:=p(t), y:=c_2(x)$. (\#\# Note that $t=c_1(x)$.)}
%%\item{$x:=p(t), y:=c_2(x)$と置く（このとき$t=c_1(x)$であることに注意）}
%\item{If $X(r)-(X(t)-\JoiV(X(t)))$ will have an isthmus of size $k+1$ or more, to prevent this situation, choose an appropriate edge $uv$ in the edge set $\{uv \mid u \in V(X(t)-\JoiV(X(t))), v \in X(y)\}$, change the color $\col(v)$ of $v$ that was colored white in Step\tume\maru{\ref{t-eliminate}} back to black, find a leaf node $w$ of $\BT$ such that $X(w)$ contains $v$, and update $X(w) := X(w) + uv$. }\label{hige}
%%\item{$X(r)-(X(t)-\JoiV(X(t)))$が長さ$k+1$以上のishmusを持ってしまう場合には（それを防ぐために）、辺集合$\{uv \mid u \in V(X(t)-\JoiV(X(t))), v \in X(y)\}$内の辺$uv$を適当に選び、
%%ステップ\maru{\ref{t-eliminate}}で白く着色した頂点$v$を黒色に戻し、$y$の子孫を代々その第二子のみ辿って行き着く$\BT$のleafノード$w$に対して$X(w):=X(w) \cup uv$に更新する}
%\item{Delete the node $t$ from $\BT$ and update $x:=y$. }\label{del-t}
%%\item{$\BT$から$t$を削除し、$x:=y$に更新する}
%\item{If $x$ is the root node of $\BT$, then update $\Joi(X(x)):=\Cen(X(x))$.}
%%\item{$x$が$\BT$のrootノードならば、$\Joi(X(x)):=\Cen(X(x))$とする}
%\item{Go to Step\tume\maru{\ref{initial-check}}.}
%%\item{ステップ\maru{\ref{initial-check}}に戻る}
%\end{alg-enumerate}
%\end{alg}

\noindent
In the algorithm $\Alg.\CutOff(G,T,\BT,f,g^{\prime})$, it is necessary to gather the set of pebbles $g^{\prime}(V(X(t)))\setminus\{0\}$, which are scattered throughout the entire $T$, onto the target subtree $X(t)$. For this purpose, the following $\Alg.\Extraction(\BT,h,t)$ is provided. 
%アルゴリズム$\Alg.\CutOff(G,T,\BT,f,h)$においては、$T$全体に散在している$(g^{\prime})^{-1}(V(X(t)))$という碁石の集合を、
%target subtreeである $X(t)$ 上に集めてくる必要がある。そのためのアルゴリズムとして、$\Alg.\Extraction(\BT,h,t)$が用意されている。
\begin{alg}[$\Alg.\Extraction(T, \BT,g^{\prime},t)$]\hfill
\begin{namelist}{OUTPUT}
\item[\sl Input]{A tree $T$ of order $n + k + 1$ with maximum isthmus size at most $k$,
a binary tree $\BT$ storing subtrees of $T$, an output configuration $g^{\prime}$ on $(T,P)$ of $\Alg.\PacCon(G,T,\BT,g)$,
and a leaf node $t$ of $\BT$.}

\item[\sl Output]{A sequence of moves $\sigma_t$.}
\end{namelist}

\begin{alg-enumerate}

\item Copy the binary tree $\BT$ and denote it by $\BT2$.

\item If $|V(X(r))| \le 3(k+1)/2$ for the root $r$ of $\BT2$,  
apply Theorem\tume\ref{sbs3} to place the pebbles
$g^{\prime}(V(X(t)))\setminus\{0\}$ on $X(t)$ according to $g^{\prime}$.
After that, change the color of every gray vertex of $T$ back to black,
and return the sequence of moves $\sigma_t$ obtained so far.
\label{initial-check2}

\item Starting from the second child $c_2(r)$ of the root $r$ of $\BT2$,
repeatedly follow the first child until a leaf node $\ell$ is reached.

\item By Lemma\tume\ref{feasible-subtree}, find a subtree $H(\ell)$ of $T$
such that

\[
|V(H(\ell))|
=
|V(X(t))| + |V(X(\ell))| + k + 1 ,
\]

$X(\ell)$ is a subgraph of $H(\ell)$,
the maximum isthmus size of $H(\ell)$ is at most $k$,
and all vertices of $H(\ell)$ are black.
Move the $k+1$ empty spaces into $V(H(\ell))$
while recording the sequence of moves.
\label{mini-puzzle}

\item Apply Theorem~\ref{sbs3} to $H(\ell)$ and fill
$X(\ell)-\JoiV(X(\ell))$
with arbitrary pebbles that are not contained in
$g^{\prime}(V(X(t)))$,
leaving no empty space there.
Record the sequence of moves.
\label{mini-puzzle-finish}

\item Set $x := p(\ell)$ and let $y$ be the child of $x$ different from $\ell$.

\item Change the color of every vertex in
$V(X(\ell)-\JoiV(X(\ell)))$
to gray.

\item If the tree
$X(r)-(X(\ell)-\JoiV(X(\ell)))$
contains an isthmus of size at least $k+1$,
then choose an edge
$uv$ with
$u \in V(X(\ell)-\JoiV(X(\ell)))$
and $v \in X(u)$.
Change the color of $u$ from gray back to black,
find a leaf node $w$ of $\BT$
such that $X(w)$ contains $v$,
and update

\[
X(w) := X(w) + uv .
\]
\label{hige2}

\item Delete the node $\ell$ from $\BT2$
and update $x := y$.
\label{del-ell}

\item If $x$ is the root of $\BT2$,
set
\[
\Joi(X(x)) := \Cen(X(x)).
\]

\item Go back to Step~\ref{initial-check2}.

\end{alg-enumerate}
\end{alg}
%\begin{alg}[$\Alg.\Extraction(T, \BT,g^{\prime},t)$]\hfill
%\begin{namelist}{OUTPUT}
%\item[\sl Input]{A tree $T$ of order $n + k + 1$ with a maximum isthmus size less than or equal to $k$,
%a binary tree $\BT$ that stores subtrees of $T$, an output configuration $g^{\prime}$ on $(T,P)$ of $\Alg.\PacCon(G, T, \BT, g)$, 
%and a leaf node $t$ of $\BT$}
%\item[\sl Output]{A sequence of moves $\sigma_t$.}
%\end{namelist}
%\begin{alg-enumerate}
%\item{Copy the binary tree $\BT$ as a binary tree named $\BT2$.}
%%\item{$\BT$を$\BT2$にコピーする}
%\item{If the number of vertices in the tree $X(r)$ stored in the root node $r$ of the binary tree $\BT$ is at most $3(k+1)/2$, use Theorem\tume\ref{sbs3} to place the pebbles 
%$g^{\prime}(V(X(t)))\setminus\{0\}$ on the tree $X(t)$ in accordance with the configuration $g^{\prime}$. Then, change the color of all the gray vertices in $T$ back to black, and return the sequence of moves $\sigma_t$ that have been applied so far. }\label{initial-check2}
%%\item{$\BT2$のrootノード$r$に格納されている木$X(r)$の頂点数が$3(k+1)/2$以下であれば、中上川先生のアルゴリズムを使って、ターゲットノード$t$に格納されている木$X(t)$上に配置すべき碁石達（$P(t)$の部分集合）を配置$h$の通りに配置し、$T$の灰色の頂点達の色をすべて黒色に戻して、これまでに施した一連のmoveの系列$\sigma_t$を返して停止}\label{initial-check2}
%\item{Find the leaf node $\ell$ by traversing the descendants of the second child $c_2(r)$ of the root node $r$ of $\BT2$, selecting only the first child node at each level of descent.}
%%\item{二分木$\BT2$のrootノード$r$の第2子$c_2(r)$の子孫を代々その第1子ノードのみを選択して辿ることにより行き着くleafノード$\ell$を見つける}
%\item{By using Lemma\tume\ref{feasible-subtree}, find a subtree $H(\ell)$ of $T$ that includes the tree $X(\ell)$ as its subgraph, has an order of $|V(X(t))| + |V(X(\ell))| + k+1$, and its maximum isthmus size of $k$ or less, consisting solely of black vertices. Then, move the $k+1$ empty spaces over $V(H(\ell))$ while recording the sequence of moves required for this task.}\label{mini-puzzle}
%%\item{Lemma\tume\ref{feasible-subtree}を用いて、木$X(\ell)$を含む位数$|V(X(t))|+|V(X(\ell))|+k-1$かつ最大isthmus長が$k$以下の黒い頂点からなる連結な部分グラフを同定し、その頂点集合上に$k$個の空白を移動させると同時に、そのmoveの系列を記録する。}\label{mini-puzzle}
%\item{By applying Theorem\tume\ref{sbs3} to the region $H(\ell)$ identified in Step\tume\maru{\ref{mini-puzzle}}, fill $X(\ell) - \JoiV(X(\ell))$ with pebbles (anything that is not included in $g^{\prime}(V(X(t)))$) without any unoccupied space, and record the sequence of moves required for this task.}\label{mini-puzzle-finish}
%%\item{ステップ\maru{\ref{mini-puzzle}}で確保した領域$H(\ell)$に中上川先生のアルゴリズムを適用し、$h(V(X(t)))$に含まれない碁石（であればなんでもよい）を$X(\ell)-\JoiV(X(\ell))$に隙間なく詰め込むと同時に、そのmoveの系列を記録する}\label{mini-puzzle-finish}
%\item{Set $x:=p(\ell)$ and set $y$ to be the child node of $x$ that is not $\ell$.}
%%\item{$x:=p(\ell)$と置き、$x$の$\ell$ではない方の子ノードを$y$と置く}
%\item{Change the colors of the vertices $V(X(\ell)-\JoiV(X(\ell)))$ to all gray.}
%%\item{$X(\ell)-\JoiV(X(\ell))$の頂点達を灰色に着色する}
%\item{If $X(r) - (X(\ell) - \JoiV(X(\ell)))$ has an isthmus of size $k+1$ or more, to prevent this situation, select an edge $uv$ from the set $\{ uv \mid u \in V(X(\ell) - \JoiV(X(\ell))), v \in X(u) \}$, change the color $\col(u)$ of $u$ from gray back to black, find a leaf node $w$ of $\BT$ such that $X(w)$ contains $v$, and update $X(w) := X(w) + uv$. }\label{hige2}
%%\item{$X(r)-(X(\ell)-\JoiV(X(\ell)))$が長さ$k+1$以上のishmusを持ってしまう場合には（それを防ぐために）、辺集合$\{uv \mid u \in V(X(\ell)-\JoiV(X(\ell))), v \in X(u)\}$内の一辺$uv$を適当に選び、$u$の色を灰色から黒に戻し、$y$の子孫を代々その第二子のみ辿って行き着く$\BT$のleafノード$w$に対して$X(w):=X(w) \cup uv$という更新を行う。}
%\item{Delete the node $\ell$ from $\BT2$, and update $x:=y$.}\label{del-ell}
%%\item{$\BT2$から$\ell$を削除し、$x:=y$に更新する}
%\item{If $x$ is the root node of $\BT2$, then set $\Joi(X(x)):=\Cen(X(x))$.}
%%\item{$x$が$\BT2$のrootノードならば、$\Joi(X(x)):=\Cen(X(x))$とする}
%\item{Go to Step\tume\maru{\ref{initial-check2}}.}
%\end{alg-enumerate}
%\end{alg}

The following Lemma\tume\ref{subtrees-containing-v} is implicitly invoked in the complexity estimation 
within the proof of Theorem\tume\ref{main}, insofar as its computational contribution is asymptotically 
negligible compared to the other components.
\begin{la}\label{subtrees-containing-v}
Step\tume\maru{\ref{hige}} and Step\tume\maru{\ref{del-t}} of $\Alg.\CutOff(G,T,\BT,f,g^{\prime})$, as well as Step\tume\maru{\ref{hige2}} and Step\tume\maru{\ref{del-ell}} of 
$\Alg.\Extraction(T, \BT,g^{\prime},t)$, can each be achieved in $\Ord(k)$. 
\end{la}

\begin{la}\label{mini-jyunbi}
The total cost of Step\tume\maru{\ref{mini-puzzle}} of $\Alg.\Extraction(T,\BT,g^{\prime},t)$ is in $\Ord(k(k+n))$.  
%上記のサブルーチン $\Alg.\Extraction(T,\BT,h,t)$のステップ\maru{\ref{mini-puzzle}}にかかるトータルの手間は、$\Ord(k(k+n))$である。
\end{la}


%%%%%%%%%%%%%%%%%%%%%%%%%%%%%%%%%%%%%%%%%%%%%%%%
%%%%%%%%%%%%%%%%%%%%%%%%%%%%%%%%%%%%%%%%%%%%%%%%
%%%%%%%%%%%%%%%%%%%%%%%%%%%%%%%%%%%%%%%%%%%%%%%%
%%%%%%%%%%%%%%%%%%%%%%%%%%%%%%%%%%%%%%%%%%%%%%%%
%%%%%%%%%%%%%%%%%%%%%%%%%%%%%%%%%%%%%%%%%%%%%%%%

We now present the main algorithm, which combines the procedures described above.
%以上に述べた一連のアルゴリズム達の機能を集積することにより、下記のmain algorithmが構成される。
\begin{alg}[{\bf Main Algorithm:} $\Alg.\Main((G,P), f, g)$]\hfill
\begin{namelist}{OUTPUT}
\item[\sl Input]{A tree $G$ with $N$ vertices, the size $k$ of a longest isthmus of $G$, the set of pebbles $P:=\{1,\ldots,n\}$ such that $n+k < N$ holds, the initial configuration $f$ and the goal configuration $g$ on the pair $(G,P)$.}
\item[\sl Output]{An explicit sequence of moves to rearrange from the initial configuration $f$ to the goal configuration $g$.}
\end{namelist}
\begin{alg-enumerate}
\item{By using Lemma\tume\ref{feasible-subtree}, find a subtree $T$ of $G$ with $n+k+1$ vertices such that the maximum size of isthmuses in $T$ is at most $k$.}
%\item{Lemma\tume\ref{feasible-subtree}を用いて、位数が$n+k+1$で最大isthmus長が$k$以下となるような$G$の部分木$T$を見つける}
\item{By using the algorithm from Lemma\tume\ref{bunkatsu}, decompose the tree $T$ into subtrees of order $k/8$ or less, and store their data in the binary tree $\BT$.}
%\item{（最適化された）ボードグラフである位数$k+n+1$の木$T$を、centroidで位数$k/8$以下のsubtreeに分解し、それを二分木データベース$\BT$に納める}
\item{By using $\Alg.\PacCon(G,T,\BT, g)$, obtain the configuration $g^{\prime}:=\phi[r]\cdot\psi(g)$ on $T$.}
%\item{$\Alg.\PacCon(G,T,\BT, g)$を用いて、$T$上の配置$g^{\prime}:=\phi[r]\cdot\psi(g)$を得る}
\item{Call $\Alg.\CutOff(G,T,\BT,f,g^{\prime})$.}
%\item{$\Alg.\CutOff(G,T,\BT,f,g^{\prime})$を呼ぶ}
\item{Return the sequence of moves $\psi^{-1}\cdot\phi[r]^{-1}\cdot\sigma$. (\#\# Note that $g=\psi^{-1}\cdot\phi[r]^{-1}\cdot\sigma(f)$.)}
%\item{$\psi^{-1}\cdot\phi[r]^{-1}\cdot\sigma$というmoveの系列を出力して停止する（$g=\psi^{-1}\cdot\phi[r]^{-1}\cdot\sigma(f)$に注意）}
\end{alg-enumerate}
\end{alg}

\section{Board Trees: Proof of Theorem\tume\ref{main}}\label{Main1Proof}

In this section, we present a proof of Theorem\tume\ref{main}.

\noindent
\begin{Proofof}{Theorem\tume\ref{main}}
Since the correctness of $\Alg.\Main((G, P), f, g)$ is straightforward, 
we focus here on analyzing its time complexity, which is shown to be 
$\Ord(N + n\D(G) + n^2\log(\min\{n,k\}))$, 
and on proving that the length of the output solution sequence is bounded by 
$\Ord(n\D(G) + n^2\log(\min\{n,k\}))$. 

It is worth noting that the initial $\Ord(N)$ term in the time complexity accounts 
for the cost of scanning the input tree $G$ and the configurations $f$ and $g$, 
in particular the positions of the pebbles. 
This cost does not contribute to the length of the solution sequence, and thus 
does not appear in the asymptotic bound for the solution length.

From a given input tree $G$, the effort of finding its subtree $T$ of order $n+k+1$ using Lemma\tume\ref{feasible-subtree} is $\Ord(n+k)$.
%与えられた入力tree $G$から、Lemma\tume\ref{feasible-subtree}を使ってその位数$n+k+1$の部分木$T$を見つける手間は$\Ord(n+k)$である。
The task of creating a binary tree database $\BT$ from the obtained $T$ using Lemma\tume\ref{BT} has a cost of $\Ord((n+k)\log(n+k))$.
%得られた$T$から、Lemma\tume\ref{BT}を用いて二分木データベース$\BT$を作る作業には、$\Ord((n+k)\log(n+k))$の手間が生じる。
Using Lemma\tume\ref{sbs0} to convert the arrangements $f$ and $g$ on $(G,P)$ into strongly related arrangements on $(T,P)$ requires a time complexity and number of moves of at most $\Ord(n\D(G))$.
%Lemma\tume\ref{sbs0}を用いて$(G,P)$上の配置$f$や$g$をstrongly relatedな$(T,P)$上の配置に変換する作業には、a time complexity and number of moves of at most $\Ord(Nn)$を要する。


In $\Alg.\Main((G,P), f, g)$, the other parts are spent on repeatedly solving sub-puzzles of size $\Ord(k)$ (Steps \maru{\ref{mini-puzzle}} and \maru{\ref{mini-puzzle-finish}} of $\Alg.\Extraction(T, \BT,h,t)$) on $(T,P)$.
%Mainアルゴリズムにおけるそれ以外のpartは、$(T,P)$上で行われるサイズ$\Ord(k)$のサブパズル（$\Alg.\Extraction(T, \BT,h,t)$のSteps \maru{\ref{mini-puzzle}} and \maru{\ref{mini-puzzle-finish}}）を繰り返し解く操作により費やされる。
According to Lemma\tume\ref{mini-jyunbi}, the repeated process of designating a suitable area of $\Ord(k)$ and gathering $(k+1)$ unoccupied spaces in that location 
(Step\tume\maru{\ref{mini-puzzle}} of $\Alg.\Extraction(T, \BT,h,t)$) for preparing those sub-puzzles has a time complexity and number of moves of at most $\Ord(k(k+n))$ for the entire $\Alg.\Extraction(T, \BT,h,t)$.
%Lemma\tume\ref{mini-jyunbi}によれば、それらのサブパズルを準備するために、適切な$\Ord(k)$の領域を定めてその場所に$k+1$個のunoccupied spacesを寄せてくる作業（$\Alg.\Extraction(T, \BT,h,t)$のStep\tume\maru{\ref{mini-puzzle}}）の繰り返しは、$\Alg.\Extraction(T, \BT,h,t)$全体として、a time complexity and number of moves of at most $\Ord(k(k+n))$を要することが分かっている。
Therefore, in $\Alg.\Main((G,P), f, g)$, the most effort and number of moves required is from the repetition of sub-puzzles of size $\Ord(k)$ (Step\tume\maru{\ref{mini-puzzle-finish}} of $\Alg.\Extraction(T, \BT,g^{\prime},t)$).
%したがって、メインアルゴリズムにおいて、もっとも多くの手間とmoveの数を必要とするのは、サイズ$\Ord(k)$のサブパズル（$\Alg.\Extraction(T, \BT,h,t)$のStep\tume\maru{\ref{mini-puzzle-finish}}）の繰り返しである。
This task is repeated $\Ord((n+k)/k)$ times each time $\Alg.\Extraction(T, \BT,g^{\prime},t)$ is called. If each of these sub-puzzles of size $\Ord(k)$ meets the prerequisites of Theorem\tume\ref{sbs3} (the proof for this part will be postponed until the end), then the total effort and number of moves for one call to $\Alg.\Extraction(T, \BT,h,t)$ is at most $\Ord(((n+k)/k) \cdot (k^2\log k))=\Ord((n+k)k\log k)$. 
%この作業は、$\Alg.\Extraction(T, \BT,h,t)$を呼ぶ度に$\Ord((n+k)/k)$回繰り返されるが、これら一つ一つのサイズ$\Ord(k)$のサブパズルがTheorem\tume\ref{sbs3}の前提条件を満たすのであれば（この部分の証明は最後に回す）、その手間とmoveの総数は、一回の$\Alg.\Extraction(T, \BT,h,t)$について、高々$\Ord(((n+k)/k) \cdot (k^2\log k))=\Ord((n+k)k\log k)$である。
$\Alg.\Extraction(T, \BT,g^{\prime},t)$ is called $\Ord((n+k)/k)$ times within $\Alg.\CutOff(G,T,\BT,f,g^{\prime})$, so the total cost and estimate of moves here is $\Ord(((n+k)/k)(k(k+n)+(n+k)k\log k ))=\Ord((n+k)^2\log k )$.
%$\Alg.\Extraction(T, \BT,h,t)$は、$\Alg.\CutOff(G,T,\BT,f,h)$の中で$\Ord((n+k)/k)$回呼ばれるから、ここでの総コストとmoveの手数の見積は、
%$\Ord(((n+k)/k)(k(k+n)+(n+k)k\log k ))=\Ord((n+k)^2\log k )$となる。

Since the operation of scaling down the input puzzles $f,g$ on $(G,P)$ to their strongly related puzzles on $(T,P)$ requires $\Ord(n\D(G))$ effort and moves, assuming $n > k$, the estimate for the total effort and number of moves in the main algorithm is $\Ord(n\D(G) +(n+k)^2 \log k)=\Ord(n\D(G)+n^2 \log k)$.
%$(G,P)$上のパズルを$(T,P)$上のパズルにスケールダウンする操作に$\Ord(Nn)$の手間がかかるので、$n>k$を仮定すると、メインアルゴリズムの手間とmoveの総数の見積は、$\Ord(Nn +(n+k)^2 \log k)=\Ord(Nn+n^2 \log k)$となる。
On the other hand, if $n \leq k$, then Step\tume\maru{\ref{mini-puzzle-finish}} of $\Alg.\Extraction(T, \BT,g^{\prime},t)$ is repeated only a constant number of times for the entire main algorithm, so the total effort estimate is $\Ord(N+n\D(G) + n^2\log n)$, just as stated in Theorem\tume\ref{sbs3}.
%一方、$n\leq k$であれば、（$\Alg.\Extraction(T, \BT,h,t)$のStep\tume\maru{\ref{mini-puzzle-finish}}）の繰り返しは、メインアルゴリズム全体として定数回しか繰り返されないため、トータルの手間の見積は、Theorem\tume\ref{sbs3}そのままに、$\Ord(Nn + n^2\log n )$となる。

Lastly, we will confirm that repeatedly performing Step\tume\maru{\ref{hige}} of $\Alg.\CutOff(G,T,\BT,f,g^{\prime})$ or Step\tume\maru{\ref{hige2}} of $\Alg.\Extraction(T, \BT,g^{\prime},t)$ does not increase the order of $X(w)$ beyond twice its original value. In this process, the edge $uv$ added to $X(w)$ is a pendant edge, and since $u$ is a leaf, by definition it is not a vertex on any isthmus. Therefore, even if this operation is repeated, no pendant edge will be added to $u$. By adding the edge $uv$, the size of a maximum isthmus has strictly decreased, which means that the degree of $v$ in the current graph is $3$ or more. Consequently, the degree of $v$ in $X(w)$ is either $3$ or a unique vertex $p$ outside $X(w)$ is adjacent to $v$. By definition, a vertex with a degree of $3$ or more is not an interior vertex of any isthmus, so as long as the degree of $v$ in the current graph is $3$, no further pendant edges will be added to $v$ through the repetition of this operation. For the degree of $v$ to again become $2$, the only option is to delete the vertex $p$ outside $X(w)$ and the edge $pv$, but since $u$ is not a vertex on any isthmus, the size of a maximum isthmus in the remaining graph will not increase by this deletion of $pv$. Therefore, the number of pendant edges that can potentially be added to each vertex of $X(w)$ is at most $1$, and since no pendant edges will be added to the newly added leaves, the order of $X(w)$ can increase by at most a factor of $2$. This implies that $k + 1 \geq k=2(2(k/8)+2(k/8)) \geq 2(|V(X(t))| + |V(X(\ell))|)$ holds, and therefore the condition of Theorem\tume\ref{sbs3} (i.e. $|V(X(t))| + |V(X(\ell))| + k + 1 \geq 3(|V(X(t))| + |V(X(\ell))|)$) is always satisfied in Step\tume\maru{\ref{mini-puzzle-finish}} of $\Alg.\Extraction(T, \BT,g^{\prime},t)$. 
%実際、$X(w)$に追加される$uv$はpendant edgeであり、$u$はleafであるから、定義によりisthmus上の頂点ではないので、この操作を繰り返しても、$u$にpendant edgeが追加されることはない。辺$uv$を追加したことにより、最大isthmus長が真に減少したということは、現在のグラフにおいて$v$の次数が$3$以上であることを意味する。したがって、$X(w)$における$v$の次数は$3$であるか、もしくは、$X(w)$の外部に唯一の頂点$p$が存在して$v$と隣接している。定義により、次数$3$以上の頂点はisthmusの内点ではないので、現在のグラフにおける$v$の次数が$3$である限り、$v$に更なるpendant edgeが（この作業の繰り返しにより）追加されることもない。$v$の次数が再び$2$になるためには、$X(w)$の外部の頂点$p$および辺$pv$を削除するしかないが、$u$はisthmus上の頂点ではないので、$pv$の削除によって残されたグラフの最大isthmus長は増加しない。以上のことにより、$X(w)$の各頂点に追加される可能性があるpendant edgeの本数は高々１回であり、追加されたleafにpendant edgeが追加されることもないので、$X(w)$の位数は高々$2$倍以下にしかならない。
%%%%%%%%%%%%%%%%%%%%%%%%%%%%%%%%%%%%%%%%%%%%%%%%%%%%%%%%%%%%%%%%%%%%%%%%%%%%%%%%%%%%%%%%%%%%%%%%%%%%%%%%%%%%%%%%%%%%%%%%%%%%
%%%%%%%%%%%%%%%%%%%%%%%%%%%%%%%%%%%%%%%%%%%%%%%%%%%%%%%%%%%%%%%%%%%%%%%%%%%%%%%%%%%%%%%%%%%%%%%%%%%%%%%%%%%%%%%%%%%%%%%%%%%%
%%%%%%%%%%%%%%%%%%%%%%%%%%%%%%%%%%%%%%%%%%%%%%%%%%%%%%%%%%%%%%%%%%%%%%%%%%%%%%%%%%%%%%%%%%%%%%%%%%%%%%%%%%%%%%%%%%%%%%%%%%%%
%%%%%%%%%%%%%%%%%%%%%%%%%%%%%%%%%%%%%%%%%%%%%%%%%%%%%%%%%%%%%%%%%%%%%%%%%%%%%%%%%%%%%%%%%%%%%%%%%%%%%%%%%%%%%%%%%%%%%%%%%%%%
%%%%%%%%%%%%%%%%%%%%%%%%%%%%%%%%%%%%%%%%%%%%%%%%%%%%%%%%%%%%%%%%%%%%%%%%%%%%%%%%%%%%%%%%%%%%%%%%%%%%%%%%%%%%%%%%%%%%%%%%%%%%
%%%%%%%%%%%%%%%%%%%%%%%%%%%%%%%%%%%%%%%%%%%%%%%%%%%%%%%%%%%%%%%%%%%%%%%%%%%%%%%%%%%%%%%%%%%%%%%%%%%%%%%%%%%%%%%%%%%%%%%%%%%%
%%%%%%%%%%%%%%%%%%%%%%%%%%%%%%%%%%%%%%%%%%%%%%%%%%%%%%%%%%%%%%%%%%%%%%%%%%%%%%%%%%%%%%%%%%%%%%%%%%%%%%%%%%%%%%%%%%%%%%%%%%%%
%%%%%%%%%%%%%%%%%%%%%%%%%%%%%%%%%%%%%%%%%%%%%%%%%%%%%%%%%%%%%%%%%%%%%%%%%%%%%%%%%%%%%%%%%%%%%%%%%%%%%%%%%%%%%%%%%%%%%%%%%%%%
%%%%%%%%%%%%%%%%%%%%%%%%%%%%%%%%%%%%%%%%%%%%%%%%%%%%%%%%%%%%%%%%%%%%%%%%%%%%%%%%%%%%%%%%%%%%%%%%%%%%%%%%%%%%%%%%%%%%%%%%%%%%
%%%%%%%%%%%%%%%%%%%%%%%%%%%%%%%%%%%%%%%%%%%%%%%%%%%%%%%%%%%%%%%%%%%%%%%%%%%%%%%%%%%%%%%%%%%%%%%%%%%%%%%%%%%%%%%%%%%%%%%%%%%%
%まず最初に与えられた大きな木$G$上のパズル（ボードサイズは$N$、碁石のサイズは$n$、最大イスムスのサイズは$k$と仮定）を、碁石の数は$n$のままでボードサイズだけを$\Ord(n+k)$にした、少し小さい部分木$T$の上のfeasibleなパズルにまで、サイズダウンします。その作業に$\Ord(Nn)$の手間が必要です。
%そしてこのスケールダウンしたパズルを$\Ord(k)$の位数の$\Ord((k+n)/k)$個の点部分木による被覆にまで分解して、それを二分木データベースとして表現します。この作業には、$\Ord(k+n)$の手間でcentroidを同定する作業を、高々$\Ord(\log(k+n))$回繰り返す必要があるので、トータルで$\Ord((k+n)\log(k+n))$の計算量がかかります。
%このデータベースを用いて、$T$のあるleafを含む位数$k$の連結成分（$T$を分解した要素の一つ）を、深さ優先探索に付随するトラバースの流れに従って、各成分から吸い出す操作（$T$全体に散らばって置かれている「ターゲット$X(t)$上に配置されるべき碁石の集合」を$X(t)$の近くに集めてきて、その後$X(t)$上に正しく配置する操作）は高々$2((n+k)/k)$回の（小さなサブパズル上での）交換により達成される。一つのターゲット$X(t)$に対して、$\Alg$.吸い出し$(\BT2,f,g^{\prime},t)$において行われる$k$個の空白の準備動きの総数は、Lemma\tume\ref{mini-jyunbi}により、
%$\Ord((n+k)k)$程度です。
%従って、この操作を高々$\Ord((n+k)/k)$回行うと全てのターゲット上（$T$の分解の要素である位数$\Ord(k)$の連結成分が$\Ord((n+k)/k)$個）の配置が終了して、$\Ord((n+k)/k)$個の成分では、それぞれ正しい位置に碁石が置かれます。
%Theorem\tume\ref{sbs3}によれば、中上川先生のアルゴリズムを用いて1回の小さいサブパズルを解く作業の手間は$\Ord(k^2 \log(k))$であり、それが一回の$\Alg$.吸い出し$(\BT2,f,g^{\prime},t)$で$\Ord((n+k)/k)$回行われるので、
%$\Ord((n+k)k \log(k))$の手間になる。Lemma\tume\ref{mini-jyunbi}の$k$個の空白の準備動きの総数$\Ord((n+k)k)$を加えても、$\Ord((n+k)k \log k )$の手間であることは変わらない。
%ターゲット$X(t)$は$\Ord((n+k)/k)$個あるので、$\Alg.$切り落とし$(\BT,f,h)$におけるトータルの計算量は$\Ord((n+k)^2 \log k )$の手間になる。
%初期パズルを$n+k+1$サイズに最適化されたパズルにスケールダウンする操作に$\Ord(Nn)$の手間がかかるので、$n>k$を仮定すると、$\Ord(Nn +(n+k)^2 \log(k))=\Ord(Nn+n^2 \log(k))$となる。
%$n\leq k$であれば、（小さなサブパズルによる）交換は全体として高々16（定数）回しか繰り返されないため、トータルの手間は一回の交換に掛かる手間の計算量のオーダーと同じであり、中上川先生のアルゴリズムの計算量をそのまま借用して、$\Ord(Nn + n^2\log n )$となる。
\end{Proofof}

%%%%%%%%%%%%%%%%%%%%%%%%%%%%%%%%%%%%%%%%%%%%%%%%%%%%%%%%%%%%%%%%%%%%%%%%%%%%%%%%
%%%%%%%%%%%%%%%%%%%%%%%%%%%%%%%%%%%%%%%%%%%%%%%%%%%%%%%%%%%%%%%%%%%%%%%%%%%%%%%%
%%%%%%%%%%%%%%%%%%%%%%%%%%%%%%%%%%%%%%%%%%%%%%%%%%%%%%%%%%%%%%%%%%%%%%%%%%%%%%%%
\section{General Board Graphs: Proof of Theorem\tume\ref{main2}}\label{Main2Proof}

In this section, we provide a proof of Theorem\tume\ref{main2}. 
Now that the efficient algorithm for the \textsc{Pebble Motion Problem on Trees} and 
its computational complexity evaluation has been clarified in the previous sections, 
the actual proof of this theorem can be obtained in a simple manner. 

\begin{Proofof}{Theorem\tume\ref{main2}}
First, computing the $2$-connected component decomposition of the input graph $G$, as well as 
constructing a breadth-first search spanning tree $T_G$ of $G$, requires 
$\Ord(|E(G)|)$ time. 
However, this computational cost is purely algorithmic overhead and does not affect the length of 
the output solution sequence in any way.

In Step\tume\maru{1} of the Main Algorithm ($\Alg.\Main((G, P), f, g)$), when finding a subtree $T$ 
of order $n + k + 1$ of $T_G$, we apply Lemma\tume\ref{feasible-subtree} separately to each part of $T_G$ 
corresponding to the $2$-connected component decomposition of $G$. 
That is, for each $2$-connected component $X$ of $G$, it suffices to ensure that if a part of the subtree $Y$ 
of $T_G$ that spans $X$ is included in $T$, then the size of that part does not exceed the maximum isthmus 
length of $Y$.
It is the same in that $T$ is decomposed into $\Ord(k)$ size subtrees and represented as binary tree data. 
The difference appears in the computational complexity evaluation. 
%我々がここで議論する2-連結ボードグラフGの上のアルゴリズムは、input board treeの代わりにGの適当な
%全域木Tを用い、パラメータkをmaxmum isthmus sizeではなくk:=N-nと定義しなおすこと以外は、木の上の
%アルゴリズムと変わるところはない。Tを細かいO(k)のsubtreeに分解し、それらを二分木データとして表現する
%ことも同じである。違いは計算量評価の方に現れる。
If the input board graph is $2$-connected then the puzzle is feasible if and only if the board graph is 
not a cycle and has at least two unoccupied spaces. Therefore, no matter how the tree $T$ is 
constructed, $T$ may inevitably have to possess an isthmus with at least $k+1$ vertices, 
and as a result, some of the subtrees in the decomposition of $T$ may have long isthmuses of size
$k+1$ or more. It should be noted that each subtree can contain at most one isthmus of size 
at least $k+1$. In such cases, to create a feasible sub-puzzle of size $\Ord(k)$
(c.f. Step\tume\ref{mini-puzzle} of $\Alg.\Extraction(T, \BT,g^{\prime},t)$), 
the puzzle must be constructed around a vertex $v$ with degree $3$ or higher located somewhere 
on a cycle of $G$ that includes the isthmus. To achieve this, the $\Ord(k)$ pebbles that make up 
the sub-puzzle must be brought around $v$. This operation requires $\Ord((k)\times \min\{n,\CL(G)\})$
moves per each sub-puzzle.
In our algorithm, since we have to solve these sub-puzzles $\Ord((\frac{n}{k})^2)$ times, 
the cost of the above operations for the entire algorithm can be estimated as 
$\Ord((\frac{n}{k})^2 \times k \times \min\{n,\CL(G)\})=\Ord(\frac{n^2\min\{n,\CL(G)\}}{k})$.
%cycle以外の$2$-connected board graphにおいては、$2$個以上の空白があればfeasibleであるから、
%$G$のspanning subtree $T$をどのように工夫して取ろうとしても、$T$が長いisthmusを含み、その結果
%$T$の分解におけるsubtreesの一部がそのisthmus上のO(N-n)のsub path達になってしまう可能性は
%避けられない。そうした場合、O(N-n)サイズのfeasibleなsub-puzzleを作るには、そのisthmusを含む
%Gの部分サイクル上のどこかにある次数３以上の頂点vを中心にしてpuzzleを構成するしかない。
%そのためには、当該sub-puzzleを構成するO(N-n)個のpebble達を、vの周囲に持ってくる必要がある。
%この操作には、一つのsub-puzzle当たり、O((N-n)n)のmoveの手間が必要となる。
%我々のアルゴリズムにおいては、これらのsub-puzzle達をO((n/(N-n))^2)回解かなければならないから、
%上記の手数のコストはアルゴリズム全体ではO((n/(N-n))^2 * (N-n)n)=O(n^3/(N-n))と見積もられる。

Note that, for a given instance $(G, P)$, the feasibility of the puzzle is guaranteed as long as 
$k$ satisfies $k \leq N - n - 1$. 
Therefore, in the above discussion, the constant $k$ in the proof may be replaced with an arbitrary 
fixed integer $\theta$ satisfying $k \leq \theta \leq N - n - 1$, without affecting the correctness of the proof.

The statement of our theorem can be directly derived from the above fact and Theorem\tume\ref{main}. 
\end{Proofof}

%%%%%%%%%%%%%%%%%%%%%%%%%%%%%%%%%%%%%%%%%%%%%%%%%%%%%%%%%%%%%%%%%%%%%%%%%%%%%%%%
%%%%%%%%%%%%%%%%%%%%%%%%%%%%%%%%%%%%%%%%%%%%%%%%%%%%%%%%%%%%%%%%%%%%%%%%%%%%%%%%
%%%%%%%%%%%%%%%%%%%%%%%%%%%%%%%%%%%%%%%%%%%%%%%%%%%%%%%%%%%%%%%%%%%%%%%%%%%%%%%%
\section{Diameters, Shortest Cycles, and Refinements}\label{refinement}

In this section, we present a refined version of the algorithm introduced in the preceding sections.
This refinement enables us to demonstrate that, for any instance, our method achieves provable asymptotic improvements—not only in the order
of the length of the resulting move sequence, but also in the overall computational complexity—when compared with existing approaches that
place pebbles sequentially, one at a time, onto predetermined vertices.

\subsection{Refined Analysis for Board Trees: Proof of Theorem\tume\ref{refinement1}}\label{refinement-tree}

We first note that the time complexity stated in Theorem\tume\ref{main} can be improved to 
$\Ord(N+n\D(G) + n^2 \log(\min\{n,k\}))$, and that the length of the solution sequence produced by 
the algorithm admits a refined upper bound of $\Ord(n\D(G) + n^2 \log(\min\{n,k\}))$.
In fact, this has already been explicitly demonstrated in the proof of Theorem\tume\ref{main}.

% 我々が本論文のprevious sectionまでに提案したアルゴリズムでは、その中のAlgorithm 4 $(\Alg.\Extraction(T,BT,g',t))$の\maru{4}において、
% 現在のboard treeから除去してもreminder graphが連結のままであるようなsub tree $X(\ell)$を指定し、
% この$X(\ell)$を含み頂点数が$|V(X(t))|+|V(X(\ell))|+k+1$である$H(\ell)$というsubtree (ここで$X(t)$は配置を確定させるターゲットのsubtree)を探して、
% $H(\ell)$というboard tree上のsub puzzleを解き、$X(\ell)$上にゴミ（$X(t)$に乗せるべきpebble以外の碁石達）を詰めて切り落とす手続き（\maru{5}）を繰り返します。
In the algorithm proposed up to the previous section of this paper, Step \maru{4} of Algorithm 4 ($\Alg.\Extraction(T, \BT,g^{\prime},t)$) performs the following procedure iteratively:
we identify a subtree $X(\ell)$ of the current board tree such that removing $X(\ell)$ does not disconnect the remainder graph. We then search for a subtree $H(\ell)$ that
contains $X(\ell)$ and has exactly $|V(X(t))| + |V(X(\ell))| + k + 1$ vertices, where $X(t)$ denotes the target subtree on which the placement of pebbles is to be correctly finalized.
Next, we solve the sub-puzzle on $H(\ell)$, pack irrelevant pebbles (i.e., pebbles not intended to be placed on $X(t)$) into $X(\ell)$, and cut off $X(\ell)$ from the board tree (\maru{5}).

% このときまず、
% １） $X(\ell)$の付け根$\JoiV(X(\ell))$と、$X(t)$の付け根$\JoiV(X(t))$とを結ぶ（木ですから『唯一の』）パス $P(\ell,t)$を考えると、$X(\ell)$のどの頂点から$X(t)$のどの頂点に至る最短路も$P(\ell,t)$を必ず使います。
% なので、$P(\ell,t)$のJointVeretx$(X(\ell))$をその一方の端点とする、$P(\ell,t)$の十分長い（少なくとも$2|V(X(t)|+1$以上の頂点数を持つ）サブパス$P'(H(\ell))$が$H(\ell)$に含まれるように、$H(\ell)$を取ることにします。
% これは、$H(\ell)$の頂点数が$|V(X(t))|+|V(X(\ell))|+k+1$ということで、十分余裕がありますので、$X(\ell)$が$X(t)$に極めて近い特別なケース以外では、必ずそのようなサブパス$P'(H(\ell))$を含むように$H(\ell)$を取ることが出来ます。
% 逆に、$P(\ell,t)$の頂点数が$2|V(X(t)|$以下であれば、$H(\ell)$として、$X(t)$を含むように取ってくることが出来ますので、その場合にはそうします。
% ２）上記１）の条件を満たすように取った$H(\ell)$のパス$P'(H(\ell))$の上に、もし、$X(t)$に置かれるべきpebbleが何個か既に乗っていた場合には、それらのpebble達達が乗っている頂点は触らずに、今$X(\ell)$上にある「$X(t)$に置かれるべきpebble達」を、$P'(H(\ell))$の「それ以外の」「できる限り$X(\ell)$から離れた位置にある」頂点のどこかに乗っている（ゴミ）pebble達と入れ替え（位置の交換）をします。
% この処理が終了した段階で、必然的に、$X(\ell)$には「ゴミのみ」が詰められることになります。
% なお、$P(\ell,t)$の頂点数が$2|V(X(t)|$以下であれば、$H(\ell)$は$X(t)$を含みますから、$X(t)$上の正しい位置の頂点に$H(\ell)$上のすべてのpebble達を正しく配置してしまえばいいです。
In the refinement of the algorithm proposed in this section, we proceed as follows:
\begin{enumerate}
      \item Consider the unique path $P(\ell, t)$ in the board tree connecting the joint vertex of $X(\ell)$, denoted by $\JoiV(X(\ell))$, and the joint vertex of $X(t)$, denoted by $\JoiV(X(t))$.
Since the underlying structure is a tree, this path is uniquely determined. Any shortest path from a vertex in $X(\ell)$ to a vertex in $X(t)$ must pass through $P(\ell, t)$.
Based on this observation, we construct $H(\ell)$ so that it contains a sufficiently long subpath $P'(H(\ell)) \subseteq P(\ell, t)$, where $P'(H(\ell))$ starts at $\JoiV(X(\ell))$
and contains at least $2|V(X(t))| + 1$ vertices. Since $H(\ell)$ is allowed to have $|V(X(t))| + |V(X(\ell))| + k + 1$ vertices, it generally has enough capacity to include such a subpath
unless $X(\ell)$ happens to be located very close to $X(t)$. In most cases, therefore, we can ensure that $H(\ell)$ includes $P'(H(\ell))$ as required. Conversely, if the total number
of vertices in $P(\ell, t)$ is at most $2|V(X(t))|$, we instead construct $H(\ell)$ so that it includes $X(t)$ itself.
      \item Suppose that some of the pebbles that are supposed to be placed on $X(t)$ are already located on vertices in $P'(H(\ell))$. In that case, we avoid moving those pebbles.
Instead, for the pebbles currently located in $X(\ell)$ that are eventually intended to be placed on $X(t)$, we swap their positions with "garbage" pebbles (i.e., pebbles not
destined for $X(t)$) located on other vertices of $P'(H(\ell))$, choosing such vertices to be as far from $X(\ell)$ as possible. After this process, all pebbles remaining in $X(\ell)$
are guaranteed to be garbage. Note that if the number of vertices in $P(\ell, t)$ is at most $2|V(X(t))|$, then $H(\ell)$ includes $X(t)$ by construction. In this case, we can directly
solve the sub-puzzle on $H(\ell)$ by placing all pebbles correctly on their target vertices within $X(t)$.
\end{enumerate}

% ２）の手続きは、「本来$X(t)$に置かれるべき$X(\ell)$上のpebble達」の個数$k_{\ell}$が、$\log k$未満である場合には、$\Ord(k^2)$の手間を$k_{\ell}$回繰り返して、トータルで$\Ord(k_{\ell} k^2)$の計算量で達成可能です。
% （ここに、我々の論文のSection 3で述べている「十分な量（pebbleの数の２倍以上）の空白の個数がある場合」の$N$-vertex board treeのケースで、任意の２つのpebbleの位置『のみを』交換する手間が、$\Ord(Nn + n^2)$であること、を使います。
% 実際、交換すべき２つのpebbleを、髭一本の場合であれば髭の両脇あたりに、髭二本のケースでは２つの髭の上にそれぞれを一つずつ、配置した後、髭を使ってそれら２つのpebbleの位置の交換をし、その後、他のpebble達について行った前処理の逆操作で状態を戻せばいいです。
% Section 3の設定では、スペースが十分確保されているおかげで、目的の２つのpebbleの交換作業中は、他のすべてのpebble達は適当な「仮置き場」に避難させておく（前処理）ことができるので、この避難を元に戻す作業を含めて、この作業は全体として$\Ord(Nn+n^2)$で達成可能です。（←この事実の証明については、Section 3のどこかに、中上川先生の方で記載して頂くことは可能でしょうか？）
% もちろん、$k_{\ell}$が、$\log k$以上である場合には、我々がすでにSection 3のTheorem8で示しているアルゴリズムを使って$\Ord(k^2 \log k)$の手間で処理が完了します。

% この１）、２）の付随条件を課しますと、$X(t)$上に置かれるべきpebble達は、初期状態で$G$のどの頂点上に置かれていたとしても、$\Ord(k)$のサイズのboard tree上のサブパズルを1回解くごとに、スタート時の頂点と自身が置かれるべき$X(t)$上のゴールの頂点とを結ぶ$G$上の最短路に沿って、$\Theta(k)$ずつ着実に歩を進める、ことになります。
% なので、この方法であれば、moveの手数は（$G$の直径を$D(G)$とすると）、$\Ord((N-n)n+\min\{k n D(G), n^2 \log(\min\{ n, k \})\})$ で達成可能であることが分かります。
The procedure described in step 2 of the above refinement can be executed with a total computational cost of $\Ord(k_{\ell} k^2)$
when the number of pebbles on $X(\ell)$ that are supposed to be placed on $X(t)$, denoted $k_{\ell}$, is less than $\log k$.
In this case, we perform a swap operation of two pebbles at most $k_{\ell}$ times, each of which costs $\Ord(k^2)$.
To achieve this, we make use of the result from Lemma\tume\ref{sbs4} in Section\tume\ref{base-cases} of this paper:
in a board tree with $N$ vertices and a sufficient number of unoccupied vertices (at least twice the number of pebbles),
swapping the positions of only two pebbles can be done in $\Ord(Nn)$ time. 
Of course, when $k_{\ell} \geq \log k$, we instead apply the algorithm already presented in Theorem\tume\ref{sbs3} of 
Section\tume\ref{base-cases}, which achieves the operation in $\Ord(k^2 \log k)$ time.

By incorporating the refinements 1) and 2) described above, the pebbles destined for $X(t)$ will, regardless of
their initial positions in the board graph $G$, make consistent progress along the shortest path in $G$ from their initial positions
to their target vertices in $X(t)$, by a distance of $\Theta(k)$ each time we solve a sub-puzzle on $H(\ell)$ such that $X(\ell)$
contains such a pebble. 

%%%%%%%%%%%%%%%%%%%%%%%%%%%%%%%%%%%%%%%%%%%%%%%%%%%%%%%%%%%%%%%%%%%%%%%%%%%%%%%%
%%%%%%%%%%%%%%%%%%%%%%%%%%%%%%%%%%%%%%%%%%%%%%%%%%%%%%%%%%%%%%%%%%%%%%%%%%%%%%%%
%%%%%%%%%%%%%%%%%%%%%%%%%%%%%%%%%%%%%%%%%%%%%%%%%%%%%%%%%%%%%%%%%%%%%%%%%%%%%%%%
At the beginning of this section, we stated that the refinement leads to an algorithm whose output exhibits an asymptotically shorter solution
--- compared to existing methods that sequentially reposition pebbles, one by one, onto predetermined locations.
In what follows, we examine this fact in greater detail.
%%%%%%%%%%%%%%%%%%%%%%%%%%%%%%%%%%%%%%%%%%%%%%%%%%%%%%%%%%%%%%%%%%%%%%%%%%%%%%%% 
%%%%%%%%%%%%%%%%%%%%%%%%%%%%%%%%%%%%%%%%%%%%%%%%%%%%%%%%%%%%%%%%%%%%%%%%%%%%%%%%
%%%%%%%%%%%%%%%%%%%%%%%%%%%%%%%%%%%%%%%%%%%%%%%%%%%%%%%%%%%%%%%%%%%%%%%%%%%%%%%%
Let $P_{X(t)}$ denote the set of pebbles that are to be placed on the target subtree $X(t)$.
For each pebble $p$ in $P_{X(t)}$, let $\SP(p)$ denote the unique path in the board tree $G$ from the vertex at which $p$ initially resides
to the vertex in $X(t)$ where $p$ is to be placed.
We associate a formal variable $x_v$ with each vertex $v$ of $G$, and define the polynomial in ${\mathbf x}:=(x_v)_{v\in V(G)}$ 
corresponding to the path $\SP(p)$ as
$$
F_{\SP(p)}({\mathbf x}) := \sum_{v \in V(\SP(p))} x_v,
$$
and 
$$
F({\mathbf x}) := \sum_{p \in P_{X(t)}} F_{\SP(p)}({\mathbf x}).
$$
By expanding $F({\mathbf x})$, let $a_v$ denote the coefficient of the variable $x_v$ corresponding to each vertex $v$ of $G$; that is,
$$
F({\mathbf x}) = \sum_{v \in V(G)} a_v x_v.
$$
Note that each $a_v$ is a non-negative integer. 
Then, we consider the following quantity:
$$
\#\SP(X(t)) := \sum_{v \in V(G)} \min\{a_v, \log k\}.
$$
The total cost of Step\tume\maru{5} in the refinement of Algorithm 4 ($\Alg.\Extraction(T, \BT, g', t)$) 
can be regarded as performing $\Ord(k^2)$ units of computational effort (and moves) up to $\Ord\big(\frac{\#\SP(X(t))}{k}\big)$ times.
Hence, we have:
\begin{prop}\label{refinement1-path-weight}
The above refinement of  Algorithm 4 $(\Alg.\Extraction(T, \BT,g^{\prime},t))$  can be stopped with a time complexity in $\Ord(k \cdot \#\SP(X(t)))$.
\end{prop}
Note that in existing approaches in the literature, where pebbles are placed one by one on the board tree,
the above complexity can grow to the order of $\Omega(k\sum_{p\in P_{X(t)}} |\SP(p)|)=\Omega(k \sum_{v \in V(G)} a_v)$.

%%%%%%%%%%%%%%%%%%%%%%%%%%%%%%%%%%%%%%%%%%%%%%%%%%%%%%%%%%%%%%%%%%%%%%%%%%%%%%%% 
%%%%%%%%%%%%%%%%%%%%%%%%%%%%%%%%%%%%%%%%%%%%%%%%%%%%%%%%%%%%%%%%%%%%%%%%%%%%%%%%
%%%%%%%%%%%%%%%%%%%%%%%%%%%%%%%%%%%%%%%%%%%%%%%%%%%%%%%%%%%%%%%%%%%%%%%%%%%%%%%%
%%%%%%%%%%%%%%%%%%%%%%%%%%%%%%%%%%%%%%%%%%%%%%%%%%%%%%%%%%%%%%%%%%%%%%%%%%%%%%%%
%%%%%%%%%%%%%%%%%%%%%%%%%%%%%%%%%%%%%%%%%%%%%%%%%%%%%%%%%%%%%%%%%%%%%%%%%%%%%%%%
%%%%%%%%%%%%%%%%%%%%%%%%%%%%%%%%%%%%%%%%%%%%%%%%%%%%%%%%%%%%%%%%%%%%%%%%%%%%%%%%
For each pebble $p \in P_{X(t)}$, it is clear that $|\SP(p)| \leq \D(T)+1 \leq \D(G)+1$, where 
$\D(T)$ and $\D(G)$ denote the diameters of the board graphs $T$ and $G$, respectively.
Therefore, we have $\#\SP(X(t)) \leq \sum_{p\in P_{X(t)}} |\SP(p)| \in \Ord(k \D(T))$.
Within the main algorithm, Algorithm 4 ($\Alg.\Extraction(T, \BT, g', t)$) is called at most $\Ord(n / k)$ times.
Therefore, the overall computational complexity contributed by this refinement is $\Ord(n k \D(T))$, which is 
at most $\Ord(n k \D(G))$.
Together with the result of Theorem\tume\ref{main}, this completes the proof of Theorem\tume\ref{refinement1}. \qed

We conclude this subsection with a brief remark on Theorem\tume\ref{refinement1}.
In our main algorithm $\Alg.\Main((G,P), f, g)$, we first reduce the pebble motion puzzle on 
the instance's board tree $G$ to a smaller puzzle on a subtree $T$ of $G$, where $T$ is 
a board graph with $k+n+1$ vertices. 
Then, we solve the puzzle on this reduced subtree $T$ and restore the resulting solution to 
the original board tree $G$ with $\Ord(N+n\D(G))$ effort and $\Ord(n\D(G))$ moves. 
Therefore, more precisely, the occurrence of $\D(G)$ inside the $\min\{\}$ expression 
in the second term of the asymptotic analysis in Theorem\tume\ref{refinement1} should 
instead be written as $\D(T)$.
However, since the instance graph is $G$ and the tree $T$ is merely an internal variable used 
within the algorithm, we have expressed the bound using $\D(G)$ rather than $\D(T)$.

As noted in the previous section, here as well, all occurrences of the parameter $k$ can be 
simultaneously replaced by any fixed integer $\theta$ satisfying $k \leq \theta \leq N - n - 1$, 
without invalidating the statement.


\subsection{Refined Analysis for General Board Graphs}\label{refinement-general}


As in Section~\ref{Main2Proof}, we first select an arbitrary vertex in the
connected board graph $G$ and construct a spanning tree $T_G$ by performing
a breadth-first search.
There is a natural relationship between the diameter $\D(G)$ of $G$
and the diameter $\D(T_G)$ of $T_G$, namely
$\D(G) \le \D(T_G) \le 2\D(G)$,
which implies $\Theta(\D(T_G)) = \Theta(\D(G))$.
We then apply the algorithm from the above 
subsection\tume\ref{refinement-tree} to this BFS spanning tree $T_G$. 
%% board graph $G$が２連結成分を含む場合にも、任意の頂点から１回幅優先探索を掛けて幅優先探索木を作れば、$G$の直径の高々２倍の直径を有する全域木$T_G$が得られます。
%% $G$と$T_G$の直径のオーダーは変わりませんので、これからは、$T_G$上での直径$D(T_G)$を$D(G)$とみなして、議論を進めます。
%% 先に示した「board graph $G$が木である場合のアルゴリズム」を$T_G$に適用することを考えます。

As already discussed in the previous section, the BFS spanning tree $T_G$ of $G$ may contain 
an isthmus of size at least $k+1$ in the subtrees that span the $2$-connected components of $G$.
Thus, continuing the same line of reasoning as in the previous section, we again encounter a potential issue:
the candidates for $H(\ell)$ may contain an isthmus of size at least $N-n$.
In this case, we fix such a candidate and denote it by $H(\ell)$.
This $H(\ell)$ contains exactly one isthmus $I(\ell)$ of size at least $N - n$, which lies along the subpath 
of the path $P(\ell, t)$ connecting the root $\JoiV(X(\ell))$ of $X(\ell)$ and the root $\JoiV(X(t))$ of $X(t)$.
If it is possible to augment $H(\ell)$ by adding appropriate edges from $G \setminus T_G$ so as to reduce the 
size of the longest isthmus and obtain a feasible board tree, then we simply add such edges to $H(\ell)$, 
render it feasible, and solve the resulting sub-puzzle.
If even this is not possible, then the graph $G$ must contain a chordless cycle $C$ of length at most $\CL(G)$
that passes through the isthmus $I(\ell)$.

In this case, we temporarily consider the structure obtained by adding $C$ to the tree $T_G$.
By simultaneously shifting the pebbles along the cycle $C$, we can reposition the pebbles on $I(\ell)$ 
around the vertices of degree at least three in $T \cup C$, thereby enabling the construction of a feasible 
sub-puzzle $H'(\ell)$.
The process of constructing $H'(\ell)$ by collectively shifting the pebbles along $C$ requires at most 
$\Ord(k \cdot \min\{\CL(G), n\})$ moves per instance of $X(\ell)$.
This reasoning is entirely analogous to the argument used in the proof of Theorem\tume\ref{main2} 
in Section\tume\ref{Main2Proof}. 
% このとき、唯一の問題点は、$H(\ell)$を作る際に、$N-n$以上のサイズのisthmusを持つものしか得られない可能性があることです。
% このとき$T_G$は、 $X(\ell)$の付け根$\JoiV(X(\ell))$と$X(t)$の付け根$\JoiV(X(t))$とを結ぶパス$P(\ell,t)$の部分パスとして、サイズ$N-n$以上のisthmus $I(\ell)$を持ちます。
% $H(\ell)$に（$T_G$以外の）$G$の適切な辺を付け加えてisthmusを短くできるのであれば、そのような（$T_G$の外の）辺を$H(\ell)$に追加して、$H(\ell)$をfeasibleにしてから、そのサブパズルを解けばよいです。
% それすらできない場合には、$G$には、このisthmus $I(\ell)$を含む長さ$\CL(G)$以下のコードレスサイクル$C$が存在しますので、$C$上でpebble達を一斉に動かすことにより、$I(\ell)$上のpebble達を$G$の次数３以上の頂点の周りに持ってくると、feasibleなサブパズル$H^{\prime}(\ell)$を作れます。
% この$C$上でpebble達を一斉に動かして$H^{\prime}(\ell)$を構成する作業は、１つの$X(\ell)$について、高々$\Ord((N-n)\min\{\CL(G),n\})$回のmoveを必要とします。
% これは、Theorem 4の証明における議論と全く同じ理屈です。

The total number of moves involved in the sequence of operations can be calculated by focusing on 
one of the $\Ord(k)$ pebbles, denoted by $p$, out of the $|X(t)|$ pebbles. Specifically:
\begin{itemize}
\item When pebble $p$ is placed on $X(\ell)$, a series of preparation moves, as described earlier, are required to construct $H^{\prime}(\ell)$. 
Afterward, solving the subproblem on $H^{\prime}(\ell)$ involves traversing a shortest path on the graph $G$ from the initial position of $p$ 
(the vertex $f^{-1}(p)$) to its goal position on $X(t)$ (the vertex ${g^{\prime}}^{-1}(p)$). 
Each step in this path moves $\Ord(k)$ units towards the goal position ${g^{\prime}}^{-1}(p)$, and this process is repeated.
\item For each pebble, the total number of subproblems to be solved is at most $\Ord(\D(G)/k)$, leading to the goal vertex.
\end{itemize}
Therefore, when aggregating the total effort for all $\Ord(k)$ pebbles on $X(t)$, the overall number of moves required 
in this process can be expressed as:
$$
\Ord(k \min\{\CL(G),n\}) \cdot \Ord(\D(G)/k) \cdot \Ord(k) = \Ord(k\min\{\CL(G),n\}\D(G)).
$$
This quantity should be regarded as an overhead added to the cost
obtained under the assumption that $G$ is a tree.

Since there are $\Ord(n/k)$ candidates for $X(t)$, 
the total overhead of the algorithm is therefore
$$\Ord(k\min\{\CL(G),n\}\D(G)) \cdot \Ord(n/k) = \Ord(n\min\{\CL(G),n\}\D(G)).$$

% 一連の作業により掛かるmoveは、$|X(t)|$個（$\Ord(N-n)$個）のpebbleの一つ $p$ に着目して計算すると、
% ・$p$が$X(\ell)$上に置かれている場合に、$H^{\prime}(\ell)$を構成するために上記の（準備の）手間のmoveが発生し、その後$H^{\prime}(\ell)$上のサブパズルを解くことで、$p$のinitial positionとしての$G$の頂点$f^{-1}(p)$と、$p$のgoal positionとなる$X(t)$上の頂点${g^{\prime}}^{-1}(p)$を結ぶ$G$上の（長さ$\Ord(D(G))$の）最短パスの上をgoal position ${g^{\prime}}^{-1}(p)$に向かって$\Ord(N-n)$ずつ歩を進める、という一連の作業が繰り返される
% ので、
% ・pebble１個あたり、トータルとして高々$\Ord(D(G)/(N-n))$回のサブパズルを解いてゴールとなる頂点上に行き着く
% 訳です。
% よって$X(t)$上のすべてのpebble（は$\Ord(N-n)$個ある）の手間を合算すると、この一連の過程で、
% $\Ord((N-n)\min\{\CL(G),n\}(D(G)/(N-n))(N-n))=\Ord((N-n)\min\{\CL(G),n\}D(G))$のmoveの手間が（board graph Gが木の場合の見積に加えて）余分に発生する、と考えることが出来ます。
% さらに、ターゲットである$X(t)$の候補は、全部で$\Ord(n/(N-n))$個あるので、この一連の過程は高々$\Ord(n/(N-n))$回繰り返されます。
% したがってこの余分な作業は、アルゴリズム全体として、$\Ord((N-n)\min\{\CL(G),n\}D(G)(n/(N-n)))=\Ord(n\min\{\CL(G),n\}D(G))$のmoveであると見積もることが出来ます。

In the previous subsection (Subsection\tume\ref{refinement-tree}), we established 
that the length of the solution sequence is 
$$
\Ord(n\D(G) + \min\{kn\D(G), n^2 \log(1+\min\{n, k\})\})
$$
when the board graph $G$ is a tree.
By adding the order of the computational cost for the preparatory moves estimated above, 
we obtain the following order for the move operations:
$$
\Ord(n\D(G) + n \min\{\CL(G), n\} \D(G) + \min\{kn \D(G), n^2 \log(1+\min\{n, k\})\}).
$$
Although similar remarks have been repeatedly made before, all occurrences of the parameter $k$ 
above can also be simultaneously replaced by any fixed integer $\theta$ satisfying $k \leq \theta \leq N - n - 1$, 
without invalidating the statement.
By comparing this result with the one derived in Theorem\tume\ref{main2} and taking the minimum of 
both orders, we obtain the overall time complexity as follows:
\begin{thm}\label{refinement2}
Assume that a shortest cycle passing through each vertex $v$ of $G$ (or $\emptyset$ if no such cycle exists) 
is provided in advance.
Let $\theta$ be an arbitrary integer such that $k \leq \theta \leq N-n-1$.
Then, for the \textsc{Pebble Motion Problem on General Graphs}, a solution sequence of length
$\Ord(n\D(G) + n \min\{\CL(G), n\} \min\{\D(G), \frac{n}{\theta}\} + \min\{n\theta\D(G), n^2 \log(1+\min\{n, \theta\})\})$
can be computed in time 
$\Ord(|E(G)| + n\D(G)+n \min\{\CL(G), n\} \min\{\D(G), \frac{n}{\theta}\} + \min\{n\theta\D(G), n^2 \log(1+\min\{n, \theta\})\}).$
\end{thm}

Theorem\tume\ref{refinement2-general-constant-blanc}, stated earlier, follows 
immediately as a corollary of Theorem~\ref{refinement2}.
\begin{Proofof}{Theorem\tume\ref{refinement2-general-constant-blanc}}
Note that $k \leq \theta < N-n \in \Ord(1)$ and hence $\{\CL(G), \D(G)\} \subset \Ord(n)$ holds in this setting.  
As a consequence, the time complexity characterization in Theorem\tume\ref{refinement2} can be simplified as 
$\Ord(|E(G)| + n\D(G)+ n \CL(G) \D(G) + n\D(G))=\Ord(|E(G)| + n \CL(G) \D(G)). $
Similarly, the solution sequence length also simplifies to $\Ord(n \CL(G) \D(G))$.
\end{Proofof}

In the proof of Theorem~\ref{refinement2}, we estimated the overhead required to eliminate the infeasibility 
caused by an excessively long isthmus in $T_G$.
More precisely, when such an isthmus appears, we consider the shortest cycle in $G$ containing this isthmus 
and restore feasibility by simultaneously shifting all pebbles along the cycle by a distance equal to the length 
of the isthmus.

This overhead also admits another upper bound.
For the reader's convenience, we recall the definition of $V(G, q)$, which denotes the set of degree-two vertices in $G$ 
that lie on paths in $G$ and are part of a contiguous sequence of at least $q$ degree-two vertices.
Vertices in $V(G, N-n-2)$ may create overly long isthmuses in the spanning tree $T_G$, serving as sources of infeasibility for the PMT.  
In the worst case, every vertex in this set induces such an isthmus in $T_G$. 
Consequently, the overhead required to move pebbles to their designated positions on a single target subtree $X(t)$ is bounded by
$\Ord\!\left(\sum_{v \in V(G, N-n-2)} \CL(G, v)\right)$.
Since the number of target subtrees is generally $\Ord(n/k)$,
setting $k = N-n$ gives that the number of target subtrees is $\Ord(n/(N-n))$.
Multiplying this bound by the number of target subtrees then gives the following estimate for the total overhead of the algorithm:
$\Ord\!\left(\frac{n}{N-n} \sum_{v\in V(G,N-n-2)} \CL(G,v)\right)$.
From this observation, we obtain the following theorem as a corollary.
\begin{thm}\label{refinement3}
For every feasible instance of the \textsc{Pebble Motion Problem on General Graphs}, 
the length of an optimal solution sequence is in 
$$\Ord\Bigl(n\D(G) + \frac{n}{N-n}\sum_{v\in V(G,N-n-2)} \CL(G,v) + n^2 \log(1+\min\{n, N-n\})\Bigr).$$
\end{thm}

Now, we restate Theorem~\ref{Restriction} for convenience.
{\setcounter{thm}{5}
\begin{thm}\label{Restriction}
Let $(G,P)$ be a feasible instance of the \textsc{Pebble Motion Problem on General Graphs}.
If the length of an optimal solution sequence is $\Theta(N^3)$, then the board graph $G$ must satisfy 
the following conditions:
\begin{enumerate}
\item $N-n = \Ord(1)$;\label{CLL1}
\item $\displaystyle \sum_{v\in V(G,N-n-2)} \CL(G,v) = \Theta(N^2)$;\label{CLL}
\item The graph $G$ contains a set $S \subseteq V(G,N-n-2)$ of $\Theta(N)$ vertices;\label{CLL3}
\item For every vertex $v \in S$, the length of the shortest cycle in $G$ passing through $v$ is $\Theta(N)$.\label{CLL4}
\end{enumerate}
\end{thm}}
\begin{Proofof}{Theorem~\ref{Restriction}}
Conditions \ref{CLL1} and \ref{CLL} follow immediately from Theorem~\ref{refinement3}.
Since both $\CL(G)$ and $|V(G,N-n-2)|$ are in $\Ord(N)$, Condition \ref{CLL} clearly implies Conditions \ref{CLL3} and \ref{CLL4}. 
\end{Proofof}

\noindent
This observation shows that the $\Theta(N^3)$ behavior of the PMP
can arise only under highly restrictive structural conditions on the instance $(G,P)$.
%%%%%%%%%%%%%%%%%%%%%%%%%%%%%%%%%%%%%%%%%%%%%%%%%%%%%%%%%%%%%%%%%%%%%%%%%%%%%%%%
%%%%%%%%%%%%%%%%%%%%%%%%%%%%%%%%%%%%%%%%%%%%%%%%%%%%%%%%%%%%%%%%%%%%%%%%%%%%%%%%
%%%%%%%%%%%%%%%%%%%%%%%%%%%%%%%%%%%%%%%%%%%%%%%%%%%%%%%%%%%%%%%%%%%%%%%%%%%%%%%%%
%
%\textcolor{red}{
%\section{A Structural Characterization of Cubic Complexity}\label{ThetaNcube} 
%}
%
%%%%%%%%%%%%%%%%%%%%%%%%%%%%%%%%%%%%%%%%%%%%%%%%%%%%%%%%%%%%%%%%%%%%%%%%%%%%%%%%
%%%%%%%%%%%%%%%%%%%%%%%%%%%%%%%%%%%%%%%%%%%%%%%%%%%%%%%%%%%%%%%%%%%%%%%%%%%%%%%%
%%%%%%%%%%%%%%%%%%%%%%%%%%%%%%%%%%%%%%%%%%%%%%%%%%%%%%%%%%%%%%%%%%%%%%%%%%%%%%%%



\section{Concluding Remarks}\label{final}

In the concluding section ``Conclusion and Open Problems'' of \cite{715921}, 
the authors remark that:
\begin{quote}
It would be useful to at least have an efficient algorithm which approximates 
the number of moves required. For it seems that only a small fraction of the 
graph puzzles actually require $\Ord(N^3)$ moves.
\end{quote}
To illustrate this phenomenon, the authors note that when the board 
graph is a square grid with $N$ vertices, the generalized $15$-puzzle can be 
solved in $\Ord(N^{3/2})$ moves.

Our Theorems \ref{refinement1}, \ref{refinement2-general-constant-blanc}, 
\ref{refinement2}, and \ref{refinement3}, established in Section~\ref{refinement}, 
provide a partial answer to this question by extending the square-grid example 
to general board graphs and by presenting algorithms that achieve efficient 
asymptotic bounds both on the length of the output move sequence and on the 
computation time.
Indeed, in the above square-grid example we have $\CL(G)=4=\Ord(1)$ and
$\D(G)\in\Ord(\sqrt{N})$. 
If the number of unoccupied vertices satisfies $N-n\in\Ord(1)$, then a solution 
sequence of length $\Ord(n\CL(G)\D(G))=\Ord(N^{3/2})$ can be computed in time 
linear in its length.

Our results can be summarized as follows.
\begin{enumerate}
\item For every instance of the \textsc{Pebble Motion Problem on Trees},  
the length of the shortest solution sequence is bounded by $\tilde{\Ord}(N^2)$.  
This result not only rectifies an invalid proof given in \cite{715921}, 
but also constitutes a genuine improvement over the state-of-the-art 
bound presented in the most recent work \cite{WOS:001177663400002}.

\item For the case of the \textsc{Pebble Motion Problem on General Graphs},
we show that the extremal behavior $\Theta(N^3)$ can occur only under highly 
restricted structural conditions.
More precisely, such behavior arises only under the following highly restricted conditions:
\begin{enumerate}
\item the number of unoccupied vertices $N-n$ is constant;
\item the graph $G$ contains a set $S$ of $\Theta(N)$ degree-two vertices, 
each of which lies in a sequence of at least $N-n-2$ consecutive degree-two 
vertices on a path of $G$;
\item for every vertex $v \in S$, the length of the shortest cycle in $G$
that passes through $v$ is $\Theta(N)$.
\end{enumerate}
This description precisely captures the nature of the examples presented 
in \cite{715921}, in which the length of the shortest solution sequence 
reaches $\Theta(N^3)$.
\end{enumerate}

More broadly, these results offer a clearer structural understanding of 
the \textsc{Pebble Motion Problem} by demonstrating that the previously known 
$\Ord(N^3)$ worst-case bound can only arise under highly restricted 
and specific structural configurations of the underlying graph.

Next, we discuss the relationship between our results on the 
\textsc{Pebble Motion Problem} (PMP) for general graphs and the 
PMP-based techniques proposed in the context of 
Multi-Agent Path Finding (MAPF).
Our results suggest a different perspective on algorithmic approaches to 
PMP in MAPF. Rather than treating special graph classes separately, 
our analysis indicates that the tree case can serve as a natural base 
structure for more general instances. In particular, computations can 
be organized around a spanning breadth-first search tree of the graph, 
while the remaining structural features of the graph are addressed only 
when needed to resolve infeasible configurations.
This viewpoint leads to a more unified treatment of PMP across different 
graph structures. In particular, our improved understanding of the 
complexity of the problem on trees—showing that it admits an 
$\Ord(N^2 \log N)$-time algorithm—suggests that tree-based formulations 
may serve as a computationally viable foundational framework. 
We believe that this perspective may provide useful guidance for the design 
of PMP solvers within MAPF systems.

Finally, we provide some intuition behind Conjecture~\ref{conjecture}.
It is already known that the shortest solution sequences for PMT
may require $\Omega(N^2)$ moves.
Moreover, there exist instances in which moving a single pebble
by a distance $\Theta(N)$ along the board tree already requires
$\Theta(N^2)$ moves.
In view of this phenomenon, establishing an upper bound of
$\tilde{\Ord}(N^2)$ on the length of the shortest solution sequence
appears difficult without relying on some form of divide-and-conquer
strategy.

To illustrate the underlying intuition, we consider the following
model case (which is also discussed in Remark~\ref{counterexample} 
in the Appendix).
Let $P$ be a sufficiently long path graph of even length, and attach
a pendant edge $uv$ to the midpoint vertex $u$ of $P$.
The resulting board graph is $G := P + uv$.
Let $P_L$ and $P_R$ denote the left and right subpaths of $P$
obtained by removing the vertex $u$.
Let $\ell := |V(P_L)| = |V(P_R)|$.
In this board tree, the maximum isthmus size equals $\ell$, and hence
at most $\ell+1$ pebbles can be placed on the board.

Consider the following reconfiguration task.
Place $\ell/2$ pebbles on $P_L$ and $\ell/2$ pebbles on $P_R$.
The goal is to transform this configuration into its reflection
across the axis determined by the edge $uv$.
More precisely, if the pebbles on $P_L$ are arranged along the path as
$p_1,p_2,\ldots,p_k$, they must be relocated to the symmetric
positions on $P_R$ in the order $p_k,\ldots,p_2,p_1$.
At the same time, the pebbles originally on $P_R$ must be moved to
the corresponding symmetric positions on $P_L$.
To analyze this process, partition the $\ell/2$ pebbles on each of
$P_L$ and $P_R$ into $m$ groups of equal size.
In order to realize the above reconfiguration, each group on $P_L$
must exchange positions near the edge $uv$ with every group on $P_R$.
Consequently, interactions may arise between essentially every pair
of these groups, giving rise to $m^2$ induced subproblems.

The main difficulty lies in the cost of preparing these subproblems
on the board tree.
Even if the subproblems are solved in a carefully chosen order,
our current intuition suggests that, at least for this example,
preparing each subproblem of size $\Theta(\ell/m)$ would require
$\Theta((\ell/m)^2)$ moves even when $m=\Ord(1)$.
Moreover, if $m=\omega(1)$, the cost of this preparation step
would likely grow significantly beyond the quadratic order. 
Since the original instance has size $\Theta(\ell)$, this intuition
suggests that each level of recursive refinement incurs a
preparation cost of
\[
\Theta((\ell/m)^2)\times\Theta(m^2)=\Theta(\ell^2)
\]
moves.
Reducing the subproblem size down to $\Ord(1)$ requires
$\Theta(\log \ell)$ levels of such refinement.
Consequently, the total number of moves spent in these preparatory
steps accumulates to $\Theta(\ell^2 \log \ell)$.
Since $\ell = \Theta(N)$ in this construction, this corresponds to
$\Theta(N^2 \log N)$ moves.

This example suggests that long degree-two corridors constitute a common 
structural source of difficulty in PMP: on trees, they appear as large isthmuses 
that induce substantial coordination costs and, when sufficiently long relative 
to the number of empty vertices, may even render the instance infeasible, 
while in general graphs, where cycles provide alternative routes, the same 
feature can lead, in extreme cases, to $\Theta(N^3)$-length solution sequences.

%%%%%%%%%%%%%%%%%%%%%%%%%%%%%%%%%%%%%%%%%%%%%%%%%%%%%%%%%%%%%%%%%%%%%%%%%%%%%%%%
%%%%%%%%%%%%%%%%%%%%%%%%%%%%%%%%%%%%%%%%%%%%%%%%%%%%%%%%%%%%%%%%%%%%%%%%%%%%%%%%
%%%%%%%%%%%%%%%%%%%%%%%%%%%%%%%%%%%%%%%%%%%%%%%%%%%%%%%%%%%%%%%%%%%%%%%%%%%%%%%%

\section*{Appendix}

\begin{rem}\label{counterexample}
The following statement on page 247, left column, second paragraph from the bottom of Kornhauser, Miller, and Spirakis \cite{715921}:
\begin{quote}
\noindent
{\bf b. G Separable and transitive}

\vspace{5pt}

``If G is a tree, then the proof of transitivity implies that at most $\Ord(n)$ moves are needed to move a pebble anywhere;
so the proof of Theorem 1, Case 2, implies an upper bound of $\Ord(n^2)$ to move all the $k < n$ pebbles.''
\end{quote}
(where $k$ denotes the number of pebbles and $n$ is the number of vertices of the board graph $G$) does not hold in general.
% The initial part of the above statement should be corrected as follows:
% \begin{quote}
%     ``If G is a tree, then the proof of transitivity implies that at most $\Ord(n^2)$ moves are needed
%     to move a pebble anywhere;''
% \end{quote}
There exist counterexamples to the initial part of the above citation, namely,
``If G is a tree, then the proof of transitivity implies that at most $\Ord(n)$ moves are needed to move a pebble anywhere;''.
This is because, when pebbles are densely packed along a path, it may be impossible
to relocate a particular pebble unless all pebbles are moved simultaneously.
For example, suppose we are given a sufficiently long path graph $P$ of even length, 
where a pendant edge $uv$ is attached to the midpoint vertex $u$ of $P$.
Let the board graph be $G:=P+uv$.
Let $P_L$ and $P_R$ denote the left and right subpaths of $P$ obtained by removing $u$, respectively, so that $P = P_L - u - P_R$.
Note that $n=|V(G)|=|V(P)|+2$ and $|V(P_L)|=|V(P_R)|=(n-2)/2$ hold. 
Suppose that all vertices of $P_L$ are occupied by pebbles without any gaps, and no other vertex contains a pebble.
In this case, the number of pebbles is $k = (n - 2)/2 = n/2 - 1$.
%For simplicity (though this is not essential to the argument), suppose that $n - 2$ is divisible by $4$.
In this configuration, by Theorem\tume{A}, this pebble motion puzzle on the board tree $G$ is feasible and thus transitive.
Consequently, all the assumptions required for Case 2 of Theorem 1 in \cite{715921} are satisfied.
In this setup, consider the task of moving the pebble $X$ initially located at the leftmost vertex of $P_L$ to the rightmost vertex of $P_R$.
This requires $\Omega(n^2)$ moves, regardless of the strategy employed.
The pebble $X$ is blocked on its right by $n/2 - 2$ pebbles, packed tightly along the path $P_L$.
%Since direct swaps of adjacent pebbles are not allowed, $X$ must be temporarily evacuated to the leaf vertex $v$ on the pendant edge $uv$ at least once during the process.
%However, to make room for $X$ to reach $v$, all the $n/2 - 2$ pebbles to its right must first be moved to the subpath $P_R$, which has $n/2 - 1$ unoccupied vertices. 
Since direct swaps of adjacent pebbles are not allowed, $X$ must be temporarily evacuated to the base vertex $u$ on the pendant edge $uv$ at least once during the process.
However, to make room for $X$ to reach $u$, all of the $n/2 - 2$ pebbles packed immediately to the right of $X$ must at some point pass through the base vertex $u$.
Indeed, this sequence of operations already demands at least 
$$
\sum_{i=1}^{n/2 - 1} i = \frac{n(n-2)}{8} \in \Omega(n^2)
$$
moves of pebbles. 
%Specifically, by the end of this process, the right half of $P_R$ will contain at least $(n - 2)/4 - 1$ pebbles.
%Since all of these pebbles originally resided on $P_L$, each must have made at least $(n - 2)/4+2$ moves to reach their positions.
%Hence, the total number of moves incurred in this stage is at least
%$$
%\left( \frac{n - 2}{4} - 1 \right) \cdot \left(\frac{n - 2}{4} +2\right)\in \Omega(n^2).
%$$
Moreover, under the setting described above, it is straightforward to verify that, starting from a configuration $f$ in which $k(=n/2 - 1)$ pebbles
labeled $a_1$ through $a_k$ are tightly packed from left to right $(a_1,\ldots, a_k)$ beginning at the leftmost end-vertex of the left subpath $P_L$
of the board graph $G$, reconfiguring them into a target configuration $g$ where the same $k$ pebbles are tightly packed from right to left
$(a_k,\ldots,a_1)$ beginning at the rightmost end-vertex of the right subpath $P_R$ of $G$, according to the algorithm proposed in \cite{715921}
requires a sequence of $\Omega(n^3)$ moves.
Although we present here only a single counterexample, % showing a lower bound  $\Omega(n^3)$ of the algorithm in \cite{715921},
it is evident that similar situations necessarily arise whenever there exists an isthmus of size $\Omega(n)$
in the $n$-vertex board tree with $\Omega(n)$ pebbles. 
% Conversely, the fact that moving a pebble to an arbitrary position can be achieved by a sequence of $\Ord(n^2)$ moves
% requires a separate proof.
% That the procedure of the algorithm in \cite{715921} indeed achieves this has been independently verified through
% replication studies in several subsequent papers \cite{WOS:001177663400002, WOS:000077614700003}. 
% Therefore, a correct revision of their argument should read:
% \begin{quote}
% \noindent
% ``If G is a tree, then the proof of transitivity implies that at most $\Ord(n^2)$ moves are needed
% to move a pebble anywhere; so the proof of Theorem 1, Case 2, implies an upper bound of
% $\Ord(n^3)$ to move all the $k < n$ pebbles.''
% \end{quote}
% That said, immediately after the above-mentioned statement in \cite{715921},
% the authors write: "The existence of biconnected subgraphs, however, can force us to an upper bound of $\Ord(n^3)$..."
% This presents a clear contrast between the case where the board graph is a tree—suggesting a complexity of $\Ord(n^2)$—and
% the case where it contains biconnected subgraphs, in which the best attainable upper bound is only $\Ord(n^3)$.
% This contrast reveals that the error identified in \cite{715921}  is not merely a typographical mistake or a momentary oversight,
% but rather stems from a fundamental misunderstanding. 
Furthermore, the facts presented here clearly indicate that,
in order to reduce the computational complexity of the problem to strictly below $\Ord(n^3)$,
it is necessary to design algorithms in a clever way that allows amortized analysis to function effectively.
% Put differently, any naive strategy—such as the one used in the algorithm of \cite{715921} ---
% that completes the reconfiguration by fixing one pebble at a time starting from the "outermost" vertices of the board graph
% and progressively reducing the size of the remaining subgraph by one at each step, cannot yield a move sequence whose length
% is truly subcubic, i.e., asymptotically better than $\Ord(n^3)$.
% Kornhauser, Miller, and Spirakis \cite{715921} の247ページの左コラムの下から２番目の段落の
% \begin{quote}
% \noindent
% {\bf b. G Separable and transitive}

% \vspace{5pt}

% If G is a tree, then the proof of transitivity implies that at most $\Ord(n)$
% moves are needed to move a pebble anywhere; so the proof of Theorem 1, Case
% 2, implies an upper bound of $\Ord(n^2)$ to move all the $k < n$ pebbles.
% \end{quote}
% という文章（ここで$k$はpebbleの数であり、$n$はボードとなるグラフの頂点数を意味している）はあからさまに間違っている。
% 彼らの上記の文章の冒頭の部分は、
% \begin{quote}
% If G is a tree, then the proof of transitivity implies that at most $\Ord(n^2)$
% moves are needed to move a pebble anywhere;    
% \end{quote}
% とすべきである。
% pebblesがパスグラフの上に一列に連なってぎっしり詰まっている状況では、そのpebbleたち全体を一斉に動かさない限り、
% その中の特定の一つのpebbleを取り出してそれを所望の場所に移動させることはできない場合がある。
% 例えば、十分に長いパス$P$の真ん中の点$u$にpendant edge$uv$が一本接続していて、$P$の$u$を除く左半分の領域にだけ隙間なくpebbleが
% 一列にびっしり詰まっている場合を考えよう。ここで$P-u$の左側のsub pathを$P_L$とし、$P-u$の右側のsub pathを$P_R$としよう。
% $P=P_L-u-P_R$であり、$P_L$上のすべての頂点の上にpebbleが置かれており、それ以外の頂点にはpebbleは置かれていない。
% その場合、$k=(n-2)/2=n/2-1$である。ついでに、本質的ではないが、簡単のため、$n-2$は$4$で割り切れる偶数であるとしよう。
% この例において、$P$の最左端の頂点に置かれているpebble $X$を、$P$の最右端の頂点の上に移動させるには、どう頑張っても
% $\Omega(n^2)$のmoveの手数が必要である。現在$X$の右隣には$n/2-2$個のpebble達がびっしりと一列に隙間なく詰まっている状況であるが、
% Xはそれらのpebble達のすべてを乗り越えて$P$の最右端の頂点に移動しなければならない。
% 隣り合うpebble同士の位置を直接入れ替えることはできないから、その場合、$P$のpendant edge $uv$の端点であるleaf $v$の上に、
% Xを（少なくとも１回以上）避難させる必要がある。
% しかしながら、それには、まず$X$の右側の$n/2-2$個の碁石すべてを、$P$上の頂点$u$の右側のsub path $P_R$
% （には$n/2-1$個の頂点の空きがある）にtemporalに移動させなければならない。
% しかしその手続きには、$\Omega(n^2)$のmoveが必要になる。
% 実際、この手続きが終了した時点で、$P_R$の右半分のsub pathには少なくとも$(n-2)/4 -1$個のpebbleが乗っているはずだが、
% これらのpebble達は元々$P_L$上に置かれていたはずであるから、その一つ一つが少なくとも$(n-2)/4$回以上のmoveを繰り返して
% この場所に到達しているはずである。
% つまりこの作業が終了した段階で、少なくとも$((n-2)/4 -1)(n-2)/4$回のmoveの手数は確定しているのである。
% 言い換えると、この作業には、$\Omega(n^2)$回のmoveの手数が必要ということになる。
% したがって、彼らの議論を正しく修正すると、
% \begin{quote}
% \noindent
% If G is a tree, then the proof of transitivity implies that at most $\Ord(n^2)$
% moves are needed to move a pebble anywhere; so the proof of Theorem 1, Case
% 2, implies an upper bound of $\Ord(n^3)$ to move all the $k < n$ pebbles.
% \end{quote}
% ということになる。すなわち、彼らが証明したのは、$\Ord(n^3)$という上界であるに過ぎない。
\end{rem}


%%%%%%%%%%%%%%%%%%%%%%%%%%%%%%%%%%%%%%%%%%%%%%%%%%%%%%%%%%%%%%%%%%%%%%%%%%%%%%%%
%%%%%%%%%%%%%%%%%%%%%%%%%%%%%%%%%%%%%%%%%%%%%%%%%%%%%%%%%%%%%%%%%%%%%%%%%%%%%%%%
%%%%%%%%%%%%%%%%%%%%%%%%%%%%%%%%%%%%%%%%%%%%%%%%%%%%%%%%%%%%%%%%%%%%%%%%%%%%%%%%

%\begin{Proofof}{Corollary~\ref{2n<N}}
%It suffices to show that, for any input board graph $G$, we can find a spanning tree 
%such that its longest isthmus consists of at most $N-n-1$ vertices, and this can be done in 
%$\Ord(N^2)$ time. 
%Let us consider the following algorithm. 
%\begin{alg-enumerate}
%\item{Find a spanning tree $T$ of $G$ that contains at least one vertex of degree 3 or higher.}\label{initial}
%\item{If $T$ has no isthmus with at least $N-n$ vertices, return $T$.}\label{check}
%\item{Let $I$ denote a unique isthmus of $T$ with at least $N-n$ vertices. 
%If every interior vertex of $I$ is a degree-$2$ vertex of $G$, then goto Step\tume\maru{\ref{big-cycle}}. 
%Find an edge of $G - T$ that is incident to an interior vertex of $I$ and denote it as $e:= uv$, where $v$ is 
%an interior vertex of $I$. 
%Let $C$ denote the unique cycle of $T+e$. 
%If $C$ has no vertex of degree $2$ in $T+e$, then let $e'$ be a single edge of $C$ that is not incident to $v$. 
%Otherwise, let $P$ be the longest path in $C$ such that all the interior vertices of $P$ have degree 2 in $T+e$, 
%and let $e'$ be one of the edges in the middle of $P$ that is not incident to $v$. 
%Set $T:=T+e-e'$ and goto Step\tume\maru{\ref{check}}.}\label{cycle}
%\item{Find an edge $e$ of $G - T$ such that the unique cycle $C$ in $T + e$ contains all of $I$.
%Take the longest path $P$ in the part of $C$ such that all the interior vertices of $P$ have degree $2$ on 
%$T+e$. Let $e'$ be one of the edges in the middle of $P$ that is not incident to $v$. Set $T:=T+e-e'$ and 
%return $T$}\label{big-cycle}
%\end{alg-enumerate}
%At Step\tume\maru{\ref{cycle}}, the maximum isthmus size of $T$ decreases strictly. 
%The proof of this is as follows:
%The statement is obviously true if $u$ is a vertex of $I$. 
%Let us consider the other case where the vertex $u$ is outside of $I$. 
%Let the two connected components obtained by deleting all interior vertices of $I$ from $G$ be denoted by $X$ and $Y$, 
%where $X$ is the component that contains vertex $u$, and $Y$ is the other component. 
%Let $p$ be the end-vertex of $I$ contained in $X$, and $q$ be the other end-vertex of $I$ contained in $Y$. 
%If $|E(P)| \geq 3$ holds, removing $e'$ will not affect the degree of any vertex with degree $3$ or more. 
%Therefore, any isthmus newly formed in $X+e-e'$ is either a path of $C$ that is edge-disjoint with $P$, 
%or a proper subpath of $P$. 
%If it is the former case, the length of the isthmus is at most $|E(P)|$, and in the latter case, it is at most $|E(P)|/2$. 
%Hence, regardless of which of the two cases it is, the number of vertices of it can only be at most $(N-1)/2$, 
%which is strictly less than $N-n$. 
%We now consider the case where the number of edges in $P$ is $2$ or fewer. 
%Let $P_1$ denote the subpath of $P$ connecting $p$ and $v$. 
%Let $X_1$ denote the connected component of $T-e'$ containing $u$, and let $X_2:=(X-e')-X_1$.
%We have that $|E(P_1)| \leq |E(P)| \leq 2$, $|V(I \cup Y)| \geq N-n+1$, and $|V(X)| \leq n$. 
%Then since $|V(X_1 + e')|\leq n$ and $|V(X_2 + P_1)| \leq n+1$ hold, the number of vertices 
%in any isthmus contained in these subgraphs is also bounded above by $N - n - 1$. 
%Therefore, in any case, if $T + e - e'$ has an isthmus with at least $N - n$ vertices, it must be a proper subpath of $I$.    
%Step\tume\maru{\ref{check}} and Step\tume\maru{\ref{cycle}} can be accomplished in $\Ord(N)$ time each, 
%while Step\tume\maru{\ref{initial}} and Step\tume\maru{\ref{big-cycle}} can be achieved in $\Ord(N^2)$ time each. 
%Since Step\tume\maru{\ref{check}} and Step\tume\maru{\ref{cycle}} are repeated at most $N/2$ times, the total time 
%complexity of this algorithm is $\Ord(N^2)$.
%\end{Proofof}

%%%%%%%%%%%%%%%%%%%%%%%%%%%%%%%%%%%%%%%%%%%%%%%%%%%%%%%%%%%%%%%%%%%%%%%%%%%%%%%%
%%%%%%%%%%%%%%%%%%%%%%%%%%%%%%%%%%%%%%%%%%%%%%%%%%%%%%%%%%%%%%%%%%%%%%%%%%%%%%%%
%%%%%%%%%%%%%%%%%%%%%%%%%%%%%%%%%%%%%%%%%%%%%%%%%%%%%%%%%%%%%%%%%%%%%%%%%%%%%%%%

\begin{Proofof}{Lemma\tume\ref{sbs0}}
We will show an algorithm of the reconstruction problem on trees for unlabeled pebbles.
%It is worth noting that in the case of unlabeled pebbles, we can echange a role of a pebble and that of an empty space.
%Hence, it is sufficient to show that there exists a sequence of moves of length $\Ord (n\D(G))$ which transfers all the pebbles from $S_1$ to $S_2$, and it can be computed in time $\Ord (N + n\D(G))$.

Let $T$ be a directed rooted tree with the root $r$.
For $v \in V(T)$, the closed descendants of $v$ is the set of vertices that are either $v$ itself or can be reached by a directed path starting from $v$.
Let us denote the closed descendants of $v$ of $T$ by ${\textrm{Des}}_{T}(v)$. 

%\begin{alg}[$\Alg.\Pac(T, \BT, x)$]\hfill
\begin{alg}[$\Alg.sbs0(G, S_1, S_2)$]\hfill
\begin{namelist}{OUTPUT}
\item[\sl Input]{A tree $G$ with $N$ vertices. Subsets of vertices $S_1, S_2  \subset V(G)$ with $|S_1| = |S_2| = n$. }
\item[\sl Output]{A sequence of moves which transfers all the pebbles from $S_1$ to $S_2$.}
\end{namelist}
\begin{alg-enumerate}
\item Choose a vertex $r$ of $G$, and build a directed rooted tree $T$ with the root $r$ such that $V(T)= V(G)$ and $E(T)=E(G)$.\label{sbs0_start0}
\item For $i=1,2$, build a directed rooted tree $T_i \subset T$ with the root $r$ such that $V(T_i) = \{ v \in V(T) \,:\, {\textrm{Des}}_{T}(v) \cap S_i \ne \emptyset \}$.
Let $\textrm{Stk}$ be a stack for delayed sequences of moves. 
\item Make a list $L$ of leaves of $T$. \label{sbs0_start1}
\item If $L = \emptyset$, then Goto step \maru{\ref{sbs0_pop}}.\label{sbs0_main}
Take a leaf $v \in L$.\\
If $v \in S_1 \cap S_2$, then:\\
$\ $ Set $S_1 = S_1 \setminus \{ v \}$, $S_2 = S_2 \setminus \{ v \}$, and update $T_1$ and $T_2$ appropriately.\\
else if $v \in S_1 \setminus S_2$, then:\\
$\ $ Let $x$ be the first vertex of $V(T_2)$ such that $x$ is on the $vr$-path.
Then there exists a vertex $y \in {\textrm{Des}}_{T_2}(x) \cap S_2$.
Let $z$ be the first vertex of $S_2$ on the $vxy$-path. 
Push the $vz$-path to $\textrm{Stk}$.\\
Set $S_1 = S_1 \setminus \{ v \}$, $S_2 = S_2 \setminus \{ z \}$, and update $T_1$ and $T_2$ appropriately.\\
else if $v \in S_2 \setminus S_1$, then:\\
$\ $ Let $x$ be the first vertex of $V(T_1)$ such that $x$ is on the $vr$-path.
Then there exists a vertex $y \in {\textrm{Des}}_{T_1}(x) \cap S_1$.
Let $z$ be the first vertex of $S_1$ on the $vxy$-path. 
Move a pebble from $z$ to $v$ in $G$.
Set $S_1 = S_1 \setminus \{ z \}$, $S_2 = S_2 \setminus \{ v \}$, and update $T_1$ and $T_2$ appropriately.\\
\item Set $T = T - v$.\label{sbs0_cut_v}
Let $u$ be the parent of $v$ in $T$.
Remove $v$ from $L$.
If $u$ is a leaf of $T$, then add $u$ to $L$.
Goto step \maru{\ref{sbs0_main}}. 
\item \label{sbs0_pop}
Pop an element, a $vz$-path, from $\textrm{Stk}$, and move a pebble from $v$ to $z$ in $G$  repeatedly, until $\textrm{Stk}$ is empty.
\end{alg-enumerate}
\end{alg}

In step \maru{\ref{sbs0_cut_v}}, $|V(T)|$ always decreases.
Hence, we have finally $L = \emptyset$.

In the algorithm, each pebble moves to a destination one by one through a path.
Since there is no other pebble on the path, the algorithm works correctly.
Moreover, the length of each path is at most $D(G)$, the total length of a sequence of moves is $\Ord (n\D(G))$.

For steps from \maru{\ref{sbs0_start0}} to \maru{\ref{sbs0_start1}}, it is sufficient to apply Depth First Search of $G$, whose running time is $\Ord (N)$.
The total running time of the algorithm is dominated by step \maru{\ref{sbs0_main}} and step  \maru{\ref{sbs0_pop}}.
In case $v \not\in S_1 \cup S_2$, it takes a constant time.
Hence, the total running time of the case is $\Ord (N)$. 
In case $v \in S_1 \cup S_2$, it takes $\Ord (\D(G))$ time to update $T_1$ and $T_2$, and it takes $\Ord (\D(G))$ steps to move a pebble.
Since this case occurs at most $n$ times throughout the entire algorithm, the total running time of the case is $\Ord (n\D(G))$.   

Therefore, the total running time of the algorithm is $\Ord (N + n\D(G))$. 
\end{Proofof} 



%%%%%%%%%%%%%%%%%%%%%%%%%%%%%%%%%%%%%%%%%%%%%%%%%%%%%%%%%%%%%%%%%%%%%%%%%%%%%%%%
%%%%%%%%%%%%%%%%%%%%%%%%%%%%%%%%%%%%%%%%%%%%%%%%%%%%%%%%%%%%%%%%%%%%%%%%%%%%%%%%
%%%%%%%%%%%%%%%%%%%%%%%%%%%%%%%%%%%%%%%%%%%%%%%%%%%%%%%%%%%%%%%%%%%%%%%%%%%%%%%%

\begin{Proofof}{Lemma\tume\ref{sbs1}}
We propose two algorithms for the proof.
The first algorithm is as follows. 

\begin{alg}[$\Alg.sbs1A(G, Q_1, Q_2)$](Partition with a rivet)\hfill
\begin{namelist}{OUTPUT}
\item[\sl Input]{A tree $G$ satisfying the assumption of Lemma\tume\ref{sbs1}.\\
A configuration $\pi \in \cF(G, P)$.\\
A partition of pebbles $P = Q_1 \cup Q_2$ such that $Q_1 \cap Q_2 = \emptyset$ and $|Q_i| = \lfloor n/2 \rfloor$ or $\lceil n/2 \rceil$ for $i=1, 2$.  
}
\item[\sl Output]{A sequence of moves which transfers $Q_i$ to $V(G_i)$ for $i = 1, 2$.}
\end{namelist}
\begin{alg-enumerate}
\item Find subsets $S_i$ of $V(G)$ for $1 \le i \le 4$, such that $|S_1| = |S_4| = \lfloor n/2 \rfloor$, $|S_2| = |S_3| = \lceil n/2 \rceil + 1$, $S_i \cap S_j \ne \emptyset$ for $1 \le i < j \le 4$, $S_1, S_2 \subset V(G_1)$, $S_3, S_4 \subset V(G_2)$, $S_2 \cap S_3 = \{ o \}$, and $S_i \cup S_{i+1}$ induces a subtree $G(i, i+1)$ of $G$ for $1 \le i \le 3$. 
\item Applying $\Alg.sbs0$, move a quarter of $P$ to $S_i$ for $1 \le i \le 4$.
Put $P_i$ as a set of pebbles on $S_i$ for $1 \le i \le 4$.
\item Consider the puzzle $(G(2,3), P_2 \cup P_3)$ and transfer half of $P_2 \cup P_3$ to $S_2$ and the other half to $S_3$ such that pebbles of $Q_1$ are moved to $S_2$ as many as possible.
Reset $P_2$ and $P_3$ as a set of pebbles on $S_2$ and $S_3$, respectively.
\item Remove $P_3 \cup P_4$ to $S_4$.
Load $P_1 \cup P_2$ to $S_2 \cup S_3$. 
Consider the puzzle $(G(2,3), P_1 \cup P_2)$ and transfer half of $P_1 \cup P_2$ to $S_2$ and the other half to $S_3$ such that pebbles of $Q_1$ are moved to $S_2$ as many as possible.
Reset $P_1$ and $P_2$ as a set of pebbles on $S_2$ and $S_3$, respectively.
\item Remove $P_1 \cup P_2$ to $S_1$.
Load $P_3 \cup P_4$ to $S_2 \cup S_3$. 
Consider the puzzle $(G(2,3), P_3 \cup P_4)$ and transfer half of $P_3 \cup P_4$ to $S_2$ and the other half to $S_3$ such that pebbles of $Q_1$ are moved to $S_2$ as many as possible.
Reset $P_3$ and $P_4$ as a set of pebbles on $S_2$ and $S_3$, respectively.
\item Remove $P_4$ to $S_4$.
Load $P_2 \cup P_3$ to $S_2 \cup S_3$. 
Consider the puzzle $(G(2,3), P_2 \cup P_3)$ and transfer half of $P_2 \cup P_3$ to $S_2$ and the other half to $S_3$ such that pebbles of $Q_1$ are moved to $S_2$ as many as possible.
Reset $P_2$ and $P_3$ as a set of pebbles on $S_2$ and $S_3$, respectively.
\end{alg-enumerate}
\end{alg}

We claim that $\Alg.sbs1A$ works correctly, it outputs a solution sequence of moves of length $\Ord (n\D(G) + n^2 \log n) $ and it can be implemented to run in time $\Ord (N + n\D(G) + n^2 \log n)$.
From steps \maru{3} to \maru{6}, $\Alg.sbs1A$ is recursively called four times with a half number of original pebbles.
The worst case is $Q_1 = P_3 \cup P_4$ initially.
In this case, after step \maru{3}, we have $Q_1 = P_2 \cup P_4$, and after step \maru{4}, we have $Q_1 = P_1 \cup P_4$, and after step \maru{5}, we have $Q_1 = P_1 \cup P_3$, and finally after step \maru{6}, we have $Q_1 = P_1 \cup P_2$.

Step \maru{2} can be implemented in $\Ord(N+n\D(G))$ time by Lemma\tume\ref{sbs0}.
Let us denote the running time of the algorithm from step \maru{3} to \maru{6} by $g_1(n)$.
We want to prove $g_1(n) \le c_1 n^2 \log_2 n$ with some constant $c_1$.
we proceed by induction on $n$.
The number of steps of the removal, loading, and restoration phase is at most $cn^2$ with some constant $c$, since $|S_i| = n/2 + \Ord (1)$ for $1 \le i \le 4$.
Then, by the inductive hypothesis, we have
\begin{eqnarray*}
g_1(n) & \le & 4 g_1(\frac{n}{2}) + c n^2 \\
 & \le & 4 c_1 (\frac{n}{2})^2 \log_2 (\frac{n}{2}) + c n^2 \\
 & = & c_1 n^2 \log_2 n - c_1 n^2 + c n^2.
\end{eqnarray*} 
By choosing $c_1 = c$, we have $g_1(n) \le c n^2 \log_2 n$, as claimed.

Next, we propose the second algorithm. 

\begin{alg}[$\Alg.sbs1B(G, \pi_1, \pi_2)$](Sort with a rivet)\hfill
\begin{namelist}{OUTPUT}
\item[\sl Input]{A tree $G$ satisfying the assumption of Lemma\tume\ref{sbs1}.\\
Configurations $\pi_1$, $\pi_2 \in \cF(G, P)$.
}
\item[\sl Output]{A sequence of moves which transforms $\pi_1$ to $\pi_2$.}
\end{namelist}
\begin{alg-enumerate}
\item Find subsets $S_i$ of $V(G)$ for $1 \le i \le 4$, such that $|S_1| = |S_4| = \lfloor n/2 \rfloor$, $|S_2| = |S_3| = \lceil n/2 \rceil + 1$, $S_i \cap S_j \ne \emptyset$ for $1 \le i < j \le 4$, $S_1, S_2 \subset V(G_1)$, $S_3, S_4 \subset V(G_2)$, $S_2 \cap S_3 = \{ o \}$, and $S_i \cup S_{i+1}$ induces a subtree $G(i, i+1)$ of $G$ for $1 \le i \le 3$. 
%Let $S = \cup_{1 \le i \le 4} S_i$. 
\item Choose a configuration $\pi_i^{\ast}$ for $i = 1, 2$ such that $|{\rm sup}(\pi_i^{\ast}) \cap S_j| = \lfloor n/4 \rfloor$ or $\lceil n/4 \rceil$ for $1 \le j \le 4$ and $\pi_i^{\ast}$ is strongly related to $\pi_i$ for $i = 1, 2$.
Transform $\pi_1$ to $\pi_1^{\ast}$.
\item Transfer $\pi_2^{\ast}(S_1 \cup S_2)$ to $S_1 \cup S_2$, and transfer $\pi_2^{\ast}(S_3 \cup S_4)$ to $S_3 \cup S_4$ by applying $\Alg.sbs1A$.
Reset $P_i$ as a set of pebbles on $S_i$ for $1 \le i \le 4$.
\item Remove $P_3 \cup P_4$ from $S_3 \cup S_4$ to $S_4$.
Load $P_1 \cup P_2$ to $S_2 \cup S_3$.
Consider the puzzle $(G(2,3), P_1 \cup P_2)$ and build a configuration $\pi_{2,1}^{\ast}$
such that $\pi_{2,1}^{\ast}$ and $\pi_2^{\ast}|_{S_1 \cup S_2}$ are strongly related.
\item Remove $P_1 \cup P_2$ from $S_2 \cup S_3$ to $S_1$.
Load $P_3 \cup P_4$ to $S_2 \cup S_3$.
Consider the puzzle $(G(2,3), P_3 \cup P_4)$ and build a configuration $\pi_{2,2}^{\ast}$ 
such that $\pi_{2,2}^{\ast}$ and $\pi_2^{\ast}|_{S_3 \cup S_4}$ are strongly related.
\item Make $\pi_2^{\ast}$.
\item Transform $\pi_2^{\ast}$ to $\pi_2$.
\end{alg-enumerate}
\end{alg}

We claim that $\Alg.sbs1B$ works correctly, it outputs a solution sequence of moves of length $\Ord (n\D(G) + n^2 \log n) $ and it can be implemented to run in time $\Ord (N + n\D(G) + n^2 \log n)$.

In step \maru{2} and step \maru{7}, since $\pi_i$ and $\pi_i^{\ast}$ are strongly related for $i = 1, 2$, 
the number of steps to transform $\pi_1$ to $\pi_1^{\ast}$ and to transform $\pi_2^{\ast}$ to $\pi_2$ are at most $\Ord (n\D(G))$, and it cab be  computed in time $\Ord (N + n\D(G)$ by applying $\Alg.sbs0$.
The number of steps of the removal, loading, and restoration phase is at most $cn^2$ with some constant $c$.
In step \maru{3}, by applying $\Alg.sbs1A$, we have at most $c_1 n^2 \log n$ time.
Let $f_1(n)$ be the algorithm's running time from step \maru{3} to step \maru{6}.
We want to prove $f_1(n) \le d_1 n^2 \log n$ with some constant $d_1$.
In step \maru{4} and step \maru{5}, by the inductive hypothesis, we have
\begin{eqnarray*}
f_1(n) & \le & 2 f_1(\frac{n}{2}) + c_1 n^2 \log n + c n^2 \\
 & \le & 2 d_1 (\frac{n}{2})^2 \log_2 (\frac{n}{2}) + c_1 n^2 \log n + c n^2 \\
 & = & (c_1 + \frac{d_1}{2} ) n^2 \log_2 n + (c - \frac{d_1}{2} ) n^2.
\end{eqnarray*} 
By choosing $d_1$ such that $d_1 \ge 2\max\{ c, c_1 \}$, we have $f_1(n) \le d_1 n^2 \log_2 n$.        
\end{Proofof} 

%%%%%%%%%%%%%%%%%%%%%%%%%%%%%%%%%%%%%%%%%%%%%%%%%%%%%%%%%%%%%%%%%%%%%%%%%%%%%%%%
%%%%%%%%%%%%%%%%%%%%%%%%%%%%%%%%%%%%%%%%%%%%%%%%%%%%%%%%%%%%%%%%%%%%%%%%%%%%%%%%
%%%%%%%%%%%%%%%%%%%%%%%%%%%%%%%%%%%%%%%%%%%%%%%%%%%%%%%%%%%%%%%%%%%%%%%%%%%%%%%%

\begin{Proofof}{Lemma\tume\ref{sbs2}}
We propose two algorithms for the proof.
The first algorithm is as follows. 

\begin{alg}[$\Alg.sbs2A(G, Q_1, Q_2)$](Partition with a long isthmus)\hfill
\begin{namelist}{OUTPUT}
\item[\sl Input]{A tree $G$ satisfying the assumption of Lemma\tume\ref{sbs2}.\\
A configuration $\pi \in \cF(G, P)$.\\
A partition of pebbles $P = Q_1 \cup Q_2$ such that $Q_1 \cap Q_2 = \emptyset$ and $|Q_i| = \lfloor n/2 \rfloor$ or $\lceil n/2 \rceil$ for $i=1, 2$.  
}
\item[\sl Output]{A sequence of moves which transfers $Q_i$ to $V(G_i)$ for $i = 1, 2$.}
\end{namelist}
\begin{alg-enumerate}
\item Find subsets $S_i$ of $V(G)$ for $1 \le i \le 4$, such that $|S_i| = \lfloor n/4 \rfloor$ or $\lceil n/4 \rceil$ for $1 \le i \le 4$, $S_i \cap S_j \ne \emptyset$ for $1 \le i < j \le 4$, $S_1, S_2 \subset V(G_1)$, $S_3, S_4 \subset V(G_2)$, and $S_i \cup S_{i+1} \cup V(I)$ induces a subtree $G(i, i+1)$ for $1 \le i \le 3$.
%Let $v_1$ and $v_2$ be the endvertices of $I$ such that $v_i \in V(G_i)$ for $i=1, 2$. 
\item Applying $\Alg.sbs0$, move a quarter of $P$ to $S_i$ for $1 \le i \le 4$.
Put $P_i$ as a set of pebbles on $S_i$ for $1 \le i \le 4$.
\item Consider the puzzle $(G(2,3), P_2 \cup P_3)$ and transfer half of $P_2 \cup P_3$ to $S_2$ and the other half to $S_3$ such that pebbles of $Q_1$ are moved to $S_2$ as many as possible.
Reset $P_2$ and $P_3$ as a set of pebbles on $S_2$ and $S_3$, respectively.
\item Consider the puzzle $(G(1,2), P_1 \cup P_2)$ and transfer half of $P_1 \cup P_2$ to $S_1$ and the other half to $S_2$ such that pebbles of $Q_1$ are moved to $S_1$ as many as possible.
Reset $P_1$ and $P_2$ as a set of pebbles on $S_1$ and $S_2$, respectively.
\item Consider the puzzle $(G(3,4), P_3 \cup P_4)$ and transfer half of $P_3 \cup P_4$ to $S_3$ and the other half to $S_4$ such that pebbles of $Q_1$ are moved to $S_3$ as many as possible.
Reset $P_3$ and $P_4$ as a set of pebbles on $S_3$ and $S_4$, respectively.
\item Consider the puzzle $(G(2,3), P_2 \cup P_3)$ and transfer half of $P_2 \cup P_3$ to $S_2$ and the other half to $S_3$ such that pebbles of $Q_1$ are moved to $S_2$ as many as possible.
Reset $P_2$ and $P_3$ as a set of pebbles on $S_2$ and $S_3$, respectively.
\end{alg-enumerate}
\end{alg}

We claim that $\Alg.sbs2A$ works correctly, it outputs a solution sequence of moves of length $\Ord (n\D(G) + n^2 \log n) $ and it can be implemented to run in time $\Ord (N + n\D(G) + n^2 \log n)$.

In step \maru{4} and step \maru{5}, $\Alg.sbs1A$ is called two times with a half number of original pebbles, and 
in step \maru{3} and step \maru{6}, $\Alg.sbs2A$ is recursively called two times with a half number of original pebbles.
The worst case is $Q_1 = P_3 \cup P_4$ initially.
In this case, after step \maru{3}, we have $Q_1 = P_2 \cup P_4$, and after step \maru{4}, we have $Q_1 = P_1 \cup P_4$, and after step \maru{5}, we have $Q_1 = P_1 \cup P_3$, and finally after step \maru{6}, we have $Q_1 = P_1 \cup P_2$.

The running time $g_2(n)$ of the algorithm is dominated by step \maru{3} - step \maru{6}.
We want to prove $g_2(n) \le n\D(G) + c_2 n^2 \log_2 n$ with some constant $c_2$.
we proceed by induction on $n$.
The number of steps of removal, loading, and restoration phase is at most $cn^2$ with some constant $c$, since $|S_1 \cup S_2| = n/2 + \Ord (1)$ and $|S_3 \cup S_4| = n/2 + \Ord (1)$.
The cost of step \maru{4} and step \maru{5} is at most $c_1 (n/2)^2 \log_2 (n/2)$ by Lemma\tume\ref{sbs1}.

Then, by the inductive hypothesis, we have
\begin{eqnarray*}
g_2(n) & \le & 2 g_2(\frac{n}{2}) + 2c_1 (\frac{n}{2})^2 \log_2 (\frac{n}{2}) + c n^2 \\
 & \le & 2 ( \frac{n}{2}D(G) + c_2 (\frac{n}{2})^2 \log_2 (\frac{n}{2}) ) + 2c_1 (\frac{n}{2})^2 \log_2 (\frac{n}{2}) + c n^2 \\
 & \le & n\D(G) + (\frac{c_2}{2} + \frac{c_1}{2}) n^2 \log_2 n + (c - \frac{c_2}{2} - \frac{c_1}{2} )n^2.
\end{eqnarray*} 
By choosing $c_2$ such that $c_2 \ge \max \{ c_1, 2c - c_1 \}$, we have $g_2(n) \le n\D(G) + c_2 n^2 \log_2 n$, as claimed.              

Next, we propose the second algorithm. 

\begin{alg}[$\Alg.sbs2B(G, \pi_1, \pi_2)$](Sort with a long isthmus)\hfill
\begin{namelist}{OUTPUT}
\item[\sl Input]{A tree $G$ satisfying the assumption of Lemma\tume\ref{sbs2}.\\
Configurations $\pi_1$, $\pi_2 \in \cF(G, P)$.
}
\item[\sl Output]{A sequence of moves which transforms $\pi_1$ to $\pi_2$.}
\end{namelist}
\begin{alg-enumerate}
\item Find subsets $S_i$ of $V(G)$ for $1 \le i \le 4$, such that $|S_i| = \lfloor n/4 \rfloor$ or $\lceil n/4 \rceil$ for $1 \le i \le 4$, $S_i \cap S_j \ne \emptyset$ for $1 \le i < j \le 4$, $S_1, S_2 \subset V(G_1)$, $S_3, S_4 \subset V(G_2)$, and $S_i \cup S_{i+1} \cup V(I)$ induces a subtree $G(i, i+1)$ for $1 \le i \le 3$.
%Let $v_1$ and $v_2$ be the endvertices of $I$ such that $v_i \in V(G_i)$ for $i=1, 2$. 
\item Choose $\pi_i^{\ast}$ for $i = 1, 2$ such that $|{\rm sup}(\pi_i^{\ast}) \cap S_j| = \lfloor n/4 \rfloor$ or $\lceil n/4 \rceil$ for $1 \le j \le 4$ and $\pi_i^{\ast}$ is strongly related to $\pi_i$ for $i = 1, 2$.
Transform $\pi_1$ to $\pi_1^{\ast}$.
\item Consider the puzzle $(G, P)$ and transfer $\pi_2^{\ast}(S_1 \cup S_2)$ to $S_1 \cup S_2$ and transfer $\pi_2^{\ast}(S_3 \cup S_4)$ to $S_3 \cup S_4$ by applying $\Alg.sbs2A$.
Reset $P_i$ as a set of pebbles on $S_i$ for $1 \le i \le 4$.
\item 
Consider the puzzle $(G(1,2), P_1 \cup P_2)$ and build a configuration $\pi_{2,1}^{\ast}$ 
such that $\pi_{2,1}^{\ast}$ and $\pi_2^{\ast}|_{S_1 \cup S_2}$ are strongly related.
\item
Consider the puzzle $(G(3,4), P_3 \cup P_4)$ and build a configuration $\pi_{2,2}^{\ast}$ 
such that $\pi_{2,2}^{\ast}$ and $\pi_2^{\ast}|_{S_3 \cup S_4}$ are strongly related.
\item Make $\pi_2^{\ast}$.
\item Transform $\pi_2^{\ast}$ to $\pi_2$.
\end{alg-enumerate}
\end{alg}

We claim that $\Alg.sbs2B$ works correctly, it outputs a solution sequence of moves of length $\Ord (n\D(G) + n^2 \log n) $ and it can be implemented to run in time $\Ord (N + n\D(G) + n^2 \log n)$.

In step \maru{2} and step \maru{7}, since $\pi_i$ and $\pi_i^{\ast}$ are strongly related for $i = 1, 2$, the number of steps to transfer $\pi_1$ to $\pi_1^{\ast}$ and to transfer $\pi_2^{\ast}$ to $\pi_2$ are at most $\Ord (n\D(G))$ by applying $\Alg.sbs0$.
In step \maru{3}, $\Alg.sbs2A$ is called, and 
in step \maru{4} and step \maru{5}, $\Alg.sbs1B$ is called two times with a half number of original pebbles.

The running time $f_2(n)$ of the algorithm is dominated by step \maru{3} - step \maru{6}.
We want to prove $f_2(n) \le n\D(G) + d_2 n^2 \log_2 n$ with some constant $d_2$.
The number of steps of removal, loading, and restoration phase is at most $cn^2$ with some constant $c$.
The cost of  step \maru{3} is at most $Nn + c_2 n^2 \log_2 n$ by $\Alg.sbs2A$.
The cost of  steps \maru{4} and \maru{5} is at most $d_1 (n/2)^2 \log_2 (n/2)$ by $\Alg.sbs1B$.
Then, we have
\begin{eqnarray*}
f_2(n) & \le & n\D(G) + c_2 n^2 \log_2 n + 2 d_1 (\frac{n}{2})^2 \log_2 (\frac{n}{2}) + c n^2 \\
 & \le & n\D(G) + ( c_2 + \frac{d_1}{2}) n^2 \log_2 n + (c - \frac{d_1}{2})n^2.
\end{eqnarray*} 
By choosing $d_2$ such that $d_2 \ge c_2 + c$, we have $f_2(n) \le n\D(G) + d_2 n^2 \log_2 n$.              
\end{Proofof}

%%%%%%%%%%%%%%%%%%%%%%%%%%%%%%%%%%%%%%%%%%%%%%%%%%%%%%%%%%%%%%%%%%%%%%%%%%%%%%%%
%%%%%%%%%%%%%%%%%%%%%%%%%%%%%%%%%%%%%%%%%%%%%%%%%%%%%%%%%%%%%%%%%%%%%%%%%%%%%%%%
%%%%%%%%%%%%%%%%%%%%%%%%%%%%%%%%%%%%%%%%%%%%%%%%%%%%%%%%%%%%%%%%%%%%%%%%%%%%%%%%
\begin{Proofof}{Theorem\tume\ref{sbs3}}
Since $G$ is a tree, $G$ has a centroid vertex or a centroid edge.
We consider two cases.
\medskip\\
\underline{Case 1.} There exists a centroid $o$ of degree at least $3$.
\medskip\\
In this case, since $N \ge 3n$, $G$ has subtrees $G_1$ and $G_2$ such that $V(G_1) \cap V(G_2) = \{ o \}$ and $|V(G_i)| \ge n + 1$ for $i=1, 2$.
Hence, by Lemma\tume\ref{sbs1}, the proof is finished.
\medskip\\
\underline{Case 2.} There exists no centroid of degree at least $3$.
\medskip\\  
In this case, the centroid is contained in an isthmus $I$.
Let $v_1$ and $v_2$ be endvertices of $I$, and let $V_0$ be the set of inner vertices of $I$.
We have two subtrees $G_1$ and $G_2$ induced by $V(G) \setminus V_0$ such that $v_i \in V(G_i)$ for $i = 1, 2$.
Put $N_i = |V(G_i)|$ for $i = 1, 2$.
By the feasibility of the puzzle, we have $n \le N_1 + N_2 - 3$.
We may assume $N_1 \le N_2$. 
Let $k= |V_0|$. 
\medskip\\
\underline{Subcase 2. 1.} $k < n$.
\medskip\\  
In this case, we have $N_2 \ge |V(G) \setminus V_0|/2 = (N-k)/2 > (N - n)/2 \ge n$.
On the other hand, since $I$ contains the centroid, we have $|V(G_1) \cup V(I)| \ge N/2 \ge (3/2)n$.
Hence, by taking $v_2$ as the rivet in the assumption of Lemma\tume\ref{sbs1}, we can apply Lemma\tume\ref{sbs1} to finish the proof.
\medskip\\
\underline{Subcase 2. 2.} $k \ge n$.
\medskip\\
Let $H$ be a subtree of $G$ induced by $V(I) \cup V(G_2)$, and 
let $K$ be a subtree of $G$ induced by $V(G_1) \cup V(I) \cup V(G'_2)$, 
where $G'_2$ is a subtree of $G_2$ containing $v_2$ with $N_1$ vertices.
Let us take positive integers $n_1$, $n_2$ such that $n_1 + n_2 = n$ and $n_i \le N_i - 1$ for $i=1,2$.  
We will use $\Alg.sbs1B$ and $\Alg.sbs2B$ to solve the puzzle $(G,P)$ as follows;
\begin{alg-enumerate}
\item{Apply $\Alg.sbs1B$ for the puzzle $(H, n_2)$ with a rivet $v_2$.}
\item{Apply $\Alg.sbs2B$ for the puzzle $(K, 2n_1)$ with an isthmus $I$.}
\item{Again, apply $\Alg.sbs1B$ for the puzzle $(H, n_2)$ with a rivet $v_2$.}
\end{alg-enumerate}
The running time of steps \maru{1} - \maru{3} is $\Ord(N + n\D(G) + n^2\log n)$ by Lemma\tume\ref{sbs1} and Lemma\tume\ref{sbs2}.
This completes the proof.
\end{Proofof}

\begin{Proofof}{Lemma\tume\ref{sbs4}}
Since $G$ is a tree, $G$ has a centroid vertex or a centroid edge.
We consider two cases.
\medskip\\
\underline{Case 1.} There exists a centroid $o$ of degree at least $3$.
\medskip\\
In this case, since $N \ge 3n$, $G$ has subtrees $G_1$ and $G_2$ such that $V(G_1) \cap V(G_2) = \{ o \}$ and $|V(G_i)| \ge n + 1$ for $i=1, 2$.
Since $\deg \, o \ge 3$, we may assume that there exists a set of $3$ neighbors $v_1, v_2, v_3$ of $o$ such that $v_1 \in G_1$ and $v_2, v_3 \in G_2$.
Set $S = \{ o, v_1, v_2, v_3 \}$.

We will use the following algorithm;
\begin{alg-enumerate}
\item Starting from $\pi$, gather all the pebbles in $V(G_1)$.
\item Move pebbles to $V(G_2)$ through $o$.
If $p$ or $q$ passes through $o$, stop them at $v_2$ or $v_3$.
Let all pebbles except $p$ and $q$ pass through $S$ and put them in $V(G_2) \setminus \{ v_2, v _3 \}$, until both $p$ and $q$ arrive at $S$.
\item Interchange $p$ and $q$ in $S$.
\item Follow all of the steps before interchanging $p$ and $q$  in reverse order to achieve the desired configuration $\pi(p\leftrightarrow q)$.
\end{alg-enumerate}
\noindent
\underline{Case 2.} There exists no centroid of degree at least $3$.
\medskip\\  
In this case, the centroid is contained in an isthmus $I$.
Let $v_1$ and $v_2$ be endvertices of $I$, and let $V_0$ be the set of inner vertices of $I$.
We have two subtrees $G_1$ and $G_2$ induced by $V(G) \setminus V_0$ such that $v_i \in V(G_i)$ for $i = 1, 2$.
Put $N_i = |V(G_i)|$ for $i = 1, 2$.
By the feasibility of the puzzle, we have $n \le N_1 + N_2 - 3$.
We may assume $N_1 \le N_2$. 
Let $k= |V_0|$. 
\medskip\\
\underline{Subcase 2. 1.} $k < n$.
\medskip\\  
In this case, we have $N_2 \ge |V(G) \setminus V_0|/2 = (N-k)/2 > (N - n)/2 \ge n$.
On the other hand, since $I$ contains the centroid, we have $|V(G_1) \cup V(I)| \ge N/2 \ge (3/2)n$.
Hence, by taking $v_2$ as the vertex $o$ in Case 2.1, we can apply the same algorithm in Case 2.1 to finish the proof.
\medskip\\
\underline{Subcase 2. 2.} $k \ge n$.
\medskip\\
Let $H$ be a subtree of $G$ induced by $V(I) \cup V(G_2)$, and 
let $K$ be a subtree of $G$ induced by $V(G_1) \cup V(I) \cup V(G'_2)$, 
where $G'_2$ is a subtree of $G_2$ containing $v_2$ with $N_1$ vertices.
Let us take positive integers $n_1$, $n_2$ such that $n_1 + n_2 = n$ and $n_i \le N_i - 1$ for $i=1,2$.  

Since $\deg \, v_i \ge 3$ for $i=1,2$, we may assume that there exists a set of vertices $u_{1,1}, u_{1,2}, u_{2,1}, u_{2,2}$ such that $u_{i,1}, u_{i,2}$ are neighbors of $v_i$ in $G_i$ for $i=1,2$.
Set $S = V(I) \cup \{ u_{1,1}, u_{1,2}, u_{2,1}, u_{2,2} \}$.

We will use the following algorithm;
\begin{alg-enumerate}
\item Starting from $\pi$, gather all the pebbles in $V_0$.
\item Move $n_i$ pebbles to $V(G_i)$ through $v_i$ for $i=1,2$.
If $p$ or $q$ passes through $v_i$, stop them at $u_{i,1}$ or $u_{i,2}$.
Let all pebbles except $p$ and $q$ pass through $v_1$ or $v_2$ and put them in $(V(G_1) \cup V(G_2))\setminus \{ u_{1,1}, u_{1,2}, u_{2,1}, u_{2,2} \}$.
\item If both $p$ and $q$ arrive at common $v_i$ for $i=1$ or $2$, then interchange $p$ and $q$ around $v_i$.  
Otherwise, we may assume $p$ arrives at $v_1$ and $q$ arrives at $v_2$.
Let all pebbles except $p$ and $q$ pass through $v_1$ or $v_2$ and put them in $(V(G_1) \cup V(G_2))\setminus \{ u_{1,1}, u_{1,2}, u_{2,1}, u_{2,2} \}$, then interchange $p$ and $q$ in $S$.
\item Follow all of the steps before interchanging $p$ and $q$ in reverse order to achieve the desired configuration $\pi(p\leftrightarrow q)$.
\end{alg-enumerate}
In all cases, the total running time of steps is $\Ord(N + n\D(G))$, since each pebble moves $\Ord(\D(G))$ times in the algorithms.
This completes the proof.
\end{Proofof}

%%%%%%%%%%%%%%%%%%%%%%%%%%%%%%%%%%%%%%%%%%%%%%%%%%%%%%%%%%%%%%%%%%%%%%%%%%%%%%%%
%%%%%%%%%%%%%%%%%%%%%%%%%%%%%%%%%%%%%%%%%%%%%%%%%%%%%%%%%%%%%%%%%%%%%%%%%%%%%%%%
%%%%%%%%%%%%%%%%%%%%%%%%%%%%%%%%%%%%%%%%%%%%%%%%%%%%%%%%%%%%%%%%%%%%%%%%%%%%%%%%

\begin{Proofof}{Lemma\tume\ref{centroid-splitting}}
If $T$ has a centroid edge, then by definition the statement of this lemma clearly holds.
Hence let us assume that $T$ has a centroid vertex $o$.
Let us consider the set of components $\cS$ of $T-o$. 
Without loss of generality, let $X$ be a graph with the fewest order, 
and $Y$ be a graph with the second-fewest order in the set $\cS$. 
Of course it is possible that $|V(X)|=|V(Y)|$. 
If $m \geq 3$, since $|V(X)|+|V(Y)| \leq \frac{2(N-1)}{3}$ holds, let us denote  
$Z:=X \cup Y$ and update $\cS:=(\cS \setminus \{X,Y\})\cup \{Z\}$. 
Repeating this operation, the size of the set S becomes $2$. 
At that stage, let us denote $S=\{G,H\}$, and let us assume $|V(G)|\leq |V(H)|$.
In this case, we can take $T_1$ to be the subgraph of $T$ induced by 
$V(G) \cup \{o\}$, and similarly $T_2$ to be the subgraph of $T$ induced by 
$V(H) \cup \{o\}$. 
Since the size of each element of the set $\cS$ is between $1$ and $n$, 
we can sort and manage each element of $\cS$ in order of its size using 
a $1\times n$-array $A(\cS)$ whose each cell is a list. 
Since the minimum of the orders of the graphs in $\cS$ is 
monotonically non-decreasing, the task of identifying a graph with 
the smallest order and a graph with the second smallest order in $A(\cS)$
takes only $\Ord(N)$ effort in total.
\end{Proofof}

%%%%%%%%%%%%%%%%%%%%%%%%%%%%%%%%%%%%%%%%%%%%%%%%%%%%%%%%%%%%%%%%%%%%%%%%%%%%%%%%
%%%%%%%%%%%%%%%%%%%%%%%%%%%%%%%%%%%%%%%%%%%%%%%%%%%%%%%%%%%%%%%%%%%%%%%%%%%%%%%%
%%%%%%%%%%%%%%%%%%%%%%%%%%%%%%%%%%%%%%%%%%%%%%%%%%%%%%%%%%%%%%%%%%%%%%%%%%%%%%%%

\begin{Proofof}{Lemma\tume\ref{feasible-subtree}}
%$T^{\prime}$が長さ$1$のパス$(K_2)$である場合には、$T^{\prime}$のいずれかの点を$v$と置き、$H:=T$とする。それ以外の場合には、$T^{\prime}$のleaf以外のすべての点を$T$上で$1$頂点$v$にまで縮約した木を$H$と置く。さらに$T^{\prime}$に対応する$H$の部分グラフを$X$と置く。
%$X$に属していない$v$の隣接頂点$u$が$H$に見つかる限り、その頂点$u$と辺$vu$を$X$に追加（$X:=X;vu$に更新）し続ける。
%$X$に属していない$v$の隣接頂点$u$が$H$に存在しなくなったら、$X$のどこかのleaf $\ell$に隣接する$X$外の頂点$w$と辺$\ell{w}$を$X$に追加する。その結果$X$に新しく出来た長さ$2$のarm $v-\ell-w$の$v$に接続する辺$v\ell$を縮約して$\ell$を$v$に同一視したものを新たに$X$に置き直し、$H$上で$v\ell$を縮約した結果得られる新たなグラフを$H$に置き直して、$v$の隣接頂点を$X$に追加するフェイズに戻る。
Let us consider the following algorithm:
%下記の一連の作業を$X$に対応する$T$の点部分グラフの位数が$d$になったところで止めて、復元して出力すればよい。
\begin{alg-enumerate}
\item{If $T^{\prime}$ consists of a single edge $e$, let $v$ be an end-vertex of $e$ and set $H:=T$. 
Otherwise, let $H$ be the graph obtained by contracting all parts of $T^{\prime}$ except the pendant edges 
to a single vertex $v$. Let $X$ be the subgraph of $H$ corresponding to $T^{\prime}$. }
%\item{$T^{\prime}$が長さ$1$のパス$(K_2)$である場合には、$T^{\prime}$のいずれかの点を$v$と置き、$H:=T$とする。それ以外の場合には、$T^{\prime}$のleaf以外のすべての点を$T$上で$1$頂点$v$にまで縮約した木を$H$と置く。さらに$T^{\prime}$に対応する$H$の部分グラフを$X$と置く。}
\item{Continue to add vertex $u$ and edge $vu$ to $X$ (i.e., update to $X:=X+vu$) as long as 
the rank of the subgraph $T(X)$ of $T$ corresponding to $X$ is strictly less than $d$, and 
a neighboring vertex $u$ of $v$ that does not belong to $X$ is found in $H$.}\label{leaf-add}
%\item{$X$に属していない$v$の隣接頂点$u$が$H$に見つかる限り、その頂点$u$と辺$vu$を$X$に追加（$X:=X+vu$に更新）し続ける。}\label{leaf-add}
\item{If the rank of the subgraph $T(X)$ of $T$ corresponding to $X$ is strictly less than $d$, do:
\begin{alg-enumerate}
\item{Find a vertex $w$ outside $X$ adjacent to some leaf $\ell$ in $X$ and add that vertex $w$ and edge $\ell{w}$ to $X$. 
The resulting $X$ has the length $2$ arm $v-\ell-w$.}
\item{Update $X$ (resp. $H$) as the graph obtained by contracting the edge $v\ell$ of $X$ (resp. $H$) and identifying $\ell$ with $v$. }
\item{Go to Step\tume\maru{\ref{leaf-add}}}
\end{alg-enumerate}
}\label{spider}
%\item{$X$のどこかのleaf $\ell$に隣接する$X$外の頂点$w$を見つけて、その頂点$w$と辺$\ell{w}$を$X$に追加する。その結果$X$に生じた長さ$2$のarm $v-\ell-w$上の辺$v\ell$を縮約して$\ell$を$v$と同一視させたグラフを$X$に取り直し、$H$の辺$v\ell$を縮約したグラフを$H$に取り直して、ステップ\maru{\ref{leaf-add}}に戻る。}\label{spider}
\item{Return the subgraph $T(X)$ of $T$ corresponding to $X$.}
\end{alg-enumerate}


Let us show the correctness of the above algorithm. 
%上記のアルゴリズムの正当性について述べる。
To do so, we need only check the fact that the maximal isthmus size of the subgraph $T(X)$ of $T$ corresponding to $X$ 
does not exceed the maximal isthmus size of $T$ by the operation of Step\tume{\maru{\ref{leaf-add}}} or Step\tume{\maru{\ref{spider}}}.
%それには、$X$に対応する$T$の部分グラフ$T(X)$の最大isthmus長が、ステップ{\maru{\ref{leaf-add}}}やステップ{\maru{\ref{spider}}}の操作によって
%$T$の最大isthmus長を上回ることがない、という事実を確認すればよい。
First, for the operation of Step\tume\maru{\ref{leaf-add}}, the maximum isthmus size of $T(X)$ is not increased, 
since the degree of the non-leaf vertices in $T(X)$ is increased.
%まず、ステップ\maru{\ref{leaf-add}}の操作は、$T(X)$におけるleaf以外の頂点の次数しか増やさないので、最大isthmus長は減ることはあっても増えることはない。
Next, the vertex $w$ that is added to $T$ in Step\tume\maru{\ref{spider}} operation is not a cut vertex because it is a leaf of $T(X)$, 
so it cannot be part of any isthmus in $T(X)$.
%次に、ステップ\maru{\ref{spider}}の操作で$T(X)$に付け加える頂点$w$は$T(X)$のleafなのでcut vertexではないから、$T(X)$のisthmusを
%構成する頂点ではあり得ない。
Therefore, the new isthmus arising in $T(X)$ at this step has one end vertex at $\ell$. 
%したがって、この段階で$T(X)$に生じる新たなisthmusは一方の端点を$\ell$に持つ。
On the other hand, just before entering the operation of Step\tume\maru{\ref{spider}}, 
the degree on $T(X)$ and the degree on $T$ are the same for any vertex other than leaves 
in $T(X)$, so this newly generated isthmus in $T(X)$ is just an isthmus of $T$ itself. 
%一方、ステップ\maru{\ref{spider}}の操作に入る直前の状態においては、$T(X)$におけるleaf以外のどの頂点についても、$T(X)$上の次数と$T$上の次数は一致しているので、$T(X)$に新たに生じたこのisthmusは$T$のisthmusに過ぎない。
It is clear that the time complexity of this algorithm is linear. 
%なお、このアルゴリズムの計算量が線形時間であることは明らかである。
\end{Proofof}

%%%%%%%%%%%%%%%%%%%%%%%%%%%%%%%%%%%%%%%%%%%%%%%%%%%%%%%%%%%%%%%%%%%%%%%%%%%%%%%%
%%%%%%%%%%%%%%%%%%%%%%%%%%%%%%%%%%%%%%%%%%%%%%%%%%%%%%%%%%%%%%%%%%%%%%%%%%%%%%%%
%%%%%%%%%%%%%%%%%%%%%%%%%%%%%%%%%%%%%%%%%%%%%%%%%%%%%%%%%%%%%%%%%%%%%%%%%%%%%%%%

\begin{Proofof}{Lemma\tume\ref{bunkatsu}}
Let us consider the operation $X$ mentioned in Lemma\tume\ref{centroid-splitting} of dividing a given tree into two subtrees. 
%$T$のcentroid vertexもしくはcentroid edgeを境にして、小さい方の部分木の位数が$n/3$以上となるような2つの部分木$T_1, T_2$（一般性を失わず、$|V(T_1)|=<|V(T_2)|$を仮定）に分けるという操作$X$について考えます。
Let us recursively apply the operation $X$ to each subtree of $T$ to the point such that the order of each subtree is greater than or equal to $k$ and that another $X$ on those subtrees will reduce the order of the resulting subtrees to strictly less than $k$.
%上記の「centroidを境に2つの部分木に分ける操作$X$」を、$T$の各部分木が「位数$k$を下回らないギリギリのサイズ」になるまで再帰的に適用します。

Consequently, if $r$ is the number of subtrees obtained so far, then $2n/k > (n-1)/(k-1) \geq r$.
%その結果、ここまでに得られた部分木の個数を$r$とすると、$2n/k > (n-1)/(k-1) \geq r$が成り立ちます。
In fact, one $X$ operation increases the number of subtrees by one.
%実際、このアルゴリズムにおいて、操作$X$を1回施すと、部分木の個数は1つ増えます。
Therefore, the total overlap of the vertices of the subtrees is at most $r-1$.
%したがって、ここまでに得られた部分木の個数を$r$とすると、部分木の点の重複度は高々$r-1$に過ぎません。
From this it follows that $((r-1)+n)/k \geq r$, and solving this inequality yields $n-1 \geq (k-1)r$.
%このことから、$((r-1)+n)/k \geq  r$ であり、この不等式を解くと、$n-1 \geq (k-1)r$ です。
Now, obviously, if $k=1$, the degree of overlap is zero, hence we can assume $k \geq 2$, so we have $2n/k > (n-1)/(k-1) \geq r$.
%今、明らかに、$k=1$ならば重複度はゼロですから$k$は$2$以上としてよいので、$2n/k > (n-1)/(k-1) \geq  r$ です。
Therefore, even if we apply one more $X$ operation to each subtree here, the number of subtrees can be kept to $4n/k$.
%よって、ここで各subtreeを高々もう一回ずつcentroidで分解したとしても、その個数は$4n/k$個で抑えられます。
Thus, by this algorithm, $T$ is covered by at most $4n/k$ subtrees satisfying the conditions of the theorem.
%従って、このアルゴリズムにより、$T$は、高々$4n/k$個の位数$k$以下の$T$の部分木達により被覆されます。
\end{Proofof}

%%%%%%%%%%%%%%%%%%%%%%%%%%%%%%%%%%%%%%%%%%%%%%%%%%%%%%%%%%%%%%%%%%%%%%%%%%%%%%%%
%%%%%%%%%%%%%%%%%%%%%%%%%%%%%%%%%%%%%%%%%%%%%%%%%%%%%%%%%%%%%%%%%%%%%%%%%%%%%%%%
%%%%%%%%%%%%%%%%%%%%%%%%%%%%%%%%%%%%%%%%%%%%%%%%%%%%%%%%%%%%%%%%%%%%%%%%%%%%%%%%

\begin{Proofof}{Lemma\tume\ref{children}}
According to the definition of the subtree stored in the second child of each node in this binary tree $\BT$ (cf. Step\tume\ref{No2Child}), for the sibling node  $c_2(p(\ell(i)))$ of $\ell(i)$, the graph $X(r)-(X(\ell(i))-\JoiV(X(\ell(i))))$ includes the graph $X(c_2(p(\ell(i))))$. 
The statement of this lemma can be directly derived from this fact. 
%第2子ノードに格納する木の選び方（ステップ{\ref{No2Child}}）から、$\ell(i)$の兄弟(Sibling)である第2子leafノード$c_2(p(\ell(i)))$について、
%$X(c_2(p(\ell(i))))$は$X(r)-(X(\ell(i))-\JoiV(X(\ell(i))))$の中に含まれる。当該命題の言明はこの事実から直ちに導かれる。
\end{Proofof}

%%%%%%%%%%%%%%%%%%%%%%%%%%%%%%%%%%%%%%%%%%%%%%%%%%%%%%%%%%%%%%%%%%%%%%%%%%%%%%%%
%%%%%%%%%%%%%%%%%%%%%%%%%%%%%%%%%%%%%%%%%%%%%%%%%%%%%%%%%%%%%%%%%%%%%%%%%%%%%%%%
%%%%%%%%%%%%%%%%%%%%%%%%%%%%%%%%%%%%%%%%%%%%%%%%%%%%%%%%%%%%%%%%%%%%%%%%%%%%%%%%

\begin{Proofof}{Lemma\tume\ref{BT}}
%Lemma~\ref{bunkatsu}の証明で用いられるアルゴリズムを用いる。
Let $\NS(\ell)$ be the set of nodes that are at distance $\ell$ from the root node of the binary tree $\BT$.
%二分木$\BT$のルートノードから距離$\ell$にあるノードの集合を$\NS(\ell)$と置く。
Since $\forall i, \sum_{x \in\ NS(i)}|V(X(x))|\leq 4n$ holds, 
according to Lemma~\ref{centroid}, for any distance $\ell$, the effort required to correctly store all the objects that should be placed in all nodes of $\NS(\ell)$ is at most $\Od(n)$. 
%$\forall i, \sum_{x \in \NS(i)}|V(X(x))|\leq 4n$であるから、Lemma~\ref{centroid}を勘案すれば、任意の距離$\ell$に対して、$\NS(\ell)$の構成には、高々$\Od(n)$程度の計算量しか掛からない。
The height of $\BT$ is $\Od(\log n)$, so the total effort required to complete $\BT$ is $\Od(n\log n)$. 
%この二分木$\BT$の高さは高々$\Od(\log n)$であるから、この二分木$\BT$を完成させるのに必要な手間は、$\Od(n\log n )$である。
\end{Proofof}

%%%%%%%%%%%%%%%%%%%%%%%%%%%%%%%%%%%%%%%%%%%%%%%%%%%%%%%%%%%%%%%%%%%%%%%%%%%%%%%%
%%%%%%%%%%%%%%%%%%%%%%%%%%%%%%%%%%%%%%%%%%%%%%%%%%%%%%%%%%%%%%%%%%%%%%%%%%%%%%%%
%%%%%%%%%%%%%%%%%%%%%%%%%%%%%%%%%%%%%%%%%%%%%%%%%%%%%%%%%%%%%%%%%%%%%%%%%%%%%%%%

\begin{Proofof}{Lemma\tume\ref{constant}}
In this case, by applying Lemma\tume\ref{feasible-subtree} and consecutively solving at most 
$\Ord(\D(G))$ feasible sub-puzzles --- each on a board tree with at most $2M+2$ vertices --- 
along the path connecting any two pebbles, it is possible to swap their positions 
without altering the positions of the other pebbles. Since the size of each board tree is 
bounded by the constant $2M+2$ in these sub-puzzles, the effort and number of moves 
to rearrange any configuration on such a board tree is also in constant order $\Ord(1)$. 
Note that, in general, after performing a suitable preprocessing of $\Ord(n\D(G))$ steps, 
the entire configuration can be completed by repeating the operation of 
'swapping the positions of two pebbles' at most $n-1$ times. 
%%%%%%%%%%%%%%%%%%%%%%%%%%%%%%%%%%%%%%%%%%%%%%%%%%%%%%%%%%%%%%%%%%%%%%%%
%%%%%%%%%%%%%%%%%%%%%%%%%%%%%%%%%%%%%%%%%%%%%%%%%%%%%%%%%%%%%%%%%%%%%%%%
%%%%%%%%%%%%%%%%%%%%%%%%%%%%%%%%%%%%%%%%%%%%%%%%%%%%%%%%%%%%%%%%%%%%%%%%
%%%%%%%%%%%%%%%%%%%%%%%%%%%%%%%%%%%%%%%%%%%%%%%%%%%%%%%%%%%%%%%%%%%%%%%%
%%%%%%%%%%%%%%%%%%%%%%%%%%%%%%%%%%%%%%%%%%%%%%%%%%%%%%%%%%%%%%%%%%%%%%%%
%%%%%%%%%%%%%%%%%%%%%%%%%%%%%%%%%%%%%%%%%%%%%%%%%%%%%%%%%%%%%%%%%%%%%%%%
%%%%%%%%%%%%%%%%%%%%%%%%%%%%%%%%%%%%%%%%%%%%%%%%%%%%%%%%%%%%%%%%%%%%%%%%
Moreover, in general, for any vertex $u$ of $G$ occupied by a pebble $A$, 
and any other vertex $v$ occupied either by a different pebble $B$ or 
by an unoccupied blank space, the positions of pebble $A$ and pebble $B$ 
(or the blank space) can be exchanged—without altering the positions of 
the other pebbles—by performing at most $\Ord\left(\frac{\D(G)}{k}\right)$
sub-puzzles of size $\Ord(k)$ along a shortest path in $G$ connecting $u$ 
and $v$, whose length is at most $\D(G)$.
%%%%%%%%%%%%%%%%%%%%%%%%%%%%%%%%%%%%%%%%%%%%%%%%%%%%%%%%%%%%%%%%%%%%%%%%
%%%%%%%%%%%%%%%%%%%%%%%%%%%%%%%%%%%%%%%%%%%%%%%%%%%%%%%%%%%%%%%%%%%%%%%%
%%%%%%%%%%%%%%%%%%%%%%%%%%%%%%%%%%%%%%%%%%%%%%%%%%%%%%%%%%%%%%%%%%%%%%%%
%%%%%%%%%%%%%%%%%%%%%%%%%%%%%%%%%%%%%%%%%%%%%%%%%%%%%%%%%%%%%%%%%%%%%%%%
%%%%%%%%%%%%%%%%%%%%%%%%%%%%%%%%%%%%%%%%%%%%%%%%%%%%%%%%%%%%%%%%%%%%%%%%
%%%%%%%%%%%%%%%%%%%%%%%%%%%%%%%%%%%%%%%%%%%%%%%%%%%%%%%%%%%%%%%%%%%%%%%%
%%%%%%%%%%%%%%%%%%%%%%%%%%%%%%%%%%%%%%%%%%%%%%%%%%%%%%%%%%%%%%%%%%%%%%%%
From these observations, it immediately follows that, in this case, even if scanning 
the entire structure of the input tree $G$ requires up to $\Ord(N)$ time, 
the \textsc{Pebble Motion Problem on Trees} can still be computed in time  
$\Ord(N+n\D(G))$, and the length of the resulting 
solution sequence is bounded by $\Ord(n\D(G)).$
\end{Proofof}


%%%%%%%%%%%%%%%%%%%%%%%%%%%%%%%%%%%%%%%%%%%%%%%%%%%%%%%%%%%%%%%%%%%%%%%%%%%%%%%%
%%%%%%%%%%%%%%%%%%%%%%%%%%%%%%%%%%%%%%%%%%%%%%%%%%%%%%%%%%%%%%%%%%%%%%%%%%%%%%%%
%%%%%%%%%%%%%%%%%%%%%%%%%%%%%%%%%%%%%%%%%%%%%%%%%%%%%%%%%%%%%%%%%%%%%%%%%%%%%%%%

\begin{Proofof}{Lemma\tume\ref{subtrees-containing-v}}
Since Step\tume\maru{\ref{hige}} (Step\tume\maru{\ref{del-t}}, resp.) of $\Alg.\CutOff(G,T,\BT,f,g^{\prime})$ and Step\tume\maru{\ref{hige2}} (Step\tume\maru{\ref{del-ell}}, resp.) of $\Alg.\Extraction(T, \BT,g^{\prime},t)$ perform essentially the same operation, we will prove it only for the former.
The change in the maximum isthmus size of $X(r)-(X(t)-\JoiV(X(t)))$ can be achieved by checking at most an $\Ord(k)$ region around $\Joi(X(t))$, which can be done in $\Ord(k)$ time. During the initial construction of the binary tree $\BT$, if we provide an array containing information for each vertex of tree $T$ about which leaf node of $\BT$ stores the subtree that includes itself (a process that can be accomplished in total at most in $\Ord(n + k)$ time), then a leaf node $w$ of $\BT$ such that $X(w)$ contains the vertex $v$ can be found in $\Ord(1)$ time. Note that, conversely, for each leaf node $\ell$ of $\BT$, if we provide data indicating the position of $\ell$ in the list associated with each vertex of $X(\ell)$ (this task also can be handled in $\Ord(n+k)$ overall), then maintaining the lists associated with the vertices of $T$ with the removal of $\ell$ from $\BT$ at Step\tume\maru{\ref{del-t}} of $\Alg.\CutOff(G,T,\BT,f,g^{\prime})$ can be done in $\Ord(k)$ time. 
\end{Proofof}


%%%%%%%%%%%%%%%%%%%%%%%%%%%%%%%%%%%%%%%%%%%%%%%%%%%%%%%%%%%%%%%%%%%%%%%%%%%%%%%%
%%%%%%%%%%%%%%%%%%%%%%%%%%%%%%%%%%%%%%%%%%%%%%%%%%%%%%%%%%%%%%%%%%%%%%%%%%%%%%%%
%%%%%%%%%%%%%%%%%%%%%%%%%%%%%%%%%%%%%%%%%%%%%%%%%%%%%%%%%%%%%%%%%%%%%%%%%%%%%%%%

\begin{Proofof}{Lemma\tume\ref{mini-jyunbi}}
Let the set of all nodes in the binary tree $\BT2$ be denoted as $\NS(\BT2)$.
%二分木$\BT2$のすべてのノードからなる集合を$\NS(\BT2)$と書くことにする。
Let $H$ be a subtree of $T$ containing all the vertices $\{\JoiV(X(x)) \mid x\in \NS(\BT2)\}$ and all the edges $\{\JoiE(X(x)) \mid x\in \NS(\BT2)\}$.  
% $\BT2$のすべてのノード$v$に格納されている$\JoiV(X(v))$および$\JoiE(X(v))$を含むような$T$の部分木$H$を考える。
The movement of the $k+1$ unoccupied spaces through the repetition of Step\tume\ref{mini-puzzle} in the algorithm is accomplished, as a whole, by having a connected unoccupied space of size $k+1$ traverse $H$ in a depth-first search order. 
%題意のステップの繰り返しは、$k$個の連結な空白領域が$H$を深さ優先探索の順番に回遊していくことにより達成される。
The estimated number of moves for one unoccupied space (that is achieved by the sequence of swapping the positions of this unoccupied space and pebbles) is $\Ord(k+n)$, so the total effort for swapping the unoccupied space of size $k+1$ is bounded by $\Ord(k(k+n))$.
%その際の空白と碁石の位置の交換の手数の見積は空白一つ当たりトータルで$\Ord(k+n)$程度だから、$k$個の空白全体の交換の手間は$\Ord(k(k+n))$で抑えられる。
%なおステップ\maru{\ref{mini-puzzle}}の繰り返しにおいて、$k$個の連結な空白領域の形状や位置が異なっている可能性はあるので、それが交換の手数に影響するか否かについて検討しておく必要がある。ある段階でのステップ\maru{\ref{mini-puzzle-finish}}の終了時の空白の位置から、次回のステップ\maru{\ref{mini-puzzle}}の新たな連結領域上に空白を移動させるための交換の回数は、サブパズルのサイズが$\Ord(k)$なので、各回につき高々$\Ord(k^2)$程度の誤差（増減）しか生じない。よってトータルでも$\Ord(k^2\cdot(k+n)/k)=\Ord(k(k+n))$程度の誤差である。
\end{Proofof}


%%%%%%%%%%%%%%%%%%%%%%%%%%%%%%%%%%%%%%%%%%%%%%%%%%%%%%%%%%%%%%%%%%%%%%%%%%%%%%%%
%%%%%%%%%%%%%%%%%%%%%%%%%%%%%%%%%%%%%%%%%%%%%%%%%%%%%%%%%%%%%%%%%%%%%%%%%%%%%%%%
%%%%%%%%%%%%%%%%%%%%%%%%%%%%%%%%%%%%%%%%%%%%%%%%%%%%%%%%%%%%%%%%%%%%%%%%%%%%%%%%



%%%%%%%%%%%%%%%%%%%%%%%%%%%%%%%%%%%%%%%%%%%%%%%%%%%%%%%%%%%%%%%%%%%%%%%%%%%%%%%%
%%%%%%%%%%%%%%%%%%%%%%%%%%%%%%%%%%%%%%%%%%%%%%%%%%%%%%%%%%%%%%%%%%%%%%%%%%%%%%%%
%%%%%%%%%%%%%%%%%%%%%%%%%%%%%%%%%%%%%%%%%%%%%%%%%%%%%%%%%%%%%%%%%%%%%%%%%%%%%%%%




%%%%%%%%%%%%%%%%%%%%%%%%%%%%%%%%%%%%%%%%%%%%%%%%%%%%%%%%%%%%%%%%%%%%%%%%%%%%%%%%
%%%%%%%%%%%%%%%%%%%%%%%%%%%%%%%%%%%%%%%%%%%%%%%%%%%%%%%%%%%%%%%%%%%%%%%%%%%%%%%%
%%%%%%%%%%%%%%%%%%%%%%%%%%%%%%%%%%%%%%%%%%%%%%%%%%%%%%%%%%%%%%%%%%%%%%%%%%%%%%%%




%%%%%%%%%%%%%%%%%%%%%%%%%%%%%%%%%%%%%%%%%%%%%%%%%%%%%%%%%%%%%%%%%%%%%%%%%%%%%%%%
%%%%%%%%%%%%%%%%%%%%%%%%%%%%%%%%%%%%%%%%%%%%%%%%%%%%%%%%%%%%%%%%%%%%%%%%%%%%%%%%
%%%%%%%%%%%%%%%%%%%%%%%%%%%%%%%%%%%%%%%%%%%%%%%%%%%%%%%%%%%%%%%%%%%%%%%%%%%%%%%%

\bibliographystyle{alphaurl}
%\bibliographystyle{plainurl}% the mandatory bibstyle
\bibliography{pebble-motion-problem-on-trees.bib}
%\nocite{*}

%%
%%\textbf{Acknowlegments.\/} 
%\begin{thebibliography}{1}
%
%\bibitem{WOS:000083999100001}
%A. Archer,
%\newblock A modern treatment of the $15$ puzzle,
%\newblock \textsl{American Math. Monthly}
%\newblock \textbf{106}\,(1999), 793--799.
%\bibitem{WOS:001177663400002}
%\newblock S. Ardizzoni, I. Saccani, L. Consolini, M. Locatelli, and B. Nebel, 
%\newblock An Algorithm with Improved Complexity for Pebble Motion/Multi-Agent Path Finding on Trees,
%\newblock \textsl{Journal of Artificial Intelligence Research}
%\newblock \textbf{79}\,(2024), 483--514.
%\bibitem{WOS:000077614700003}
%V. Auletta, A. Monti, M. Parente, and P. Persiano,
%\newblock A linear-time algorithm for the feasibility of pebble motion on trees,
%\newblock \textsl{Algorithmica}
%\newblock \textbf{23}\,(1999), 223--245.
%\bibitem{WOS:000417242300006}
%R. B. Bapat, S. Fujita, S. Legay, Y. Manoussakis, Y. Matsui, T. Sakuma, and Z. Tuza, 
%\newblock Safe Sets, Network Majority on Weighted Trees,
%\newblock \textsl{Networks.}
%\newblock \textbf{71}\,(2018), 81--92.
%\bibitem{WOS:000254460900008}
%G. Calinescu, A. Dumitrescu, and J. Pach, 
%\newblock Reconfigurations in Graphs and Grids,
%\newblock \textsl{SIAM J. Discrete Math.}
%\newblock \textbf{22}\,(2008), 124--138.
%%\bibitem{CEBMENMC2021}
%%B. Cimring,  R. El Sabeh, M. Bacvanski, S. Maaz,  I. El Hajj, N. Nishimura,  AE, Mouawad, and A. Cooper
%%\newblock Efficient algorithms to solve atom reconfiguration problems. I. Redistribution-reconfiguration algorithm, 
%%\newblock \textsl{Physical Review A}
%%\newblock \textbf{108}\,Issue 2,(2023), 023107.
%\bibitem{zbMATH05966071}
%A. Fink, R. Guy,
%\newblock Rick's tricky six puzzle: $S_5$ sits specially in $S_6$,
%\newblock \textsl{Math. Magazine}
%\newblock \textbf{82}\,(2009), 83--102.
%\bibitem{WOS:000301306200015}
%S. Fujita, T. Nakamigawa, and T. Sakuma,
%\newblock Colored Pebble Motion on Graphs,
%\newblock \textsl{European J. of Combin.}
%\newblock \textbf{33}\,(2012), 884--892.
%\bibitem{zbMATH06431549}
%S. Fujita, T. Nakamigawa, and T. Sakuma,
%\newblock Pebble Exchange on Graphs,
%\newblock \textsl{Disc. Appl. Math.}
%\newblock \textbf{184}\,(2015), 281--287.
%\bibitem{F1949}
%R. Frucht,
%\newblock Graphs of degree three with a given abstract group,
%\newblock \textsl{Canad. J. Math.}
%\newblock \textbf{1}\,(1949), 365--378.
%\bibitem{aa887662-c3e2-30b0-a475-bcd8c655206a}
%Wm. Woolsey Johnson and William E. Story,
%\newblock Notes on the $15$ puzzle,
%\newblock \textsl{American J. Math.}
%\newblock \textbf{2}\,(1879), 397--399.
%\bibitem{zbMATH07360934}
%T. Kato, T. Nakamigawa, and T. Sakuma,
%\newblock Pebble exchange group of graphs,
%\newblock \textsl{European J. of Combin.}
%\newblock \textbf{95}\,(2021), 103325.
%\bibitem{715921}
%D. Kohnhauser, G. Miller, and P. Spirakis, 
%\newblock Coordinating pebble motion on graphs, the diameter of permutation groups, and applications,
%\newblock \textsl{Proceedings of the 25-th Symposium on Foundations of Computer Science},
%\newblock (FOCS '84), 241--250.
%\bibitem{loyd1959mathematical}
%S. Loyd,
%\newblock \textsl{Mathematical Puzzles of Sam Loyd},
%\newblock Dover, New York, 1959.
%\bibitem{365740}
%C. Papadimitriou, P. Raghavan, M. Sudan, and H. Tamaki,
%\newblock Motion planning on a graph,
%\newblock \textsl{Proceedings of the 35-th Symposium on Foundations of Computer Science},
%\newblock (FOCS '94), 511--520.
%\bibitem{zbMATH04155895}
%D. Ratner and M. Warmuth,
%\newblock Finding a shortest solution for the ($N \times N$)-extension of the 15-puzzle is intractable,
%\newblock \textsl{J. Symbolic Computation}
%\newblock \textbf{10}\,(1990), 111--137.
%\bibitem{zbMATH03446929}
%R. M. Wilson,
%\newblock Graph puzzles, homotopy, and the alternating group,
%\newblock \textsl{J. Combin. Theory Ser. B}
%\newblock \textbf{16}\,(1974), 86--96.
%\end{thebibliography}
%
\end{document}
